% !TeX program = xelatex
\documentclass[11pt]{article}
\usepackage[a4paper,top=2.1cm,bottom=2.1cm,left=2.2cm,right=2.2cm]{geometry}
\usepackage[T1]{fontenc}
\usepackage[utf8]{inputenc}
\usepackage[table,dvipsnames]{xcolor}
\usepackage{amsmath,amssymb,amsthm}
\usepackage{booktabs,longtable,array,ragged2e,enumitem}
\usepackage{titlesec,fancyhdr,mdframed}
\usepackage[most]{tcolorbox}
\usepackage[htt]{hyphenat}
\usepackage[colorlinks=true,linkcolor=NavyBlue,urlcolor=NavyBlue,citecolor=NavyBlue]{hyperref}

% Box colours by role, each taken from the palette sig-brand.sty defines. The
% four semantic annotation colours below are outside that palette on purpose:
% the prose names them (blue for Midnight's layer, red for secrets) so they
% must stay recognisable whatever the brand theme.
\colorlet{boxtitle}{SigDark}
\colorlet{boxrule}{SigShell}
\colorlet{boxfill}{SigShell!20!SigCream}
\colorlet{boxbg}{SigCream}
\colorlet{pubbg}{SigBrand}
\colorlet{pubrule}{SigHeader}
\colorlet{addrule}{SigDark}
\definecolor{mntbg}{HTML}{D3E6F7}     % Midnight-provided component fill
\definecolor{mntrule}{HTML}{3E7CB1}   % Midnight-provided component border
\definecolor{bordo}{RGB}{128,0,32}    % secret values in protocol text


\titleformat{\section}{\Large\bfseries\color{hdblue}}{\thesection}{0.7em}{}
\titleformat{\subsection}{\normalsize\bfseries\color{hdblue}}{\thesubsection}{0.5em}{}
\titleformat{\subsubsection}{\normalsize\bfseries\color{hdblue}}{\thesubsubsection}{0.5em}{}
\pagestyle{fancy}\fancyhf{}
\newcolumntype{L}[1]{>{\RaggedRight\arraybackslash}p{#1}}
\renewcommand{\arraystretch}{1.18}

% Document-wide environments, theorem kinds and notation. Colours live in
% palette.tex, packages and page layout in preamble.tex.

% =====================================================================
%  Boxes and lists
% =====================================================================
\newcommand{\code}[1]{\texttt{\small #1}}

\newcounter{expfig}
\newenvironment{ExpFig}[2]%
 {\refstepcounter{expfig}\label{#1}%
  \begin{mdframed}[linecolor=boxrule,linewidth=0.7pt,backgroundcolor=boxfill,
    innertopmargin=7pt,innerbottommargin=7pt,innerleftmargin=9pt,innerrightmargin=9pt,
    skipabove=9pt,skipbelow=3pt]\small
  \begin{center}\textbf{#2}\end{center}\vspace{-2pt}\noindent}
 {\end{mdframed}\vspace{-4pt}\begin{center}\footnotesize Fig.~\arabic{expfig}\end{center}\vspace{4pt}}

\newenvironment{ProcBox}[1]%
 {\begin{mdframed}[linecolor=boxrule,linewidth=0.8pt,backgroundcolor=boxbg,
    innertopmargin=6pt,innerbottommargin=6pt,innerleftmargin=8pt,innerrightmargin=8pt,
    frametitle={\normalsize\bfseries\color{boxtitle}#1},frametitlebackgroundcolor=boxfill,
    frametitlerule=true,frametitlerulewidth=0.4pt,skipabove=8pt]\small}%
 {\end{mdframed}}

\newenvironment{PubBox}[1]%
 {\begin{mdframed}[linecolor=pubrule,linewidth=0.9pt,backgroundcolor=pubbg,
    innertopmargin=6pt,innerbottommargin=6pt,innerleftmargin=8pt,innerrightmargin=8pt,
    skipabove=7pt,skipbelow=7pt]\small
  \textbf{\textsf{#1}}\par\vspace{2pt}}
 {\end{mdframed}}

% A circuit box numbers itself C1, C2, ... through \ref.
\newcounter{cktctr}
\renewcommand{\thecktctr}{C\arabic{cktctr}}
\newenvironment{CktBox}[2]%
 {\refstepcounter{cktctr}\label{#1}%
  \begin{mdframed}[linecolor=boxtitle,linewidth=0.8pt,backgroundcolor=boxbg,
    innertopmargin=6pt,innerbottommargin=6pt,innerleftmargin=8pt,innerrightmargin=8pt,
    frametitle={\color{boxtitle}Circuit~\thecktctr{} --- #2},
    frametitlebackgroundcolor=boxfill,frametitlerule=true,frametitlerulewidth=0.4pt,
    skipabove=9pt,skipbelow=9pt]\small\sloppy}%
 {\end{mdframed}}

\newenvironment{TxtBox}[2]%
 {\label{#1}%
  \begin{mdframed}[linecolor=boxtitle,linewidth=0.8pt,backgroundcolor=boxbg,
    innertopmargin=6pt,innerbottommargin=6pt,innerleftmargin=8pt,innerrightmargin=8pt,
    frametitle={\normalsize\bfseries\color{boxtitle} #2},
    frametitlebackgroundcolor=boxfill,frametitlerule=true,frametitlerulewidth=0.4pt,
    skipabove=9pt,skipbelow=9pt]\small\sloppy}%
 {\end{mdframed}}

\newcounter{stepctr}[subsection]
\newenvironment{Steps}%
 {\begin{list}{\footnotesize\textcolor{boxtitle}{Step \arabic{stepctr}.}}%
   {\usecounter{stepctr}\setlength{\leftmargin}{3.6em}\setlength{\itemsep}{3pt}%
    \setlength{\parsep}{0pt}\setlength{\topsep}{3pt}\setlength{\labelwidth}{3.2em}}}%
 {\end{list}}

\newenvironment{OutList}%
 {\smallskip\noindent\textbf{\textit{Output.}}
  \begin{list}{}{\setlength{\leftmargin}{3.0em}\setlength{\labelwidth}{2.6em}%
   \setlength{\labelsep}{0.4em}\setlength{\itemsep}{2pt}\setlength{\parsep}{0pt}%
   \setlength{\topsep}{2pt}\setlength{\itemindent}{0pt}}}%
 {\end{list}}
\newcommand{\oitem}[1]{\item[$#1$\,:]}
\newcommand{\fld}[1]{\smallskip\noindent\textbf{\textit{#1}}\ }

% Blue shading marks components Midnight provides.
\newcommand{\mnhl}[1]{\colorbox{mntbg}{\small #1}}
\newcommand{\MidNote}[1]{\par\smallskip\noindent
  \fcolorbox{mntrule}{mntbg}{\parbox{\dimexpr\linewidth-2\fboxsep-2\fboxrule\relax}{\small #1}}\par\smallskip}
\newcommand{\MidNoteT}[1]{\fcolorbox{mntrule}{mntbg}{\small #1}}
\newcommand{\bluedot}{\textcolor{mntrule}{$\blacksquare$}\,}

\newcommand{\lnk}[2]{\hyperlink{#1}{\ensuremath{#2}}}
\newcommand{\cktref}[1]{\hyperref[#1]{Circuit~\ref*{#1}}}

% =====================================================================
%  Theorem-like environments
% =====================================================================
\theoremstyle{definition}
\newtheorem{definition}{Definition}
\newtheorem{remark}{Remark}
\newtheorem{assumption}{Usage assumption}
\newtheorem{gap}{Gap}
\newtheorem{obligation}{Proof obligation}
\newtheorem{theorem}{Theorem}
\newtheorem{lemma}{Lemma}

% =====================================================================
%  Editorial marks
% =====================================================================
% \DRAFTfalse silences the inline author notes for the build that goes out.
\newif\ifDRAFT
\DRAFTtrue
\newcommand{\ale}[1]{\textcolor{red}{#1}}
\ifDRAFT\else\renewcommand{\ale}[1]{}\fi
\newcommand{\remove}[1]{}
\newcommand{\upd}[1]{#1}
\newcommand{\ADDEDINL}[1]{\textcolor{addrule}{#1}}
\newcommand{\newversion}[1]{{\textcolor{boxtitle}{\textbf{Differences with the latest release (\latestTag):}} #1}}

% Colour roles inside protocol text.
\newcommand{\secret}[1]{\textcolor{bordo}{#1}}
\newcommand{\freshrn}[1]{\textcolor{Mulberry}{#1}}

% =====================================================================
%  Notation
% =====================================================================

% --- the vault release this document analyses, and the latest release ---
\newcommand{\analysedTag}{erc20-vault-v0.3.0}
\newcommand{\latestTag}{erc20-vault-v0.4.0}
\newcommand{\vaultRepoUrl}{https://github.com/sig-net/midnight-examples}
\newcommand{\analysedContractUrl}{\vaultRepoUrl/blob/\analysedTag/examples/erc20-vault/contract/src/erc20-vault.compact}
\newcommand{\analysedReadmeUrl}{\vaultRepoUrl/blob/\analysedTag/examples/erc20-vault/README.md}
\newcommand{\latestContractUrl}{\vaultRepoUrl/blob/\latestTag/examples/erc20-vault/contract/src/erc20-vault.compact}

% --- general ---
\newcommand{\secpar}{\lambda}
\newcommand{\negl}{\mu}
\newcommand{\Adv}{\mathcal{A}}
\newcommand{\Chal}{\mathcal{C}}
\newcommand{\distprf}{\mathcal{B}}
\newcommand{\numusers}{N}
\newcommand{\bits}[1]{\{0,1\}^{#1}}
\newcommand{\assert}{\textbf{assert}~}
\newcommand{\concat}{\,\|\,}
\newcommand{\bc}{,\penalty0\ }
\newcommand{\Epsilon}{\varepsilon}
\newcommand{\Payment}{Transaction}
\newcommand{\payment}{transaction}
\newcommand{\zkswap}{ZKSwap}
\newcommand{\SigSys}{\mathsf{SigNet}}

% --- pseudorandom functions ---
\newcommand{\PRFfam}{\mathsf{F}}
\newcommand{\prfkey}{\mathsf{k}}
\newcommand{\prfKeySp}{\mathcal{K}}
\newcommand{\prfDom}{\mathcal{X}}
\newcommand{\prfRng}{\mathcal{Y}}
\newcommand{\RandFun}{\mathsf{RF}}
\newcommand{\AdvPRF}{\mathsf{Adv}^{\mathsf{prf}}}
\newcommand{\AdvMUPRF}{\mathsf{Adv}^{\mathsf{mu\text{-}prf}}}
\newcommand{\oracle}{\mathcal{O}}

% --- cryptographic primitives ---
\newcommand{\Hash}{\mathsf{HashK}}
\newcommand{\PersistantHash}{\mathsf{PersistantHash}}
\newcommand{\Kec}{\mathsf{HashP}}
\newcommand{\SHA}{\mathsf{SHA\text{-}256}}
\newcommand{\EcdsaVf}{\mathsf{EcdsaVerify}}
\newcommand{\NIZK}{\Pi_{\mathsf{zk}}}
\newcommand{\zkSetup}{\mathsf{ZK.Setup}}
\newcommand{\zkProve}{\mathsf{ZK.Prove}}
\newcommand{\zkVerify}{\mathsf{ZK.Verify}}
\newcommand{\zkSim}{\mathsf{ZK.Sim}}
\newcommand{\crs}{\mathsf{crs}}
\newcommand{\td}{\mathsf{td}}
\newcommand{\CS}{\mathsf{CS}}
\newcommand{\csGen}{\mathsf{CS.Gen}}
\newcommand{\csCom}{\mathsf{CS.Com}}
\newcommand{\ck}{\mathsf{ck}}
\newcommand{\MAC}{\mathsf{MAC}}
\newcommand{\macTag}{\mathsf{Tag}}
\newcommand{\macVer}{\mathsf{Vrfy}}
\newcommand{\SIG}{\mathsf{SIG}}
\newcommand{\sigKG}{\mathsf{SIG.KGen}}
\newcommand{\sigSign}{\mathsf{SIG.Sign}}
\newcommand{\sigVer}{\mathsf{SIG.Vrfy}}
\newcommand{\AdvZK}{\mathsf{Adv}^{\mathsf{zk}}}
\newcommand{\AdvSnd}{\mathsf{Adv}^{\mathsf{snd}}}
\newcommand{\AdvHide}{\mathsf{Adv}^{\mathsf{hide}}}
\newcommand{\AdvBind}{\mathsf{Adv}^{\mathsf{bind}}}
\newcommand{\AdvMAC}{\mathsf{Adv}^{\mathsf{euf\text{-}cma}}}
\newcommand{\AdvSIG}{\mathsf{Adv}^{\mathsf{euf\text{-}cma}}_{\mathsf{akd}}}

% --- ledgers, transcripts and hybrids ---
\newcommand{\Lpub}{\mathcal{L}_{\mathsf{pub}}}
\newcommand{\Lshd}{\mathcal{L}_{\mathsf{shd}}}
\newcommand{\LExecute}{\mathsf{Execute}}
\newcommand{\LRead}{\mathsf{Read}}
\newcommand{\Tpub}{T_{\mathsf{pub}}}
\newcommand{\Tshd}{T_{\mathsf{shd}}}
\newcommand{\Tall}{T}
\newcommand{\hybrid}{\ensuremath{\mathsf{Hyb}}}
\newcommand{\contract}{\ensuremath{\mathsf{contract}}}
\newcommand{\contractinput}{\ensuremath{\mathsf{input}}}
\newcommand{\blockheight}{\ensuremath{\mathsf{blockheight}}}
\newcommand{\kind}{\ensuremath{\mathsf{kind}}}
\newcommand{\outlt}{\ensuremath{\mathsf{length}}}

% --- addresses, accounts, values ---
\newcommand{\pubAddr}{\mathsf{pubAddr}}
\newcommand{\shdAcct}{\mathsf{shdAcct}}
\newcommand{\destAddr}{\mathsf{destAddr}}
\newcommand{\srcAcct}{\shdAcct}
\newcommand{\recvAcct}{\mathsf{recvAcct}}
\newcommand{\setlAcct}{\mathsf{setlAcct}}
\newcommand{\rcpt}{\mathsf{rcpt}}
\newcommand{\usrTag}[1]{\mathsf{user}(#1)}
\newcommand{\candZero}{\mathsf{shAcct}_0}
\newcommand{\candOne}{\mathsf{shAcct}_1}
\newcommand{\candB}{\mathsf{shAcct}_b}
\newcommand{\callerZKey}{\shdAcct}
\newcommand{\callerZKeySk}{\mathsf{shdAcctSKey}}
\newcommand{\tokType}{\mathsf{tokType}}
\newcommand{\tokTypeB}{\mathsf{tokType}'}
\newcommand{\tokTypeIn}{\tokType_{\mathrm{in}}}
\newcommand{\tokTypeOut}{\tokType_{\mathrm{out}}}
\newcommand{\amt}{\mathsf{amt}}
\newcommand{\amtB}{\mathsf{amt}'}
\newcommand{\amtRem}{\mathsf{amtRemainder}}
\newcommand{\amtIn}{\mathsf{amtSpent}}
\newcommand{\amtSpent}{\mathsf{amtSpent}}
\newcommand{\amtMax}{\mathsf{amtMax}}
\newcommand{\amtOut}{\mathsf{amtOut}}
\newcommand{\amtChg}{\mathsf{amtChange}}
\newcommand{\amtShares}{\mathsf{amtShares}}
\newcommand{\amtAssets}{\mathsf{amtAssets}}
\newcommand{\bal}{\mathsf{vaultBal}}
\newcommand{\ledgerBal}{\mathsf{ledgerBal}}
\newcommand{\colr}{\mathsf{color}}
\newcommand{\colrU}{\colr_{\mathsf{u}}}
\newcommand{\colrW}{\colr_{\mathsf{w}}}

% --- coins ---
\newcommand{\coinClaim}{\mathsf{newMintedCoin}}
\newcommand{\coinRem}{\mathsf{remainderCoin}}
\newcommand{\coinIn}{\mathsf{surrCoin}}
\newcommand{\coinMint}{\mathsf{newCoin}}
\newcommand{\coinMintW}{\mathsf{newCoinShares}}
\newcommand{\coinMintU}{\mathsf{newCoinAssets}}
\newcommand{\coinChange}{\mathsf{changeCoin}}
\newcommand{\coinBurnOut}{\mathsf{burnOutCoin}}
\newcommand{\coinFund}{\mathsf{fundCoin}}
\newcommand{\coinCommit}{\mathsf{coinCommit}}
\newcommand{\coinNull}{\mathsf{coinNull}}
\newcommand{\evolveNonce}{\mathsf{evolveNonce}}
\newcommand{\tokenTypeOf}{\mathsf{tokenTypeOf}}

% --- procedure interface ---
\newcommand{\pp}{\mathsf{pp}}
\newcommand{\Setup}{\mathsf{Setup}}
\newcommand{\Deposit}{\mathsf{Deposit}}
\newcommand{\Withdraw}{\mathsf{Withdraw}}
\newcommand{\Swap}{\mathsf{Swap}}
\newcommand{\Supply}{\mathsf{Supply}}
\newcommand{\Redeem}{\mathsf{Redeem}}
\newcommand{\act}{\mathsf{act}}
\newcommand{\auxx}{\mathsf{aux}}
\newcommand{\privIn}{\mathsf{privIn}}
\newcommand{\prvOut}{\mathsf{privOut}}
\newcommand{\outc}{\mathsf{outcome}}
\newcommand{\stpub}{\mathsf{state}_{\mathsf{pub}}}
\newcommand{\stshd}{\mathsf{state}_{\mathsf{shd}}}
\newcommand{\trans}{\mathsf{transcript}}
\newcommand{\succeeded}{\mathsf{succeeded}}
\newcommand{\failed}{\mathsf{failed}}
\newcommand{\pending}{\mathsf{pending}}
\newcommand{\schemaIn}{\mathsf{schemaIn}}
\newcommand{\schemaOut}{\mathsf{schemaOut}}
\newcommand{\okbit}{\mathsf{ok}}
\newcommand{\notif}{\mathsf{notify}}

% --- oracles ---
\newcommand{\ONewPub}{\mathcal{O}\mathsf{NewPubAcct}}
\newcommand{\ONewShd}{\mathcal{O}\mathsf{NewShdAcct}}
\newcommand{\OMint}{\mathcal{O}\mathsf{Mint}}
\newcommand{\OTransPub}{\mathcal{O}\mathsf{TransferPub}}
\newcommand{\OTransShd}{\mathcal{O}\mathsf{TransferShd}}
\newcommand{\OShdTrans}{\mathcal{O}\mathsf{ShieldedTransfer}}
\newcommand{\ODep}{\mathcal{O}\mathsf{Deposit}}
\newcommand{\OWd}{\mathcal{O}\mathsf{Withdraw}}
\newcommand{\OSw}{\mathcal{O}\mathsf{Swap}}
\newcommand{\OSup}{\mathcal{O}\mathsf{Supply}}
\newcommand{\ORed}{\mathcal{O}\mathsf{Redeem}}

% --- games ---
\newcommand{\unlinkgameG}{\mathsf{Exp\text{-}WithdrawUnlink}}
\newcommand{\lendunlinkgame}{\mathsf{Exp\text{-}LendUnlink}}
\newcommand{\unlinkgame}{\mathsf{Exp\text{-}Unlink}}

% --- the general unlinkability experiment's parameters ---
\newcommand{\colIn}{\tokTypeIn}
\newcommand{\valIn}{\amt}
\newcommand{\parm}{\auxx}

% --- a circuit call's output components ---
\newcommand{\Dvault}{\Delta_{\mathsf{vault}}}
\newcommand{\Dshd}{\Delta_{\mathsf{coin}}}
\newcommand{\stmt}{x}
\newcommand{\wit}{w}
\newcommand{\Rel}{R}
\newcommand{\zkproof}{\pi}
\newcommand{\zkproofs}{\mathsf{zk proofs}}
\newcommand{\comIO}{\mathsf{comCallIO}}

% --- circuits ---
\newcommand{\cktInit}{\code{initialise}}
\newcommand{\cktApprove}{\code{approveRouter}}
\newcommand{\cktApproveStata}{\code{approveStata}}
\newcommand{\cktApproveWrap}{\code{approveStata}}
\newcommand{\cktStartDeposit}{\code{startDeposit}}
\newcommand{\cktDeposit}{\cktStartDeposit}
\newcommand{\cktCompleteDeposit}{\code{completeDeposit}}
\newcommand{\cktStartWithdraw}{\code{startWithdraw}}
\newcommand{\cktWithdraw}{\cktStartWithdraw}
\newcommand{\cktCompleteWithdraw}{\code{completeWithdraw}}
\newcommand{\cktRefundWithdraw}{\code{refundWithdraw}}
\newcommand{\cktStartSwap}{\code{startSwap}}
\newcommand{\cktSwap}{\cktStartSwap}
\newcommand{\cktCompleteSwap}{\code{completeSwap}}
\newcommand{\cktRefundSwap}{\code{refundSwap}}
\newcommand{\cktStartSupply}{\code{startSupply}}
\newcommand{\cktSupply}{\cktStartSupply}
\newcommand{\cktCompleteSupply}{\code{completeSupply}}
\newcommand{\cktRefundSupply}{\code{refundSupply}}
\newcommand{\cktStartRedeem}{\code{startRedeem}}
\newcommand{\cktRedeem}{\cktStartRedeem}
\newcommand{\cktCompleteRedeem}{\code{completeRedeem}}
\newcommand{\cktRefundRedeem}{\code{refundRedeem}}
\newcommand{\cktFlush}{\code{flushQueue}}
\newcommand{\cktSend}{\code{sendX}}
\newcommand{\cktQueueAtt}{\code{queueAttestationN}}
\newcommand{\cktFailGate}{\mathsf{failGate}}
\newcommand{\opReceive}{\code{receiveShielded}}
\newcommand{\opMint}{\code{mintShieldedToken}}
\newcommand{\mintShieldedToken}{\ensuremath{\mathtt{mintShieldedToken}}}

% --- on-chain objects ---
\newcommand{\vaultAddr}{\mathsf{vaultAddr}}
\newcommand{\vaultEvmAddr}{\mathsf{vaultEvmAddr}}
\newcommand{\derivedAddr}{\mathsf{userDerivedAddr}}
\newcommand{\burnAddr}{\mathsf{burnAddr}}
\newcommand{\routerAddr}{\mathsf{routerAddr}}
\newcommand{\underAddr}{\mathsf{underlyingAddr}}
\newcommand{\wrapAddr}{\mathsf{wrapperAddr}}
\newcommand{\wrapOf}[1]{\mathsf{wrap}(#1)}
\newcommand{\mpcRootPk}{\mathsf{mpcRootPk}}
\newcommand{\mpcRespKey}{\mathsf{mpcRespKey}}
\newcommand{\Derive}{\mathsf{Derive}}
\newcommand{\keyPath}{\mathsf{keyPath}}
\newcommand{\keyPathLit}{\text{``vault''}}
\newcommand{\keyVer}{\mathsf{keyVersion}}
\newcommand{\chainId}{\mathsf{chainId}}
\newcommand{\initFlag}{\mathsf{initialized}}
\newcommand{\sigNonce}{\mathsf{reqCounter}}
\newcommand{\req}{\mathsf{req}}
\newcommand{\rid}{\mathsf{reqId}}
\newcommand{\reqIdOf}{\Kec}
\newcommand{\attSig}{\mathsf{attSig}}
\newcommand{\attDigest}{\mathsf{attDigest}}
\newcommand{\evmOut}{\mathsf{evmOut}}
\newcommand{\evmTxOf}{\mathsf{evmTx}}
\newcommand{\evmNonce}{\mathsf{evmNonce}}
\newcommand{\gasEnv}{\mathsf{gasEnvelope}}
\newcommand{\feeTier}{\mathsf{feeTier}}
\newcommand{\domSep}{\mathsf{domSep}}
\newcommand{\userComm}{\mathsf{userComm}}
\newcommand{\refundComm}{\mathsf{refundComm}}
\newcommand{\deplComm}{\mathsf{deployerComm}}

% --- ledger maps and settle views ---
\newcommand{\reqMap}{\mathsf{reqMap}}
\newcommand{\depositMap}{\mathsf{depositMap}}
\newcommand{\swapMap}{\mathsf{swapMap}}
\newcommand{\supplyMap}{\mathsf{supplyMap}}
\newcommand{\redeemMap}{\mathsf{redeemMap}}
\newcommand{\depositView}{\mathsf{depositView}}
\newcommand{\wdView}{\mathsf{wdView}}
\newcommand{\swapView}{\mathsf{swapView}}
\newcommand{\supplyView}{\mathsf{supplyView}}
\newcommand{\redeemView}{\mathsf{redeemView}}
\newcommand{\wdMarkerMap}{\wdView}
\newcommand{\swapMarkerMap}{\swapView}
\newcommand{\supplyMarkerMap}{\supplyView}
\newcommand{\redeemMarkerMap}{\redeemView}
\newcommand{\outputAttestationBuffer}{\ensuremath{\mathtt{outputAttestationBuffer}}}
\newcommand{\inputAttestationBuffer}{\ensuremath{\mathtt{inputAttestationBuffer}}}
\newcommand{\inputRequestBuffer}{\ensuremath{\mathtt{inputRequestBuffer}}}
\newcommand{\outputRequestBuffer}{\ensuremath{\mathtt{outputRequestBuffer}}}
\newcommand{\lastseen}{\ensuremath{\mathtt{globalLastSeen}}}
\newcommand{\vaultNonce}{\ensuremath{\mathtt{vaultAccountNonce}}}
\newcommand{\inIndex}{\mathsf{inIndex}}
\newcommand{\ownComm}{\mathsf{ownComm}}

% --- fields of a settle view ---
\newcommand{\fComm}{\mathsf{comm}}
\newcommand{\fErc}{\mathsf{erc20}}
\newcommand{\fAmt}{\mathsf{amt}}
\newcommand{\fTokIn}{\mathsf{tokenIn}}
\newcommand{\fTokOut}{\mathsf{tokenOut}}
\newcommand{\fAmtOut}{\mathsf{amtOut}}
\newcommand{\fAmtMax}{\mathsf{amtMax}}
\newcommand{\fShares}{\mathsf{shares}}

% --- secrets, nonces, randomness ---
\newcommand{\setupstate}{\rho}
\newcommand{\seed}{\mathsf{seed}}
\newcommand{\gen}{\mathsf{g}}
\newcommand{\usersk}{\mathsf{userSk}}
\newcommand{\coinsk}{\mathsf{coinSk}}
\newcommand{\coinpk}{\mathsf{coinPk}}
\newcommand{\encsk}{\mathsf{encSk}}
\newcommand{\encpk}{\mathsf{encPk}}
\newcommand{\mintNonce}{\mathsf{mintNonce}}
\newcommand{\burnNonce}{\mathsf{burnCoinNonce}}
\newcommand{\changeNonce}{\mathsf{changeNonce}}

% --- identifiers exactly as they appear in the Compact and Solidity source ---
\newcommand{\srcStartDeposit}{\code{startDeposit}}
\newcommand{\srcStartSupply}{\code{startSupply}}
\newcommand{\srcStartRedeem}{\code{startRedeem}}
\newcommand{\srcFailGate}{\code{assertAttestedFailureOutput}}
\newcommand{\srcReqMap}{\code{signBidirectionalEventMap}}
\newcommand{\srcDepositMap}{\code{depositEventMap}}
\newcommand{\srcSwapMap}{\code{swapEventMap}}
\newcommand{\srcSupplyMap}{\code{supplyEventMap}}
\newcommand{\srcRedeemMap}{\code{redeemEventMap}}
\newcommand{\srcDepositView}{\code{depositSettleViews}}
\newcommand{\srcWdView}{\code{withdrawSettleViews}}
\newcommand{\srcSwapView}{\code{swapSettleViews}}
\newcommand{\srcSupplyView}{\code{supplySettleViews}}
\newcommand{\srcRedeemView}{\code{redeemSettleViews}}
\newcommand{\srcInitFlag}{\code{initialised}}
\newcommand{\srcUnderlying}{\code{stataUnderlying}}
\newcommand{\srcWrapper}{\code{stataToken}}
\newcommand{\srcDepositSel}{\code{deposit(uint256,address)}}
\newcommand{\srcRedeemSel}{\code{redeem(uint256,address,address)}}


\begin{document}

\title{SigNet's Public-to-Private Cross-Chain \Payment System\\[2pt]
  {\large\mdseries Formal Treatment of its Security Analysis}}

\maketitle
\noindent{\color{SigInkSoft}\footnotesize Alessandra Scafuro, \href{mailto:alessandra@sig.network}{alessandra@sig.network}}\par\vspace{\baselineskip}
\begin{abstract}
    This document establishes the foundation for analyzing the soundness and unlinkability of SigNet’s Public-to-Private Cross-Chain System.  The system comprises several complex components: the MPC systems, the Midnight  shielded-transaction layer and shielded-computation layer, the Ethereum ledger.
    After establishing the foundations, this document  focuses exclusively on the  \href{https://github.com/sig-net/midnight-examples/blob/erc20-vault-v0.3.0/examples/erc20-vault/contract/src/erc20-vault.compact}{vault contract} \emph{Version 0.1}\footnote{ This is an   \href{https://github.com/sig-net/midnight-examples/blob/erc20-vault-v0.3.0/examples/erc20-vault/contract/src/erc20-vault.compact}{\textbf{earlier version}} of the  system. A  \href{https://github.com/sig-net/midnight-examples/blob/dev/examples/erc20-vault/contract/src/erc20-vault.compact}{new version} is available and contains modifications that are not formally analyzed here. Nevertheless, the proof strategy used in this paper should extend to the new version. }. 
    (For a high-level description of the system logic consult the
   \href{ https://github.com/sig-net/midnight-examples/blob/dev/examples/erc20-vault/README.md}{ERC20 Vault readme} available on our repository.)
 
     Since the vault is written in compact language (provided by Midnight) and it relies on primitives provided by the Midnight shielded-transaction layer, our proof explicitly identifies the cryptographic properties that we assume from the Midnight components, though we do not formally prove that the Mindnight components achieve them.
    The document is organized as follows.
    
    \begin{enumerate}
        \item   Section~\ref{sec:definition}: Formal definition of  Public-to-Private Cross-chain Payment System.
       This section establishes a formal model for a Public-to-Private Cross-Chain \Payment System over two abstract ledgers: a public ledger ($\Lpub$) and a private ledger ($\Lshd$). It defines core operations (e.g., Deposit, Withdraw, Swap, Supply, Redeem) and their inputs/outputs description, providing a generalized framework adaptable to any pair of chains.
        
        \item  Section~\ref{ssec:unlinakbility}: Formal definition of \emph{Unlinkability} in a \emph{Public-to-Private} Cross-chain \Payment System. This section expresses  the  unlinkability that a public-to-private chain system can achieve via a Game-Based definition, standard in cryptographic proofs. We preferred the game-based approach over the simulation-based approach in the Universal Composability Framework for simplicity. We believe that the game clearly identifies what is protected (and what is leaked) in a public-to-private chain system.
        \item  Section~\ref{sec:components} and Section~\ref{sec:instantiation}:  Implementation. We provide  an abridged   description of the  \href{https://github.com/sig-net/midnight-examples/blob/erc20-vault-v0.3.0/examples/erc20-vault/contract/src/erc20-vault.compact}{vault contract} (Version 0.1). 
      Our implementation  uses  shielded-transaction layer and shielded-computation layer  provided by Midnight. When useful, we use \MidNoteT{blue-shaded boxes} to highlight that we are using a midnight component, and therefore we rely on its security and soundness.
      Furthermore, our contract uses the SygNet's MPC signing system that is treated here as a \emph{trusted signing system}. We also assume that all components related to emitting and processing chain events are trusted.
        \item Section~\ref{sec:assumptions}: Assumptions. We state  the cryptographic assumptions we use and the ones that we inherit from the     Midnight private circuit-layer  and shielded-coin layer.
    
        \item  Section~\ref{sec:proof}: A semi-formal analysis of the security of implementation. We prove that our cross-chain system achieves the unlikability properties (under the assumption that the network is anonymized) and  the soundness  property.
        In each section, we state the model and cryptographic assumptions we require for the proof. 

        \item Section~\ref{sec:conclusions}: Conclusions. Finally,  we distill essential  takeaways from this treatment. We explain what our analysis means for a user.
    \end{enumerate} 

\end{abstract}
\section{Version analyzed in this document }
This document analyses an  earlier version of the system (\href{https://github.com/sig-net/midnight-examples/blob/erc20-vault-v0.3.0/examples/erc20-vault/contract/src/erc20-vault.compact}{link})
 which   used a  counter-based  approach to prevent replay attacks.

The new version (\href{https://github.com/sig-net/midnight-examples/blob/dev/examples/erc20-vault/contract/src/erc20-vault.compact}{link}) addresses replay attacks with a new design based on timing information (updating the last-seen block), and introduces a mechanism to process many requests in parallel,  while preventing contention. 

The new version introduces public procedures (circuits) to organize requests sent to the MPC and attestations received from it.  These  procedures do not use any private inputs,  can be executed by anyone and are anonymous (i.e., the identity of the caller is not leaked by the circuit computation). As such, they do  not affect the rationale behind the unlinkability  proofs provided in this document.
The soundness approach is expected to remain unchanged, but this must be confirmed with a formal security proof.

\vspace{5pt}
In the document, we will use  the paragraph-header: ``
\newversion{...}''
to  describe the differences with the new version.

\newpage

\section{Private-Public Cross-chain \Payment System: Formal Definition}
\label{sec:definition}
This section describes the functionality that  a Private-Public Cross-chain \Payment  System provides formally. The system is defined in terms of the commands it provide. Each command is defined via the input-output interface.


A cross-chain \payment system is a bridge between a public  EVM-ledger  $\Lpub$ (e.g., Ethereum) and a private (shielded) ledger $\Lshd$ (e.g., Midnight). The users of the cross-chain system are users of possess account on both chains. The system and the users interact with each chain using the available interface.


\begin{definition}[Public-to-Private Cross-chain \Payment System]
\label{def:system}
A \emph{public-to-private cross-chain defi system} $\Pi$, relative to $\Lpub$ and $\Lshd$,
consists of the following PPT procedures.
Every procedure has oracle access to $\Lpub$ and $\Lshd$.

Each procedure outputs the following  tuple:
\[
\bigl(\ \stpub',\ \stshd',\ \trans, \outc; \ \prvOut\ \bigr),
\]
where  $\stpub',\ \stshd'$ are the resulting state  updates for  each ledger; $\trans$ is the public transcript the execution of the call produces; 
 $\outc$ indicates if the function terminates  successfully, or it failed (functions are not guaranteed to terminate)  and $\prvOut$ is a  private output the caller retains.

\begin{enumerate}[leftmargin=1.8em,itemsep=6pt]
\item $(\pp, \setupstate) \gets \Pi.\Setup^{\Lpub, \Lshd}(1^{\secpar})$\\
  Initialises the system on both ledgers and outputs public parameters $\pp$. Run once, by a
  trusted party.  \textbf{$\pp$ is an implicit input to every procedure below,} and fixes the rules and  parameters of the cross-chain payment system.

\item $\Pi.\Deposit^{\Lpub, \Lshd}\bigl([\pubAddr]\bc \shdAcct\bc \tokType\bc \amt\bc
      \auxx \,;\, \privIn\bigr)$.\\
  Called by the holder of the  $\shdAcct$, whose secret is in $\privIn$, $\auxx$ carries execution parameters of the chains. 
  On $\succeeded$: $\amt$ units of $\tokType$  is moved to a vault on $\Lpub$ and credited to 
  $\shdAcct$.  On $\failed$ no operation is performed. 
  $\pubAddr$ is the address on $\Lpub$ that funded the deposit to the vault. The parenthesis [] emphasize that $\shdAcct$ doesn't need to be the holder of  $\pubAddr$ -- though in the practical instantiation it is~\footnote{In our implementation, an address $\pubAddr$ must send funds to an intermediate, application-derived account, that is controlled by the Bridge. Such intermediate account is the used to deposit the amount $\amt$ to the vault address. We explicitly list $\pubAddr$ as a public genesis address (independent on the implementation), to stress that the link between the accounts on the two ledgers is revealed by a deposit}.
\item $\Pi.\Withdraw^{\Lpub, \Lshd}\bigl(\shdAcct\bc \destAddr*\bc \tokType\bc \amt\bc
      \auxx \,;\, \privIn\bigr)$.\\
  Called by the holder of the shielded-chain account $\srcAcct$, whose secret is in $\privIn$.
  If $\shdAcct$ owns at least $\amt$ units of $\tokType$ on the vault,  $\amt$ units of $\tokType$ leave $\srcAcct$ on $\Lshd$ and $\amt$ units are
  credited to $\destAddr$ on $\Lpub$. On $\failed$: $\shdAcct$'s balance is unchanged.


  \item $\Pi.\Swap^{\Lpub, \Lshd}\bigl(\shdAcct\bc \tokType\bc \amt\bc
    \tokTypeB\bc \amtB\bc \auxx \,;\, \privIn\bigr)$.\\
  Called by the holder of $\srcAcct$. On $\succeeded$: $\amt$ units of $\tokType$ are swapped with $\amt'$ units  of $\tokType'$.
  On $\failed$: $\shdAcct$'s balance is unchanged.
 


\item $\Pi.\Supply^{\Lpub, \Lshd}\bigl(\shdAcct\bc [\tokType,\tokTypeB]\bc \amt\bc \auxx \,;\, \privIn\bigr)$.\\
  Called by the holder of the shielded-chain account $\shdAcct$, whose secret is in $\privIn$.
  If $\shdAcct$ owns at least $\amt$ units of $\tokType$ on the vault, then on $\succeeded$:
  $\amt$ units of $\tokType$ leave $\shdAcct$ on $\Lshd$, the vault supplies them to the lending
  market on $\Lpub$, and $\amtShares$ units of $\tokTypeB$ are credited to $\shdAcct$.
  On $\failed$: $\shdAcct$'s balance is unchanged.
  \smallskip
  $\amtShares$ is determined by the state of
  $\Lpub$ at the moment of execution,  and
  is reported back in $\trans$. 
  [$\tokType,\tokTypeB$] are optional parameters (our implementation of $\Supply$ they are fixed and defined in the initial parameters $\pp$.).

\item $\Pi.\Redeem^{\Lpub, \Lshd}\bigl(\shdAcct\bc [\tokType']\bc \amtShares\bc \auxx
      \,;\, \privIn\bigr)$.\\
  Called by the holder of $\shdAcct$. If $\shdAcct$ owns at least $\amtShares$ units of
  $\tokType'$ on the vault, then on $\succeeded$: $\amtShares$ units of $\tokType'$
  leave $\shdAcct$ on $\Lshd$, the vault redeems them on $\Lpub$, and $\amtAssets$ units of
  $\tokType$  (principal plus accrued interest) are credited to $\shdAcct$.
  On $\failed$: $\shdAcct$'s balance is unchanged.
\end{enumerate}
\end{definition}


\section{Definition of Unlinkability.}
\label{ssec:unlinakbility}
A Public-to-Private  Cross-chain \Payment  System is designed for actions 
to be initiated on a private chain and to have effect on a vault on the  public chain.
The latter, requires that amounts of the transactions be revealed. With that in mind, there are two notions of unlinkability that one could aim for:
\begin{itemize}
    \item \emph{Transaction unlinkability}: can an observer  tie a
withdrawal to a user? This connection could depend on amounts and on the order of events, and cannot be a property of a protocol: if the public chain records that the vault received  only two deposits, one of  $10$ from $\pubAddr_{A}$  and one deposit of $6$ from $\pubAddr_{B}$ later  a withdrawal of $7$  to a  fresh new account  $\pubAddr^*$,  leaks 
that the withdrawal request came from the user controlling $\pubAddr_{A}$, regardless of the cryptography involved in the withdrawal process.
\item \emph{Cryptographic  unlinkability}:  do the protocol messages give any \emph{additional} help, besides what can be inferred from the public amounts displayed on the public chain?
\end{itemize} 

In this section we target the latter.
We provide   the notion of \emph{Withdrawer's Unlinkability} that states that it should be computationally infeasible from any PPT adversary to guess the identity of the 
shielded account associated to a withdrawal, even given full knowledge of the deposit and withdrawal/swaps/supply/redeems performed by all honest users (and even actively participating as a malicious user in the system). We capture this   definition with a formal cryptographic game,  the ithdrawer unlinkability game $\unlinkgame_{\Pi,\Adv}(1^{\secpar})$.

In the experiment, there is a challenger that  (honestly) simulates the  cross-chain system and the two ledgers for the adversary. The adversary  can choose the operations that are executed  in the cross-chain system: it can create honest accounts, and make them perform any action (i.e., deposits, withdrawals,  swaps, supply redeem)  
on amount and types of her choice.
In the formal game this is modeled by the fact that the adversary has oracle access to  each action provided by the system and by the two chains. The oracles are described in Section~\ref{ssec:oracles}.

At the challenge phase, the adversary selects two shielded accounts and one amount $v$.
If both accounts have at least a balance of $v$,  the challenger will choose one at random and run the withdrawer function function on it, revealing the resulting transcript to the adversary.
The goal of the adversary now is to guess which bit was chosen by the challenger.
Before outputting his guess, the adversary can continue interact with the  system arbitrarily, and the system will respond to the request  \emph{under the condition that both accounts in the challenge phase  have a $-v$ amount in their balance.}
The adversary wins the game if it guesses correctly with a probability better than half.

\vspace{5pt}
\noindent
\textbf{The power of the adversary.}
We consider an adversary that has full view of all public protocol messages and 
Ledger's messages.
Furthermore, we consider an adversary that can have any type of side channel information about the users. For instance, the adversary could know that 
a certain action (e.g., Withdraw, Deposit, Supply, etc) was issued by a specific user.
Formally, we model this knowledge by giving to the adversary oracle access
to all the calls $\Deposit$, $\Withdraw$, $\Swap$, $\Supply$, $\Redeem$ provided by our interface. The adversary can decide which users perform which action and choose the amount,.
Furthermore, the adversary is always aware of all operations that take place to each ledger. For the shielded ledger, the adversary could have side channel information about the user who is performing a certain shielded transaction at a certain point in time.
Once again, we model this knowledge by giving the adversary the power to choose which user performs which operations on each chain. This power is modeled by the Oracles (\S~\ref{ssec:oracles}).

\begin{remark}[No Network Information]
We assume an adversary that collects information solely at the application layer and explicitly exclude network-layer threats (e.g., user IP addresses, routing metadata). Since network-level data can deanonymize messages regardless of application-layer cryptographic protections, restricting the adversary's scope to the application layer is a standard assumption in privacy-preserving systems.
\end{remark}

\subsection{Oracles}
\label{ssec:oracles}
As mentioned above, in our security model we consider an adversary who could have side channel information about certain users taking certain actions at certain time, even in the shielded ledger.
To formally modeled this in our security experiment, we give the adversary access to oracles to call actions for the cross-chain system, and oracles to the ledgers.
In the security experiment,  the \emph{challenger} simulates the oracle cross-chain system and the ledgers  for the adversary.
The challenger $\Chal$ maintains the relevant states to simulate the various systems,  and in particular, it maintains a balance sheet $\bal[\cdot][\cdot]$, recording for every account
it created how much of each type that account holds, updated on every call; both ledgers; and the
transcript $\Tall = (\Tpub, \Tshd)$. Every oracle call appends to $\Tall$ everything the two
ledgers publish because of it, and $\Adv$ reads $\Tall$ at all times.

\medskip
\subsubsection{Ledger oracles.}
\label{sec:ledgeroracles}
The adversary $\Adv$  has the power to create new ledger users (for both public and private ledger), and to instruct any operation allowed on each chain and observe the result of those activities.
This power is captured by giving the adversary oracle calls to
actions  (e.g., Create, Mint, Transfer, etc.)  over $\Lpub$ and $\Lshd$.
Below we describe the ledger oracles provided to the adversary.


\medskip
\noindent\textbf{Oracles for the public ledger $\Lpub$},
\begin{itemize}[leftmargin=1.4em,itemsep=3pt]
\item $\ONewPub() \to \pubAddr_i$: it allows to create new account on the public chain. The ledger
      creates it, $\Chal$ retains the account secret, and the public address created $\pubAddr_i$ is returned to $\Adv$.

\item $\OTransPub(\pubAddr, \pubAddr', \tokType, \amt)$ it allows to request a  public-chain transfer from  $\pubAddr$ to  $\pubAddr'$ of $\amt$ of $\tokType$.
 \item $\Lpub.\LExecute(\pubAddr, \contract, \contractinput)$: moves $\amt$ units of $\tokType$
\item $\Tpub \gets \Lpub.\LRead()$ It returns the full history of the ledger.
\end{itemize}

 \medskip
\noindent\textbf{Oracles for the shielded ledger $\Lshd$},
\begin{itemize}[leftmargin=1.4em,itemsep=3pt]
\item $\ONewShd() \to \shdAcct_i$, it allows to create new account on the shielded chain. The ledger
      creates it, $\Chal$ retains the account secret, and the public part of the shielded address $\shdAcct_i$ is returned to $\Adv$.
\item $\OMint(\shdAcct_i, \tokType, \amt)$: it allows to  instruct a transfer $\amt$ to the private pool, for account $\shdAcct_i$, on $\Lshd$. The secret associated to this call are maintained by the challenger. 

\item $\OTransShd(\shdAcct_i, \shdAcct_j, \tokType, \amt)$: this oracle allows a \emph{public} transfer on $\Lshd$. 
\item $\OShdTrans(\shdAcct, \shdAcct', \tokType, \amt)$: this oracle allows  $\Adv$ to request a \emph{shielded} transfer on $\Lshd$ on accounts of their choice $\shdAcct, \shdAcct'$.   

\item $\Lshd.\LExecute(\shdAcct_i, \contract,  \contractinput)$: this oracle call allows  $\Adv$ to request execution of $\contract$ on behalf of  $\shdAcct_i$  on public input $\contractinput$. The private input is maintained by the challenger.
 % Let $\Leak(\shdAcct_i,\cdot,\cdot)$ be the leakage output from such a call (defined in Definition~\ref{def:fingerprint}).
\item $\Tshd \gets \Lshd.\LRead()$ It returns the full history of the ledger.
\end{itemize}
\medskip
\noindent Separately, $\Adv$ may create and operate accounts of its own directly through the ledger
interfaces, on either chain, and post its own transactions. The
oracles above exist so that $\Adv$ can also direct accounts whose secrets it does \emph{not} hold.
We call an account \emph{honest} if it was created by the oracles.


\subsubsection{Cross-chain oracles.}
\label{sec:crosschainoracles}
Using the cross-chain oracles the adversary can orchestrate the interactions between honest parties, choosing which parties need to perform which operations, choosing amounts and auxiliary information $\auxx$ (such as gas prices etc).
The cross-chain oracles are described below.
The adversary decides the inputs on which the oracles to be run on.
The challenger then honestly executes  those calls on behalf of honest users (i.e., the challenge knows the private state of honest users) and the resulting transcripts are returned to the adversary.

\begin{itemize}[leftmargin=1.4em,itemsep=3pt]
\item $\ODep(\pubAddr, \shdAcct, \tokType, \amt, \auxx, \outc)$:
The challenger runs $\Deposit$ on behalf of $\shdAcct$, and using  $\pubAddr$ as the source account on the public chain.
On $\succeeded$, add $\amt$ to $\bal[\shdAcct][\tokType]$ 
and update the ledger balance of  $\pubAddr$. 
\item $\OWd(\shdAcct, \destAddr*, \tokType, \amt)$: runs $\Withdraw$ on behalf of
      the holder of $\shdAcct$, provided $\bal[\shdAcct][\tokType] \ge \amt$. On $\succeeded$,
      moves $\amt$ from $\bal[\shdAcct][\tokType]$ to $\ledgerBal[\pubAddr][\tokType]$.
\item $\OSw(\shdAcct\bc \tokType\bc \amt\bc \recvAcct\bc
      \tokTypeB\bc \amtB\bc \auxx)$: runs $\Swap$
      on behalf of the holder of $\shdAcct$, provided $\bal[\shdAcct][\tokType] \ge \amt$. On
      $\succeeded$, subtracts $\amt$ from $\bal[\shdAcct][\tokType]$ and adds $\amtB$ to
      $\bal[\recvAcct][\tokTypeB]$.
\item $\OSup(\shdAcct_i\bc \tokType\bc \amt\bc \auxx\bc \outc)$: runs $\Supply$ on behalf of the
      holder of $\shdAcct$, provided $\bal[\shdAcct][\tokType] \ge \amt$. On $\succeeded$,
      subtracts $\amt$ from $\bal[\shdAcct]$ $[\tokType]$ and adds $\amtShares$ to
      $\bal[\shdAcct][\tokType']$, where $\amtShares$ is the amount $\Lpub$ returned and is
      recorded in $\Tall$, and $\tokType'$ is the exchange token
      defined by the $\Supply$ protocol. On $\failed$, no balance changes.
\item $\ORed(\shdAcct_i, \bc \amtShares\bc \auxx\bc \outc)$: runs $\Redeem$ on
      behalf of the holder of $\shdAcct$, provided
      $\bal[\shdAcct][\tokTypeB] \ge \amtShares$. On $\succeeded$, subtracts $\amtShares$
      from $\bal[\shdAcct][\tokTypeB]$ and adds $\amtAssets$ to
      $\bal[\shdAcct][\tokType]$, with $\amtAssets$ as returned by $\Lpub$ and recorded in $\Tall$.
      On $\failed$, no balance changes.

      \textbf{Note.} The $\amtAssets$ is a parameter dictated by $\Lpub$ (and not an input to the oracle). This reflect the scenario where the exchange rate is defined by the market, and it is independent  on by the user who attempts to redeem.

      $\diamond$ Future  definitions could model a scenario where  additional criteria are used (e.g., a user must prove that it belong to a whitelist)
\end{itemize}

\noindent A call whose stated balance condition fails is not executed and $\Chal$ returns $\bot$.


\subsection{Withdrawer unlinkability Definition}
\label{ssec:unlink}


\subsubsection{Experiment \texorpdfstring{$\unlinkgame$}{Exp-Unlink}}
\begin{ExpFig}{fig:unlink}{$\unlinkgame_{\Pi,\Adv}(1^{\secpar})$}

\vspace{10pt}
$\Adv$ has access to  the ledger oracles  and cross-chain oracle defined in \S\ref{ssec:oracles} and simulated by the challenger.
Let  $\Adv^{\mathcal{O}(\cdot)}$ denote the adversary having access to 
such oracles.


\begin{itemize}[leftmargin=1.4em,itemsep=3pt]
\item \textbf{Initialization.} $(\pp,\setupstate) \gets \Setup^{\Lpub,\Lshd}(1^{\secpar})$; $\bal := 0$;
      $\Tall := \emptyset$. $\Adv$ is given $\pp$.
\item \textbf{Pre-challenge Training.}
$\Adv$ has access to the oracles  described above:  
      \[(\candZero\bc \candOne\bc \tokType\bc \amt\bc \auxx^*\bc \outc^*\bc
      \mathsf{st}) \gets \Adv^{\mathcal{O}}(\pp),\]
      where $\candZero, \candOne$ are honest shielded accounts (created through the $\ONewShd$ oracle.
\item \textbf{Challenge}
  \begin{itemize}[nosep,leftmargin=1.4em]
  \item if $\candZero$ or $\candOne$ is not honest, output $0$; // C1. if one user is controlled by the adversary, distinguishing is trivial
  \item if $\bal[\candZero][\tokType] < \amt$ or $\bal[\candOne][\tokType] < \amt$, output $0$;
  // If the amounts are different, the adversary can easily tell from amount posted on the public chain $\Lpub$.
   % \item if $\Leak(\candZero, \tokType, \amt) \neq \Leak(\candOne, \tokType, \amt)$, output $0$.
   %  // If the spending fingerprint of the two users is different, the adversary could distinguish what user was selected.
   
  \item otherwise: 
  \begin{enumerate}
      \item sample $b \gets \{0,1\}$;
      
      \item create a fresh  public-chain address $\destAddr^*$ and execute
      $$\Withdraw(\candB, \destAddr^*, \tokType, \amt, \auxx \,;\, \privIn_b) $$

   with the output  $$\stpub*',\ \stshd*',\ \trans^*, \outc^*$$ returned to $\Adv$.     
  \item if $\outc^* =\succeeded$:\\ Update the balance state for both account:\\
  $\bal[\candZero][\tokType] \mathrel{-}= \amt$ \emph{and}
        $\bal[\candOne][\tokType] \mathrel{-}= \amt$; \\
        and use these amounts to simulate all the oracles \S~\ref{sec:crosschainoracles},\S~\ref{sec:ledgeroracles}.
    \item if $\outc =\failed$ balances are not updated.     
  \end{enumerate}
    \end{itemize}
\item \textbf{Post-Challenge Training.} $\Adv$ continues to have access to all the  oracles above, with $\Chal$ answering every cross-chain oracle calls against the balance sheet $\bal$ updated   in the challenge phase.


$$b' \gets \Adv^{\mathcal{O}}(\mathsf{st}, \Tall)$$
\item \textbf{Output} The  output of the experiment is $1$ iff $b = b'$.
\end{itemize}
\end{ExpFig}

\begin{definition}[Withdrawer unlinkability]
\label{def:unlink}
A cross-chain payment protocol has $\Pi$ satisfies \emph{withdrawer unlinkability} if for all PPT $\Adv$ there is a negligible $\negl$ such
that
$$\Pr[\unlinkgame_{\Pi,\Adv}(1^{\secpar}) = 1] \le \frac{1}{2} + \negl(\secpar)$$
\end{definition}


\begin{remark}[Pending Requests.]
\label{rem:failpending}
In the Unlinkability Experiment  we assume that the cross-chain oracle calls that involve \emph{honest parties}  always terminate (either with success or failure).
% \ale{Add discussion about unviable transactions}
\end{remark}


\subsection{ Swapper, Supplier, Swapper and Redeemer Unlinkability}
\label{sssec:lendgame}
The exact same experiment can be applied to every action of the  cross-chain transaction system.
Write $\act \in \{\Swap, \Supply, \Redeem\}$, and for a given action $\act$ let $\tokTypeIn$ and $\tokTypeOut$ be the input and output types. 
 Let $\amt$ be the input amount and let
$\amt^{\mathrm{out}}$ denote the amount $\Lpub$ returns, which the challenger learns when it runs the call.
We provide general  the experiment in Figure~\ref{fig:unlinkGeneral}.


\begin{ExpFig}{fig:unlinkGeneral}{$\unlinkgameG_{\Pi,\Adv}(1^{\secpar})$}

\vspace{10pt}
$\Adv$ has access to all oracles of \S\ref{ssec:oracles}.

\begin{itemize}[leftmargin=1.4em,itemsep=3pt]
\item \textbf{Initialization.} $(\pp,\setupstate) \gets \Setup^{\Lpub,\Lshd}(1^{\secpar})$;
      $\bal := 0$; $\Tall := \emptyset$. $\Adv$ is given $\pp$.
\item \textbf{Pre-challenge Training.} $\Adv$ has access to the cross-chain oracles described above
      (as well as the ledger oracles).
      \[(\act\bc \candZero\bc \candOne\bc \tokTypeIn\bc \amt\bc \tokTypeOut\bc\auxx^*\bc \outc^*\bc
      \mathsf{st}) \gets \Adv^{\mathcal{O}}(\pp),\]
      where $\act \in \{\Withdraw, \Swap, \Supply, \Redeem\}$, the parameters $\parm$ are
      those the chosen action takes, and $\candZero, \candOne$ are honest shielded accounts.

\item \textbf{Challenge}
  \begin{itemize}[nosep,leftmargin=1.4em]
  \item if $\candZero$ or $\candOne$ is not honest, output $0$;
  \item if $\bal[\candZero][\colIn] < \valIn$ or $\bal[\candOne][\colIn] < \valIn$, output $0$;
%  \item if $\Leak(\candZero, \colIn, \valIn) \neq \Leak(\candOne, \colIn, \valIn)$, output $0$;
  \item otherwise:
  \begin{enumerate}
      \item sample $b \gets \{0,1\}$;
      \item if $\act = \Withdraw$, create a fresh public-chain address $\destAddr^*$ and
            include it in $\parm$;
      \item execute $$\act(\candB\bc \parm\bc \auxx^* \,;\, \privIn_b)$$
            with the output $$\stpub*',\ \stshd*',\ \trans^*,\ \outc^*$$ returned to $\Adv$;

    \item 
     Let $\mathsf{st}^{(0)}$ and $\mathsf{st}^{(1)}$ denote the states of the honest
      parties in the two worlds: $\mathsf{st}^{(\beta)}$ is the state that results from
      running $\act$ on behalf of $\mathsf{shAcct}_{\beta}$. The challenger executes the
      challenge in world $b$ and \textbf{retains both} $\mathsf{st}^{(0)}$ and
      $\mathsf{st}^{(1)}$.
  \end{enumerate}
  \end{itemize}
\item \textbf{Post-Challenge Training.} $\Adv$ continues to have access to all the oracles above,  with $\Chal$ answering every cross-chain oracle call as follow:
      \begin{itemize}[nosep,leftmargin=1.4em,topsep=2pt]
      \item if $q$ is executable in \emph{both} $\mathsf{st}^{(0)}$ and
            $\mathsf{st}^{(1)}$, execute it in world $b$ and return the resulting
            transcript;
      \item otherwise return $\bot$.
      \end{itemize}
      $$b' \gets \Adv^{\mathcal{O}}(\mathsf{st}, \Tall)$$
\item \textbf{Output.} The output of the experiment is $1$ iff $b = b'$.
\end{itemize}
\end{ExpFig}

\begin{definition}[Action Unlinkability]
\label{def:lendunlink}
$\Pi$ has \emph{action unlinkability} if for every $\act \in \{\Swap,\Supply,\Redeem\}$ and
every PPT $\Adv$ there is a negligible $\negl$ with
$$\Pr[\lendunlinkgame_{\Pi,\Adv}(1^{\secpar}) = 1] \le \frac12 + \negl(\secpar).$$
\end{definition}



\section{Our Implementation: Components}
\label{sec:components}

The implementation is the ERC20 Vault contract.  Consult the \href{\analysedReadmeUrl}{ERC20 Vault README}  for a
high-level description of the contract logic and components.


\subsection{System Components}
\noindent
\textbf{Trusted Bridge $\SigSys$'s MPC.}
The instantiation of the cross-chain system relies on $\SigSys$'s MPC system acting as
a trusted bridge that listen for events (actions such as $\Deposit,\Swap,$ etc) on  Midnight
and executed them on Ethereum.
Between a request   (e.g., $\Deposit$) and its
settlement  circuit, three steps run outside any proof : the committee signs the requested EVM transaction;
a relayer broadcasts it; the committee observes the result and posts a signed attestation of it back
to Midnight.


$\SigSys$'s MPC  Signing System  is trusted to satisfy:
    \begin{enumerate}
        \item  \textbf{MPC committee integrity:}  the committee signs only what the protocol authorises.
        \item \textbf{Attestation faithfulness:}  MPC only provide attestations for events that have been executed on $\Lpub$ (Ethereum). The MPC attests failed only if the transaction did not execute.
    \end{enumerate}.

\begin{center}
\begin{tabular}{@{}L{2.7cm}L{5.4cm}L{6.3cm}@{}}
\toprule
\textbf{Definition} & \textbf{Instantiated by} & \textbf{Notes} \\
\midrule
$\Lpub$ &  Ethereum & accounts are EVM addresses; types are ERC20
  contracts; everything public \\
  \midrule
$\Lshd$ & Midnight  &  enables private payments  and execution of contract with private inputs.\\
\midrule
$\SigSys$  & SigNetwork MPC &  It manages the vault account on Ethereum and provides the attestations for it.\\
% \midrule
% $\Setup$ & deploy, $\cktInit$, $\cktApprove$ & \S\ref{ssec:setupinst} \\
% $\Deposit$ & $\cktDeposit$ $\to$ MPC $\to$ $\cktCompleteDeposit$ & \S\ref{ssec:depinst} \\
% $\Withdraw$ & $\cktWithdraw$ $\to$ MPC $\to$ $\cktCompleteWithdraw$ or $\cktRefund$ &
%   \S\ref{ssec:wdinst} \\
% $\Swap$ & $\cktSwap$ $\to$ MPC $\to$ $\cktCompleteSwap$ or $\cktRefund$ & \S\ref{ssec:swapinst} \\
\bottomrule
\end{tabular}
\end{center}


% $\Setup$ additionally publishes $\SigSys$'s public key ${\sf PK}_\SigSys$ in $\pp$.
% \medskip


\subsubsection{Key derivation}
\label{ssec:derive}

$\SigSys$ holds a single root key pair $(\mathsf{sk}_{\mathsf{root}}, \mpcRootPk)$ on
secp256k1, with $\mpcRootPk = \mathsf{sk}_{\mathsf{root}} \cdot G$ for the group generator
$G$ of order $n$. Every account it signs with is an \emph{additively derived} child of that
root, scoped by the requesting contract and by a $32$-byte path.


\begin{definition}[Additive key derivation]
\label{def:derive}
Let $\varepsilon : \{0,1\}^{*} \to \mathbb{Z}_n$ be
\[
  \Epsilon(\vaultAddr, \keyPath) \;=\;
  \Kec\bigl(\, \text{prefix} \,\|\, \chainId \,\|\, \vaultAddr \,\|\, \keyPath \,\bigr)
  \bmod n ,
\]
The derived key pair is
\begin{flalign*}
  &\mathsf{sk}_{\mathsf{der}} \;=\; \mathsf{sk}_{\mathsf{root}} + \varepsilon(\vaultAddr, \keyPath) \bmod n ,& \\
  &\Derive(\mpcRootPk\bc \vaultAddr\bc \keyPath) \;=\; \mpcRootPk + \Epsilon(\vaultAddr, \keyPath) \cdot G .&
\end{flalign*}
Where an EVM account is wanted, $\derivedAddr$ is the standard address encoding of the
derived public key: the low $20$ bytes of $\Kec$ applied to its uncompressed encoding.
\end{definition}

\noindent Three properties of Definition~\ref{def:derive} are used throughout.

\begin{itemize}[leftmargin=1.4em,itemsep=3pt,topsep=3pt]
\item \textbf{Public computability.} $\Derive$ takes only public inputs and addition on the
      curve is public, so \emph{anyone} can compute $\derivedAddr$ from $\mpcRootPk$,
      $\vaultAddr$ and $\keyPath$. The private half requires $\mathsf{sk}_{\mathsf{root}}$
      and therefore exists only inside $\SigSys$. This is what makes a deposit's
      $\keyPath = \userComm(\usersk)$ sufficient to expose the depositor's EVM address to an
      observer of $\Lshd$ alone.
\item \textbf{Contract scoping.} $\vaultAddr$ enters the derivation string, so no contract
      can reach another contract's derived accounts.
\item \textbf{Path separation.} Within one contract, distinct paths give unrelated keys.
      The vault uses exactly three: $\userComm(\usersk)$ for a depositor's account,
      the constant $\keyPathLit$ for the vault's own pooled account, and the fixed string
      $\text{``midnight response key''}$ for $\mpcRespKey$, which signs attestations and
      never signs transactions.
\end{itemize}

\noindent The key version $\keyVer$ selects which root key is used and is carried in every
request record; it is suppressed above, and throughout the document, since a single version
is in use.


\medskip
\noindent
\textbf{Derived Addresses:}

The value $\keyPath$ can be:
\begin{itemize}
    \item User address: $\derivedAddr$ uses $\keyPath=\userComm(\usersk)$ (defined below)
    \item Vault address:  $\vaultEvmAddr$  uses $\keyPath=\mathtt{''vault''}$
    \item  MPC Response Address: $\mpcRespKey$  uses $\keyPath=\mathtt{''midnight response key''}$. The attestation key $\mpcRespKey$ is derived the
same way with a fixed protocol string and pinned at $\cktInit$.
\end{itemize}


\subsection{User's Keys}
\label{ssec:userskeys}
A user  $U_i$ of  the cross-chain system holds:
\begin{itemize}
    \item { \bf Application secret key:} $\usersk_i$: a private key for the cross-chain system: $\usersk_i\leftarrow\{0,1\}^{\lambda}$. This is randomly chosen upon the first Deposit request.
      \item \textbf{Public account }$\pubAddr_i$  on $\Lpub$; (used to fund the derived account).
      \item \textbf{Shielded Account } ($\callerZKey_i$, $\callerZKeySk_i$) on $\Lshd$. \\
      In Midnight a shielded account consists of two public keys:
\begin{enumerate}
    \item \textbf{Coin Key-Pair} (Spending Key):
    \begin{enumerate}
        \item $\coinsk$ = $\SHA\bigl(\text{``midnight:csk''} \concat \seed\bigr)$
        \item $\coinpk = \SHA\bigl(\text{``midnight:zswap-pk[v1]''} \concat \coinsk\bigr)$
    \end{enumerate}

    \item \textbf{Encryption Key-Pair} (Discovery key):
    \begin{enumerate}
          \item $\encsk$: a group element generated via Poseidon hash function.
        \item $\encpk = \gen^{\encsk}$
    \end{enumerate}
\end{enumerate}
\noindent So
\[
  \shdAcct \;=\; \bigl(\,(\coinsk, \coinpk),\;\; (\encsk, \encpk)\,\bigr)
\]

\begin{center}
\begin{tabular}{@{}L{1.7cm}L{1.7cm}L{11.0cm}@{}}
\toprule
\textbf{Key} & \textbf{Exposure} & \textbf{Use} \\
\midrule
$\coinpk$ & public & \textbf{ownership}: used as a rcpt when Minting \\
$\coinsk$ & secret & \textbf{spending}:  used as sender when Spending \\
$\encpk$  & public & \textbf{payability}: used to compute ciphertext in Spending a coin\\
$\encsk$  & secret & \textbf{scanning} : used to decrypt in order to Receive a coin \\
\bottomrule
\end{tabular}
\end{center}
\end{itemize}


The  circuits  of our contract only use the $\coinpk$ as a parameter to mint shielded coins.
The other keys are used by the Midnight's circuits.


\subsubsection{User Authentication Values}
\label{ssec:derived}
\label{ssec:hashcomputation}
We use the following hash functions:
\begin{itemize}
    \item $\Hash$: A keyed hash function, used to compute our authentication tags. We require that this function satisfies Pseudorandomness (Def.~\ref{def:prf}) and Multi-user Pseudorandomness (Def.~\ref{def:muprf}) in addition to collision-resistance. Recall also that collision-resistance implies one-wayness (pre-image resistance).

    \begin{align*}
\hypertarget{fn:userComm}{\userComm(\usersk)} &= \Hash\bigl(\text{``vault:user:''}\bc \usersk\bigr)
  &&\text{\small identity; the MPC path in Deposit. }\\
\hypertarget{fn:refundComm}{\refundComm(\usersk, \rid)} &= \Hash\bigl(\text{``vault:refund:''}\bc \usersk\bc \rid\bigr)
  &&\text{\small the authentication tag.}
\end{align*}
We instantiate $\Hash$ with the \texttt{transientHash} implemented by Midnight~\footnote{\href{https://github.com/midnightntwrk/midnight-ledger/blob/cc909e042a17695f7ba62fc195a8d770d150b63e/transient-crypto/src/hash.rs}{Lines L75-L81}.}(based on Poseidon~\cite{poseidon}).
Poseidon follows a sponge-like approach and is a good candidate for Multi-user PRF security~\cite{KeyedDuplex17}.

    \item $\Kec$: A collision resistance hash function used to compute digests.
    \begin{align*}
\hypertarget{fn:domSep}{\domSep(\tokType)} &= \Hash\bigl(\text{``erc20:vault:''}\bc \tokType\bigr)
  &&\text{\small one shielded type per ERC20}\\
 \hypertarget{fn:reqIdOf}\rid &= \Kec(\req)
   &&\text{\small over the whole request record (\textbf{except skemas})}\\
%\hypertarget{fn:attDigest}{\attDigest(\rid, \evmOut)} &= \Kec(\rid \concat \evmOut)
%\newversion{}
\end{align*}
\end{itemize}


The \textbf{outcome} the MPC attests, is computed using the following formula:
 \begin{align*}
 \hypertarget{fn:attDigest}{\attDigest(\rid, \evmOut)}  = & \Kec(\rid \concat \evmOut.\kind \| \evmOut.\outlt \| \evmOut.\blockheight)
\end{align*}
  % &&\text{\small what the committee attests}
We instantiate $\Hash$ with the \texttt{transientHash} implemented by Midnight.

\subsection{The vault's public state}
\label{ssec:state}
\begingroup
\setlength{\LTpre}{6pt}\setlength{\LTpost}{6pt}
\begin{longtable}{@{}L{0.01cm}L{2.9cm}L{4.7cm}L{7.5cm}@{}}
\toprule
& \textbf{Here} & \textbf{Source name} & \textbf{Role} \\
\midrule
\endfirsthead
\multicolumn{4}{@{}l}{\small\itshape\ldots continued from previous page}\\
\toprule
& \textbf{Here} & \textbf{Source name} & \textbf{Information Held} \\
\midrule
\endhead
\midrule
\multicolumn{4}{r@{}}{\small\itshape continued on next page \ldots}\\
\endfoot
\bottomrule
\endlastfoot
 & $\reqMap$ & \srcReqMap & withdraw and approve requests, keyed by $\rid$; the MPC reads
  them here.  \\
& $\depositMap$ & \srcDepositMap & deposit requests keyed by $\rid$ \\
& $\swapMap$ & \srcSwapMap & swap requests keyed by $\rid$ \\
& $\supplyMap$ & \srcSupplyMap & supply requests keyed by $\rid$\\
& $\redeemMap$ & \srcRedeemMap & redeem requests  keyed by $\rid$\\
\addlinespace[3pt]
 & $\depositView$ & \srcDepositView & pending-\emph{deposit} record: the depositor
  commitment, the ERC20 and the amount $\cktCompleteDeposit$ mints \\
 & $\wdView$ & \srcWdView & pending-\emph{withdrawal} record: the refund commitment, the
  ERC20 and the amount.  \\
 & $\swapView$ & \srcSwapView & pending-\emph{swap} record: the refund commitment,
  $\tokType$, $\tokTypeB$, $\amtOut$ and $\amtMax$.  \\
& $\supplyView$ & \srcSupplyView & pending-\emph{supply} record: the refund commitment
  and the amount.\\
& $\redeemView$ & \srcRedeemView & pending-\emph{redeem} record: the refund commitment
  and the share count.  \\
\addlinespace[3pt]
& $\mpcRespKey$ & \code{mpcResponseKey} & signature verification key for attestation \\
& $\sigNonce$ & \code{signetRequestNonce} & a counter making otherwise identical requests distinct;
  its value publicly orders every request the vault has made \\
 & $\initFlag$ & \srcInitFlag & one-shot setup flag \\
& $\vaultEvmAddr$ & \code{vaultEvmAddress} & the contract-enforced recipient of every deposit \\
& $\routerAddr$ & \code{uniswapRouter} & the only spender $\cktApprove$ approves, and the only
  destination a swap request sends to \\
& $\underAddr$ & \srcUnderlying & the lent ERC20; gives $\cktStartSupply$ its burn color and
  $\cktCompleteRedeem$ its mint color \\
& $\wrapAddr$ & \srcWrapper & the interest-bearing wrapper; gives $\cktStartRedeem$ its burn color
  and $\cktCompleteSupply$ its mint color. The only spender $\cktApproveStata$ approves, and the
  only destination a supply or redeem request sends to \\
\end{longtable}
\endgroup

\noindent Also pinned at setup and public: the vault's own address $\vaultAddr$, the EVM chain
identifier, the Signet singleton, and the deployer's identity commitment


\vspace{5pt}
\newversion{The latest release replaces $\sigNonce$ and the five settle-view maps with a
request queue shared by every action: $\inputRequestBuffer$ and $\outputRequestBuffer$ hold
requests before and after nonce assignment, $\inputAttestationBuffer$ and
$\outputAttestationBuffer$ hold verified attestations before and after they are folded into
the block-height watermark $\lastseen$, $\vaultNonce$ is the vault EVM account's next nonce,
and \code{evictionMap} maps a sent $\rid$ back to its queue entry. Each action keeps a map of
its queued arguments (\code{depositArgsMap}, \code{withdrawArgsMap}, \ldots) and a map of its
sent requests (\code{bidirectionalDepositMap}, \ldots). A queue entry carries exactly one
user-dependent value, the ownership commitment
$\ownComm(\inIndex, \usersk) = \Hash(\usersk \concat \inIndex \concat \text{``vault:request:''})$
over a caller-chosen random index $\inIndex$, which stands in for $\refundComm(\usersk, \rid)$.
Everything else in these tables is bookkeeping (nonces, block heights, argument hashes and
attestation digests) that does not depend on the caller, so the tables do not alter the
privacy proof. They preserve the property that each request is unique and processed once:
completion consumes the entry, and a second identical request is refused while the first is
open.}
%-------------------------------------------------------------------------------------
%-------------------------------------------------------------------------------------
\subsection{The shielded Layer: Midnight Computation}
\label{ssec:shieldedcomputation}

Midnight provides shielded coin payments and private circuit evaluation.
We highlight with  \mnhl{shaded blue} the computation that are strictly Midnight's layer  and not the vault's own logic.

\vspace{5pt}
\noindent
\textbf{Midnight shielded-coin layer.}
In Midnight,  a raw secret coin is the tuple:
$c=(\mathsf{nonce}, \mathsf{color}, \mathsf{value})$.
Here are the functions applied to the secret coin:
\begin{align*}
\hypertarget{fn:coinCommit}{\coinCommit(c, \mathsf{rcpt})} &= \PersistantHash\bigl(\text{cc-tag} \concat \mathsf{nonce} \concat
  \mathsf{color} \concat \mathsf{value} \concat \mathsf{rcpt}\bigr)\\
\hypertarget{fn:coinNull}{\coinNull(c, \mathsf{sndr})} &= \PersistantHash\bigl(\text{cn-tag} \concat \mathsf{nonce} \concat
  \mathsf{color} \concat \mathsf{value} \concat \mathsf{sndr}\bigr)\\
\hypertarget{fn:tokenTypeOf}{\tokenTypeOf(\domSep(\tokType), \vaultAddr)} &= \PersistantHash\bigl(\text{tok-tag} \concat
  \domSep(\tokType) \concat \vaultAddr\bigr)
\end{align*}


% {\em User's coins vs contract coins.} For a \emph{user}-held coin the
% recipient  holds the coin \emph{public} key and the sender slot the coin \emph{secret} key, so
% an observer holding $\coinCommit$ cannot compute $\coinNull$. For a \emph{contract}-held coin both
% slots hold the same public contract address.


\medskip
\noindent\textbf{ Midnight Private-Circuit Computation.}
Midnight's contract computation is performed through circuits.
Every argument of a circuit is a \emph{private} input
of the proof.
A call to a  circuit $C$ on private arguments $\mathsf{in}$ and witness $\wit$ produces public outputs and state changes
for the vault contract $\Dvault$ and the Midnight chain $\Dshd$;
it also produces a statement $\stmt$ and a zk proof $\zkproof$.
\[
  \bigl(\; \Dvault,\;\; \Dshd,\;\; \stmt,\;\; \zkproof \;\bigr).
\]

\begin{center}
\begin{tabular}{@{}L{1.6cm}L{5cm}L{5cm}@{}}
\toprule
\textbf{Output} & \textbf{It contains} & \textbf{Chain Updates} \\
\midrule
$\Dvault$ & the list of assignments the call makes to the vault's public entries of
  \S\ref{ssec:state} & \textbf{the vault's} state on $\Lshd$ \\
\addlinespace[2pt]
$\Dshd$ & the list of coin commitments, nullifiers and mint records the call declares &
  \bluedot\textbf{Midnight's} state, not the vault's.\\
\addlinespace[2pt]
$\stmt$ & the public statement: the tuple an observer of $\Lshd$ reads off the call &
  --- \\
\addlinespace[2pt]
$\zkproof$ & a proof that $\Rel_C(\stmt, \wit)$ holds & --- \\
\bottomrule
\end{tabular}
\end{center}

 $\Dvault$, $\Dshd$ and $\stmt$  respect the following inclusion:

\[
  \Dvault \;\subseteq\; \stmt, \qquad \Dshd \;\subseteq\; \stmt,
  \qquad\text{and in general}\qquad \stmt \;\supsetneq\; \Dvault \cup \Dshd .
\]
\begin{itemize}[nosep,leftmargin=1.4em,topsep=3pt]
\item $\stmt$ also carries values that  do not update any state: the
      identifier $\rid$,  the commitment $\comIO$, the
      attestation the caller presented, and $\zkproof$ itself.
\item $\Rel_C$ is a \emph{predicate over $\stmt$ and $\wit$}, not a state update. This statement is verified inside the proof. Everything in $\stmt$ is visible.
\end{itemize}

\medskip
\noindent
\textbf{Midnight   Commitments.} We write $\comIO$ to denote the commitments  to the inputs and outputs of the circuit being called which are produced by the circuit computation.
The statement $\stmt$ contains $\comIO$.


\medskip
\noindent\textbf{How to read a circuit box.} \emph{Arguments} are the  inputs to the circuit (private);
\emph{witness} are the secret used to prove relations,   \emph{reads} are the public entries consulted; \emph{steps}
are the assertions and computations in source order;  A
\bluedot marked line or a blue block  highlights  Midnight's functions.


% \subsection{The attested outcome value}
% \label{ssec:outcomeval}

% Every settlement call consumes an attested outcome $\evmOut$, presented by the caller and
% authenticated only through $\attDigest(\rid, \evmOut)$. Exactly two kinds of value occur.

% \begin{definition}[$\succeeded$ and $\failed$]
% \label{def:evmout}
% Fix a request $\rid$ recorded by a request call. The attested outcome $\evmOut$ posted by $\SigSys$
% for $\rid$ is one of:
% \begin{itemize}[nosep,leftmargin=1.4em,topsep=2pt]
% \item a \emph{success value}, written $\evmOut \in \succeeded$: the encoding, under the response
%       schema $\schemaOut$ named in $\req$, of the return value of the EVM call that $\req$
%       requested. Its width and contents depend on the schema --- a one-byte boolean for a transfer,
%       an eight-byte integer for a swap, a supply or a redeem.
% \item the \emph{failure value}, written $\failed$: a single fixed protocol constant, independent of
%       the schema and of the request, posted when the requested transaction never executed. It is
%       the value $\failed$ that gates $\cktRefund$.
% \end{itemize}
% $\succeeded$ and $\failed$ are disjoint, and this disjointness is what routes an outcome to one
% settlement circuit rather than another.
% \end{definition}

%-------------------------------------------------------------------------------------
% ---------------------------------------------------------------------


\section{Our Implementation: Circuits}
\label{sec:instantiation}

In this section we show our implementation of the interface of Definition~\ref{def:system}.

\noindent
\textbf{Two calls per flow.}
Each procedure of Definition~\ref{def:system} is performed by the circuits below.
Each action: $\Deposit$, $\Withdraw$,  $\Swap$, $\Supply$, $\Redeem$, is performed by \emph{two} calls to the vault contract, with the $\SigSys$ bridge in
between. The first call records what the MPC should sign and takes
the values relevant to the call (transaction parameters). The second presents $\SigSys$'s attestation of what
happened on $\Lpub$ (Ethereum), and complete the request on $\Lshd$ (Midnight). Request calls are prefixed with \textsf{start}; settlement calls are prefixed with \textsf{complete}
The following table maps each call to their circuits.

\begin{center}
\begin{tabular}{@{}L{2.4cm}L{11.8cm}@{}}
\toprule
\textbf{Procedure} & \textbf{Circuits that implement it} \\
\midrule
$\Setup$    & $\cktInit$ (\cktref{ckt:initialize}); $\cktApprove$ (\cktref{ckt:approverouter}); $\cktApproveStata$ (\cktref{ckt:approvestata}) \\
$\Deposit$  & $\cktStartDeposit$ (\cktref{ckt:deposit}), then $\cktCompleteDeposit$ (\cktref{ckt:completedeposit}) \\
$\Withdraw$ & $\cktStartWithdraw$ (\cktref{ckt:startwithdraw}), then $\cktCompleteWithdraw$ (\cktref{ckt:completewithdraw}) or $\cktRefundWithdraw$ (\cktref{ckt:refundwithdraw}) \\
$\Swap$     & $\cktStartSwap$ (\cktref{ckt:startswap}), then $\cktCompleteSwap$ (\cktref{ckt:completeswap}) or $\cktRefundSwap$ (\cktref{ckt:refundswap}) \\
$\Supply$   & $\cktStartSupply$ (\cktref{ckt:startsupply}), then $\cktCompleteSupply$ (\cktref{ckt:completesupply}) or $\cktRefundSupply$ (\cktref{ckt:refundsupply}) \\
$\Redeem$   & $\cktStartRedeem$ (\cktref{ckt:startredeem}), then $\cktCompleteRedeem$ (\cktref{ckt:completeredeem}) or $\cktRefundRedeem$ (\cktref{ckt:refundredeem}) \\
\bottomrule
\end{tabular}
\end{center}


\newversion{In the latest release the global counter $\sigNonce$ is dropped. A request id
is made distinct by the EVM transaction nonce instead, and freshness comes from the
block-height watermark $\lastseen$. Transaction nonces are not assigned at request time by
the caller of {\code{startX}} (for any action X except Deposit, whose caller names the nonce
of their own derived account), but by a dedicated circuit that any party can call.

Specifically, each circuit {\code{startX}} only validates and records the transaction
parameters, without the EVM nonce and without the request id, under a caller-chosen random
index $\inIndex$. This \emph{partial} request is added to the buffer $\inputRequestBuffer$ of
requests waiting for nonce assignment. Every so often the queue is processed by the circuit
$\cktFlush$, which assigns the vault account's next nonce $\vaultNonce$ (a counter only
$\cktFlush$ reads or writes) to each vault-signed request and stamps every request with the
current $\lastseen$, and by the per-action circuits $\cktSend$ (\code{sendDeposit},
\code{sendWithdraw}, \ldots), which build the complete request from the flushed entry, compute
$\rid = \reqIdOf(\req)$ and notify the MPC. On the way back, the circuits $\cktQueueAtt$
verify the MPC's attestation signature and queue the attestation, $\cktFlush$ folds it into
$\outputAttestationBuffer$ and raises $\lastseen$ to its block height, and {\code{completeX}}
consumes the request together with its attestation.
$\cktFlush$, $\cktSend$ and $\cktQueueAtt$ are permissionless, take no witness and can be
executed by everyone.
Note, these circuits only add nonce and time-stamping information that is not user-dependent.
As such they do not alter the privacy analysis described in Sec.~\ref{sec:proof}. They
implement efficiency optimisations, contention avoidance and replay-attack protection. The
refund circuits ({\code{refundX}}) are folded into {\code{completeX}}, which refunds on a
failed or unviable verdict, and every {\code{completeX}}, the success branch of
$\cktCompleteWithdraw$ included, is gated on the requester's ownership commitment.
}


\medskip
\noindent
\textbf{Notification to $\SigSys$.} Every request call ends with a
cross-contract call to the Signet singleton, which emits an event the MPC polls. We write
$\notif(\vaultAddr, \rid, M)$ for that event: it carries the vault's address, the request
identifier, and the identity of the ledger map $M$ in which the request record was stored, so
that the MPC knows where to read $\req$ from.


\subsection{Setup: \texorpdfstring{$\Setup^{\Lpub,\Lshd}(1^{\secpar})$}{Setup}}
\label{ssec:setupinst2}

\begin{ProcBox}{$\pp \gets \Setup^{\Lpub,\Lshd}(1^{\secpar})$}
\sloppy Run once by the deployer, who holds a secret $\usersk_D$ of their own.
\begin{Steps}
\item \textbf{Deploy the contract} on $\Lshd$ with two constructor arguments: the deployer's
      identity $\deplComm = \lnk{fn:userComm}{\userComm(\usersk_D)}$ and the address of the Signet
      singleton. The contract's own address $\vaultAddr$ becomes known here, which is why the
      remaining values cannot be pinned in the constructor.
\item \textbf{Derive off-chain}
      \[\begin{aligned}
        \vaultEvmAddr &\gets \Derive(\mpcRootPk\bc \vaultAddr\bc \keyPathLit),\\
        \mpcRespKey &\gets \Derive(\mpcRootPk\bc \vaultAddr\bc \text{``midnight response key''}).
      \end{aligned}\]
\item \textbf{Run $\cktInit$} --- \cktref{ckt:initialize}.
\item \textbf{Run $\cktApprove$} once per ERC20 that will ever be swapped ---
      \cktref{ckt:approverouter}.
\end{Steps}
Output the public parameters:\\
$\pp = (\vaultAddr, \vaultEvmAddr, \mpcRespKey, \chainId, \routerAddr, \text{singleton address})$.
\end{ProcBox}


\begin{CktBox}{ckt:initialize}{$\cktInit$}
\fld{Arguments.} $\vaultEvmAddr$, $\routerAddr$, \ADDEDINL{$\underAddr$, $\wrapAddr$},
$\chainId$, $\mpcRespKey$.
\fld{Witness.} $\usersk$.
\fld{Reads.} $\initFlag$, $\deplComm$.
\fld{Steps.}
\begin{enumerate}[nosep,leftmargin=1.7em,topsep=2pt]
\item $\assert \initFlag = 0$.
\item $\assert \lnk{fn:userComm}{\userComm(\usersk)} = \deplComm$.
\item $\assert \chainId > 0$, \; $\routerAddr \neq 0$,
      \ADDEDINL{\; $\underAddr \neq 0$, \; $\wrapAddr \neq 0$}.
\item Write $\initFlag \gets 1$ and the six pinned values.
\end{enumerate}
\begin{OutList}
\oitem{\Dvault} $\initFlag \gets 1$; $\vaultEvmAddr$, $\routerAddr$,
  \ADDEDINL{$\underAddr$, $\wrapAddr$}, $\chainId$, $\mpcRespKey$ written.
\oitem{\Dshd} $\emptyset$.
\oitem{\stmt} $(\cktInit,\; \Dvault,\; \comIO)$.
\oitem{\wit} $\usersk$.
\oitem{\Rel} $\userComm(\usersk) =\deplComm$.
\end{OutList}
\end{CktBox}


\begin{CktBox}{ckt:approverouter}{$\cktApprove$}
\fld{Arguments.} $\tokType$, $\evmNonce$, $\keyVer$.
\fld{Witness.} none.
\fld{Reads.} $\initFlag$, $\routerAddr$, $\vaultAddr$, $\sigNonce$, $\chainId$.
\fld{Steps.}
\begin{enumerate}[nosep,leftmargin=1.7em,topsep=2pt]
\item $\assert \initFlag \geq 1$ and $\tokType \neq 0$.
\item $\evmTxOf \gets$ ``approve $\routerAddr$ for an unlimited amount of $\tokType$''. \emph{The
      spender and the amount come from pinned state; the caller chooses only the token.}
\item $\req \gets (\vaultAddr\bc \sigNonce\bc \keyVer\bc \keyPath{=}\keyPathLit\bc \evmTxOf\bc
      \chainId)$; \; $\rid \gets \lnk{fn:reqIdOf}{\reqIdOf(\req)}$; \;
      $\assert \rid \notin \reqMap$.
\end{enumerate}
\begin{OutList}
\oitem{\Dvault} $\sigNonce \mathrel{+}= 1$; \; $\reqMap[\rid] \gets \req$ ---
  \emph{never removed by any circuit}.
\oitem{\Dshd} $\emptyset$.
\oitem{\stmt} $(\cktApprove,\; \req,\; \rid,\; \notif(\vaultAddr, \rid, \reqMap))$.
\oitem{\wit} $\bot$.
\oitem{\Rel} $\mathsf{true}$ - the witness is empty and every argument is in $\stmt$, so
  $\zkproof$ certifies a computation an observer can redo.
\end{OutList}
\end{CktBox}

% ---------------------------------------------------------------------
\subsection{Deposit}
\label{ssec:depinst2}

\begin{ProcBox}{$\Pi.\Deposit^{\Lpub, \Lshd}\bigl([\pubAddr]\bc \shdAcct\bc \tokType\bc \amt\bc
      \auxx \,;\, (\secret{\usersk}, \freshrn{\mintNonce})\bigr)$.}
\sloppy
\emph{Private input.} secret key $\usersk$; and a fresh $\mintNonce$, used at Step~4.\\
\emph{Public arguments.} $\tokType$ is an ERC20 contract,
and $\auxx = (\evmNonce, \gasEnv, \keyVer)$.
The receiving \zkswap{} account $\shdAcct$  is not public per-se, but we
do not aim to protect it.

\begin{Steps}
\item \textbf{Fund the derived account:} outside any circuit, on $\Lpub$. The caller pays $\amt$ of $\tokType$  to  the MPC controlled account $\derivedAddr$ on $\Lpub$. Where,
      $\derivedAddr = \Derive\bigl(\mpcRootPk\bc \vaultAddr\bc
      \lnk{fn:userComm}{\userComm(\usersk)}\bigr)$. See \S\ref{ssec:derived}.
\item \textbf{Request call: $\cktDeposit$} (\cktref{ckt:deposit}). Recomputes
      $\keyPath \gets \lnk{fn:userComm}\userComm(\usersk)$ from the witness and request
      $\text{transfer}(\vaultEvmAddr, \amt)$ under $\rid$. \emph{$\Dshd = \emptyset$: no coin
      moves.}
\item \textbf{MPC}. The committee signs on behalf of  $\derivedAddr$ a transfer to $\vaultEvmAddr$, and posts $\attSig$ over
      $\lnk{fn:attDigest}{\attDigest(\rid, \evmOut)}$.
\item \textbf{Settlement call: $\cktCompleteDeposit$} (\cktref{ckt:completedeposit}). Verifies
      $\attSig$, consumes $\reqMap[\rid]$, authenticated the request on $\userComm(\usersk) = \req.\keyPath$, and mints
      $\amt$ to $\shdAcct$ using randomness $\mintNonce$.
\end{Steps}

\medskip
Outcome: $\succeeded$ once Step~4 completes; $\failed$ if the EVM transfer did not execute, in which
case the value never left $\derivedAddr$;
Private output: $\usersk,\mintNonce$.
\end{ProcBox}


\begin{CktBox}{ckt:deposit}{$\srcStartDeposit$}
\fld{Arguments.} $\evmNonce$, $\gasEnv$, $\keyVer$, $(\tokType, \amt)$.
\fld{Witness.} $\usersk$.
\fld{Reads.} $\initFlag$, $\vaultEvmAddr$, $\vaultAddr$, $\sigNonce$, $\chainId$.
\fld{Steps.}
\begin{enumerate}[nosep,leftmargin=1.7em,topsep=2pt]
\item $\assert \initFlag \geq 1$, \; $\tokType \neq 0$, \; $\gasEnv > 0$, \;
      $0 < \amt \leq 2^{64}-1$.
\item $\keyPath \gets \lnk{fn:userComm}{\userComm(\usersk)}$.
\item $\evmTxOf \gets$ ``$\text{transfer}(\vaultEvmAddr, \amt)$ on $\tokType$'', with $\evmNonce$,
      $\gasEnv$, on $\chainId$.
\item $\req \gets (\vaultAddr\bc \sigNonce\bc \keyVer\bc \keyPath\bc \evmTxOf\bc \chainId)$; \;
      $\rid \gets \lnk{fn:reqIdOf}{\reqIdOf(\req)}$; \;
      \ADDEDINL{$\assert \rid \notin \depositMap$}.
\end{enumerate}
\begin{OutList}
\oitem{\Dvault} $\sigNonce \mathrel{+}= 1$; \; \ADDEDINL{$\depositMap[\rid] \gets \req$}; \;
  \ADDEDINL{$\depositView[\rid] \gets (\keyPath,\; \tokType,\; \amt)$}.
\oitem{\Dshd} $\emptyset$.
\oitem{\stmt} $(\cktStartDeposit,\; \req,\; \rid,\; \depositView[\rid],\;
  \notif(\vaultAddr, \rid, \depositMap),\; \comIO)$, where $\req$ exposes
  $\tokType, \amt, \vaultEvmAddr, \evmNonce, \gasEnv, \sigNonce$ and
  $\keyPath = \userComm(\usersk)$.
\oitem{\wit} $\usersk$.
\oitem{\Rel} $\req.\keyPath = \userComm(\usersk)$ \;$\wedge$\;
  $\depositView[\rid].\fComm = \userComm(\usersk)$.
\end{OutList}
\end{CktBox}


\begin{CktBox}{ckt:completedeposit}{$\cktCompleteDeposit$}
\fld{Arguments.} $\rid$, $\attSig$, $\evmOut$, $\freshrn{\mintNonce}$,  $\shdAcct$.
\fld{Witness.} $\usersk$.
\fld{Reads.} $\initFlag$, $\mpcRespKey$, \ADDEDINL{$\depositMap$, $\depositView[\rid]$}.
\fld{Steps.}
\begin{enumerate}[nosep,leftmargin=1.7em,topsep=2pt]
\item $\assert \initFlag \geq 1$ and $\evmOut$ reports success.
\item $\assert \EcdsaVf\bigl(\lnk{fn:attDigest}{\attDigest(\rid, \evmOut)}\bc \attSig\bc
      \mpcRespKey\bigr)$ --- \textbf{gate 1}, the authentication of the MPC.
\item \ADDEDINL{$\assert \rid \in \depositMap$; \; remove $\depositMap[\rid]$} ---
      \textbf{gate 3}, at most once, enforced by the ledger through $\Dvault$, not by $\Rel$.
\item \ADDEDINL{$\assert \lnk{fn:userComm}{\userComm(\usersk)} = \depositView[\rid].\fComm$;} \;
      remove $\depositView[\rid]$ --- \textbf{gate 2}, enforces depositor-only.
\item \ADDEDINL{$\amt \gets \depositView[\rid].\fAmt$, \;
      $\tokType \gets \depositView[\rid].\fErc$}.
\end{enumerate}
\MidNote{\textbf{Midnight operation.} Execute Midnight Circuit $\mintShieldedToken$ on input
$\domSep, \amt, \mintNonce, \callerZKey$. Let $\coinMint =
\bigl(\;\freshrn{\mintNonce}, \mathsf{color}, \amt \bigr)$}
where  $\mathsf{color}=\lnk{fn:tokenTypeOf}{\tokenTypeOf}\bigl(\domSep(\tokType),
\vaultAddr\bigr)$
\begin{OutList}
\oitem{\Dvault} \ADDEDINL{$\depositMap[\rid]$ removed; $\depositView[\rid]$ removed}.
\oitem{\Dshd} \bluedot mint record $\bigl(\lnk{fn:domSep}{\domSep(\tokType)},\, \amt\bigr)$;
  \bluedot output  $\lnk{fn:coinCommit}{\coinCommit(\coinMint, \shdAcct)}$.
\oitem{\stmt} $\bigl(\cktCompleteDeposit,\; \rid,\; \attSig,\;
  (\domSep(\tokType), \amt),\; \coinCommit(\coinMint, \callerZKey)\bigr)$.
\oitem{\wit} $(\usersk,\; \mintNonce,\; \shdAcct)$.
\oitem{\Rel} $\userComm(\usersk)$ $= \depositView[\rid].\fComm$ \;$\wedge$\;
  $\EcdsaVf(\attDigest(\rid,\evmOut)$, $\attSig$, $\mpcRespKey)$ \;$\wedge$\;
  the published commitment opens to $(\mintNonce, \cdot, \amt, \callerZKey)$.
\end{OutList}
\end{CktBox}


\begin{PubBox}{Transcript from  $\Deposit$}
\textbf{On $\Lpub$:} [$\pubAddr \to \derivedAddr$],  and
$\derivedAddr \to \vaultEvmAddr$ (signed by the committee), both in the clear.

\smallskip
\textbf{On $\Lshd$:} $\stmt_{\cktDeposit}$ and
$\stmt_{\cktCompleteDeposit}$.

\smallskip
\textbf{An observer therefore holds} $\userComm(\usersk)$, hence $\derivedAddr$, hence the edge from
the caller's own EVM account to their identity in this system; $\tokType$ and $\amt$; $\rid$, so \textbf{the settlement call is  linked
to the request call}. Hidden: $\usersk$, $\mintNonce$.
\end{PubBox}
% ---------------------------------------------------------------------
\subsection{Withdraw}
\label{ssec:wdinst2}

\begin{ProcBox}{$\Withdraw^{\Lpub, \Lshd}\bigl(\shdAcct\bc \destAddr*\bc \tokType\bc \amt\bc
      \auxx \,;\, \privIn\bigr)$}\label{box:withdraw}
\sloppy
\emph{Private input.} $\privIn$ holds $\usersk$; the secret key $\coinsk$ associated to the $\zkswap$ account $\srcAcct$; a
fresh $\burnNonce$, used at Step~1; and, only if the flow ends in $\failed$, a fresh $\mintNonce$
used at Step~4.
\emph{Public arguments.} $\auxx = (\evmNonce, \keyVer)$ only: the gas envelope is fixed in the
contract, because the \emph{vault} pays for this transaction.

 \begin{Steps}
 \item \textbf{Request call: $\cktWithdraw$} (\cktref{ckt:startwithdraw}).
 The caller hands the contract  a coin descriptor   $\coinIn = \bigl(\burnNonce,\;\mathsf{color}, \amt\bigr)$; the circuit burns it, records  $\text{transfer}(\destAddr^*, \amt)$ under $\rid$ with $\keyPath = \keyPathLit$ o, and add the authentication tag to the map:  $\wdView[\rid] \gets \lnk{fn:refundComm}{\refundComm(\usersk, \rid)}$.
   \smallskip

 \item \textbf{MPC}. The MPC signs the evm request with the key derived
       at $\keyPath = \keyPathLit$, so from $\vaultEvmAddr$, the pool's own account. Once the request is executed,  the MPC signs an attestation $\attSig$ of the outcome $\evmOut$.
\item \textbf{Settlement call on success: $\cktCompleteWithdraw$}, success branch  (\cktref{ckt:completewithdraw}). \textbf{$\Dshd = \emptyset$ and $\wit = \bot$ on this branch}.
\item \textbf{Settlement call on failure: $\cktCompleteWithdraw$'s failure branch, or $\cktRefundWithdraw$}
      -  \cktref{ckt:completewithdraw} and \cktref{ckt:refundwithdraw}, both which gates on
      $\lnk{fn:refundComm}{\refundComm(\usersk, \rid)}$ against the stored marker  and re-mints $\amt$ under a fresh
      $\mintNonce$. The two circuits differ only in which attested output routes to them: an
      executed transfer that returned false goes to $\cktCompleteWithdraw$; a transaction that never
       executed goes to $\cktRefundWithdraw$.
 \end{Steps}

\medskip
Outcome: $\succeeded$ via Step~3, $\failed$ via Step~4. Private output: $\usersk$.
\end{ProcBox}


\begin{CktBox}{ckt:startwithdraw}{$\cktStartWithdraw$}
\fld{Arguments.} $\evmNonce$, $\keyVer$, $(\tokType, \amt, \destAddr)$, and
$\coinIn = (\burnNonce, \mathsf{color}, \mathsf{value})$.
\fld{Witness.} $\usersk$ (and $\coinsk$ at the wallet layer).
\fld{Reads.} $\initFlag$, $\vaultAddr$, $\sigNonce$, $\chainId$.
\fld{Steps.}
\begin{enumerate}[nosep,leftmargin=1.7em,topsep=2pt]
\item $\assert \initFlag \geq 1$, \; $\tokType \neq 0$, \; $0 < \amt \leq 2^{64}-1$.
\item $\assert \coinIn.\mathsf{color} =
      \lnk{fn:tokenTypeOf}{\tokenTypeOf(\domSep(\tokType), \vaultAddr)}$ \; and \;
      $\coinIn.\mathsf{value} = \amt$.
\item $\evmTxOf \gets$ ``$\text{transfer}(\destAddr, \amt)$ on $\tokType$'', with $\evmNonce$ and a
      contract-fixed gas envelope, on $\chainId$.
\item $\keyPath \gets \keyPathLit$.
\item $\req \gets (\vaultAddr\bc \sigNonce\bc \keyVer\bc \keyPath\bc \evmTxOf\bc \chainId)$; \;
      $\rid \gets \lnk{fn:reqIdOf}{\reqIdOf(\req)}$; \; $\assert \rid \notin \reqMap$.
\end{enumerate}
\MidNote{\textbf{Midnight operations: burn.}
$\mnhl{\code{receiveShielded}}(\coinIn)$ then
$\mnhl{\code{sendImmediateShielded}}(\coinIn, \burnAddr, \amt)$.}
\begin{OutList}
\oitem{\Dvault} $\sigNonce \mathrel{+}= 1$; \; $\reqMap[\rid] \gets \req$; \;
  \ADDEDINL{$\wdView[\rid] \gets \bigl(\lnk{fn:refundComm}{\refundComm(\usersk, \rid)},\;
  \tokType,\; \amt\bigr)$}.
\oitem{\Dshd} \bluedot $\coinCommit(\coinIn, \vaultAddr)$, \;
  \bluedot $\coinNull(\coinIn, \vaultAddr)$, \;
  \bluedot $\coinCommit(\coinBurnOut, \burnAddr)$, \;
  \bluedot $\{\coinNull(\coinFund_i, \coinsk)\}_i$ and $\coinCommit(\coinChange, \srcAcct)$.
\oitem{\stmt} $\bigl(\cktStartWithdraw,\; \req,\; \rid,\; \wdView[\rid],\; \Dshd,\;
  \notif(\vaultAddr, \rid, \reqMap),\; \comIO\bigr)$.
  \ADDEDINL{The view publishes $\tokType$ and $\amt$.}
\oitem{\wit} $(\usersk,\; \burnNonce,\; \coinsk)$.
\oitem{\Rel} $\coinIn.\mathsf{color} = \tokenTypeOf(\domSep(\tokType), \vaultAddr)$ \;$\wedge$\;
  $\coinIn.\mathsf{value} = \amt$ \;$\wedge$\;
  $\wdView[\rid].\fComm = \refundComm(\usersk, \rid)$.
\end{OutList}
\end{CktBox}


\begin{CktBox}{ckt:failgate}{$\cktFailGate$ - shared across refund circuits (\srcFailGate)}
Not exported. Called first by $\cktRefundWithdraw$, $\cktRefundSwap$, $\cktRefundSupply$ and
$\cktRefundRedeem$.
\fld{Arguments.} $\rid$, $\attSig$, $\evmOut$ of width $5$.
\fld{Reads.} $\initFlag$, $\mpcRespKey$.
\fld{Steps.}
\begin{enumerate}[nosep,leftmargin=1.7em,topsep=2pt]
\item $\assert \initFlag \geq 1$.
\item $\assert \EcdsaVf\bigl(\lnk{fn:attDigest}{\attDigest(\rid, \evmOut)}\bc \attSig\bc
      \mpcRespKey\bigr)$.
\item $\assert \evmOut = \failed$.
\end{enumerate}
\begin{OutList}
\oitem{\Dvault} none of its own.
\oitem{\Dshd} $\emptyset$.
\oitem{\stmt} $(\rid,\; \attSig,\; \evmOut)$, merged into the caller circuit's statement.
\oitem{\wit} $\bot$.
\oitem{\Rel} $\EcdsaVf(\attDigest(\rid,\evmOut), \attSig, \mpcRespKey)$ \;$\wedge$\;
  $\evmOut = \failed$.
\end{OutList}
\end{CktBox}


\begin{CktBox}{ckt:completewithdraw}{$\cktCompleteWithdraw$}
\fld{Arguments.} $\rid$, $\attSig$, $\evmOut$ of width $1$, $\mintNonce$.
\fld{Witness.} $\usersk$ - \emph{failure branch only}.
\fld{Reads.} $\initFlag$, $\mpcRespKey$, \ADDEDINL{$\wdView[\rid]$}, $\reqMap$.
\fld{Steps.}
\begin{enumerate}[nosep,leftmargin=1.7em,topsep=2pt]
\item $\assert \initFlag \geq 1$ and
      $\assert \EcdsaVf\bigl(\lnk{fn:attDigest}{\attDigest(\rid, \evmOut)}\bc \attSig\bc
      \mpcRespKey\bigr)$.
\item \ADDEDINL{$\assert \rid \in \wdView$}.
\item Remove $\reqMap[\rid]$. \ADDEDINL{The record is not read.}
\item $\okbit \gets$ the success bit of $\evmOut$.
\item \ADDEDINL{If $\okbit = 0$:
      $\assert \lnk{fn:refundComm}{\refundComm(\usersk, \rid)}$
      $= \wdView[\rid].\fComm$, then mint
      $\coinMint = \bigl(\mintNonce,\, \tokenTypeOf(\domSep(\wdView[\rid].\fErc), \vaultAddr),\,
      \wdView[\rid].\fAmt\bigr)$
      to the caller's own shielded account.}
\item Remove $\wdView[\rid]$.
\end{enumerate}
\begin{OutList}
\oitem{\Dvault} $\reqMap[\rid]$ removed; \ADDEDINL{$\wdView[\rid]$ removed}.
\oitem{\Dshd} $\okbit = 1$: $\emptyset$. \quad $\okbit = 0$: \bluedot mint record
  $\bigl(\domSep(\wdView[\rid].\fErc),\, \wdView[\rid].\fAmt\bigr)$; \;
  \bluedot $\coinCommit(\coinMint, \text{caller})$.
\oitem{\stmt} $\bigl(\cktCompleteWithdraw,\; \rid,\; \attSig,\; \okbit,\; \Dshd\bigr)$.
\oitem{\wit} $\okbit = 1$: $\bot$. \quad $\okbit = 0$: $(\usersk, \mintNonce)$.
\oitem{\Rel} $\okbit = 1$: $\EcdsaVf(\attDigest(\rid,\evmOut), \attSig, \mpcRespKey)$. \quad
  $\okbit = 0$: the same, $\wedge$ $\refundComm(\usersk, \rid)$ $= \wdView[\rid].\fComm$ $\wedge$
  the published commitment opens to $(\mintNonce, \cdot, \wdView[\rid].\fAmt, \text{caller})$.
\end{OutList}
\fld{$\wit = \bot$ on the success branch.} The proof certifies a statement with no secret in it, so
any third party holding the attestation can produce the same call. $\mintNonce$ is an argument on
both branches; on the success branch it is simply unused.
\end{CktBox}


\begin{CktBox}{ckt:refundwithdraw}{ $\cktRefundWithdraw$}
\fld{Arguments.} $\rid$, $\attSig$, $\evmOut$ of width $5$, $\mintNonce$.
\fld{Witness.} $\usersk$.
\fld{Reads.} $\initFlag$, $\mpcRespKey$, $\wdView[\rid]$, $\reqMap$.
\fld{Steps.}
\begin{enumerate}[nosep,leftmargin=1.7em,topsep=2pt]
\item Run $\cktFailGate(\rid, \attSig, \evmOut)$.
\item $\assert \rid \in \wdView$.
\item $\assert \lnk{fn:refundComm}{\refundComm(\usersk, \rid)} = \wdView[\rid].\fComm$ (authenticating the withdrawer).
\item Remove $\reqMap[\rid]$ and $\wdView[\rid]$.
\end{enumerate}
\MidNote{\textbf{Midnight operation.} Mint, to the caller's own shielded account,
\[\coinMint = \bigl(\mintNonce,\;
\lnk{fn:tokenTypeOf}{\tokenTypeOf(\domSep(\wdView[\rid].\fErc), \vaultAddr)},\;
\wdView[\rid].\fAmt\bigr).\]}
\begin{OutList}
\oitem{\Dvault} $\reqMap[\rid]$ removed; $\wdView[\rid]$ removed.
\oitem{\Dshd} \bluedot mint record
  $\bigl(\domSep(\wdView[\rid].\fErc),\, \wdView[\rid].\fAmt\bigr)$; \;
  \bluedot $\lnk{fn:coinCommit}{\coinCommit(\coinMint, \text{caller})}$.
\oitem{\stmt} $\bigl(\cktRefundWithdraw,\; \rid,\; \attSig,\; \Dvault,\; \Dshd\bigr)$.
\oitem{\wit} $(\usersk,\; \mintNonce)$.
\oitem{\Rel} $\refundComm(\usersk, \rid)$ $= \wdView[\rid].\fComm$ \;$\wedge$\; the relation of
  $\cktFailGate$ \;$\wedge$\; the published commitment opens to
  $(\mintNonce, \cdot, \wdView[\rid].\fAmt, \text{caller})$.
\end{OutList}
\end{CktBox}


\begin{PubBox}{What $\Withdraw$ publishes}
\textbf{On $\Lpub$:} on success, one transfer $\vaultEvmAddr \to \destAddr$ of $\amt$ of $\tokType$.
\emph{Nothing tied to $\srcAcct$ appears.}

\smallskip
\textbf{On $\Lshd$:} $\stmt_{\cktWithdraw}$, the settlement
statement. Of these, the components that involve the caller at all are: the marker
$\refundComm(\usersk, \rid)$;  the coins sent to the burn address $\burnAddr$ in the $\cktStartWithdraw$ (the identity of the payer for this coins is protected by Mindnight's library under the assumption that $\burnNonce$ is private and is random), and -- in case of failure -- the new coins minted to the caller of $\cktRefundWithdraw$ (the identity of the payer for this coins is protected by Mindnight's library under the assumption that $\mintNonce$ is private and is random).


\smallskip
\textbf{Hidden:} which shielded account funded the call/ asked for a refund.
\end{PubBox}


%---------------------------------------------------------------------
%---------------------------------------------------------------------------


\subsection{Swap}
\label{ssec:swapinst2}

\begin{ProcBox}{$\Pi.\Swap^{\Lpub, \Lshd}\bigl(\shdAcct\bc \tokType\bc \amt\bc
    \tokTypeB\bc \amtB\bc \auxx \,;\, \privIn\bigr)$}
\sloppy
\emph{Private input.} As $\Withdraw$: $\usersk$, $\coinsk$ of $\srcAcct$, a fresh $\burnNonce$ at
Step~1, and a fresh $\mintNonce$ at Step~3 and a fresh $\changeNonce$ in Step 4.).
\emph{Public arguments.} $\amt$ is a \emph{maximum} input $\amtMax$ and $\amtB$ an exact output
$\amtOut$; (in an honest use $\srcAcct$ is the recipient of the swapped coin);
$\auxx = (\evmNonce, \keyVer, \feeTier)$.
\begin{Steps}
\item \textbf{Request call: $\cktSwap$} - \cktref{ckt:startswap}. Same shape as $\cktWithdraw$: a   coin descriptor for $\amtMax$ of $\tokType$ is surrendered and burned, an exact-output exchange is
      recorded under $\rid$ in the map $\swapMap$, and
      $\swapView[\rid] \gets \lnk{fn:refundComm}{\refundComm(\usersk, \rid)}$ is recorded.
\item \textbf{MPC}. As in withdraw,   the MPC signs the swap request on behalf of  $\vaultEvmAddr$. Once the request is executed,  the MPC signs an attestation $\attSig$ of the outcome $\evmOut$.
\item \textbf{Settlement call on success: $\cktCompleteSwap$} - \cktref{ckt:completeswap}.
If the transaction is successful, the user can claim the new-colored coin and the change, after proving knowledge of the pre-image of $\refundComm$.

\item \textbf{Settlement call on failure: $\cktRefundSwap$}. If the transaction reverts on $\Lpub$, user calls  Circuit \cktref{ckt:refundswap} to re-mint a coin with the full $\amtMax$ of the original color, after proving  knowledge of the pre-image of $\refundComm$.
\end{Steps}

\medskip
Outcome: $\succeeded$ via Step~3, $\failed$ via Step~4.
Private output: the new coins minted.
\end{ProcBox}


\begin{CktBox}{ckt:startswap}{$\cktStartSwap$}
\fld{Arguments.} $\evmNonce$, $\keyVer$, $(\tokType, \tokTypeB, \feeTier, \amtOut, \amtMax)$, and
$\coinIn = (\burnNonce, \mathsf{color}, \mathsf{value})$.
\fld{Witness.} $\usersk$ (and $\coinsk$ at the wallet layer).
\fld{Reads.} $\initFlag$, $\routerAddr$, $\vaultEvmAddr$, $\vaultAddr$, $\sigNonce$, $\chainId$.
\fld{Steps.}
\begin{enumerate}[nosep,leftmargin=1.7em,topsep=2pt]
\item $\assert \initFlag \geq 1$, \; $\tokType \neq 0$, \; $\tokTypeB \neq 0$, \;
      $0 < \amtOut \leq 2^{64}-1$, \; $0 < \amtMax \leq 2^{64}-1$.
\item $\assert \coinIn.\mathsf{color} =
      \lnk{fn:tokenTypeOf}{\tokenTypeOf(\domSep(\tokType), \vaultAddr)}$ \; and \;
      $\coinIn.\mathsf{value} = \amtMax$.
\item $\evmTxOf \gets$ ``exact-output exchange of at most $\amtMax$ of $\tokType$ for exactly
      $\amtOut$ of $\tokTypeB$ at $\feeTier$, recipient  is $\vaultEvmAddr$, no price bound'',
      to $\routerAddr$, contract-fixed gas envelope.
\item $\keyPath \gets \keyPathLit$, as in $\cktStartWithdraw$.
\item $\req \gets (\vaultAddr\bc \sigNonce\bc \keyVer\bc \keyPath\bc \evmTxOf\bc \chainId)$; \;
      $\rid \gets \lnk{fn:reqIdOf}{\reqIdOf(\req)}$; \; $\assert \rid \notin \swapMap$.
\end{enumerate}
\MidNote{\textbf{Midnight operations: burn}, exactly as in \cktref{ckt:startwithdraw}:
$\mnhl{\code{receiveShielded}}(\coinIn)$ then
$\mnhl{\code{sendImmediateShielded}}(\coinIn, \burnAddr, \amtMax)$.}
\begin{OutList}
\oitem{\Dvault} $\sigNonce \mathrel{+}= 1$; \; $\swapMap[\rid] \gets \req$; \;
  \ADDEDINL{$\swapView[\rid] \gets \bigl(\lnk{fn:refundComm}{\refundComm(\usersk, \rid)},\;
  \tokType,\; \tokTypeB,\; \amtOut,\; \amtMax\bigr)$}.
\oitem{\Dshd} as \cktref{ckt:startwithdraw}, with $\amtMax$ in place of $\amt$: \;
  \bluedot $\coinCommit(\coinIn, \vaultAddr)$, \; \bluedot $\coinNull(\coinIn, \vaultAddr)$, \;
  \bluedot $\coinCommit(\coinBurnOut, \burnAddr)$, \;
  \bluedot $\coinNull(\coinFund, \coinsk)$ and $\coinCommit(\coinChange, \srcAcct)$.
\oitem{\stmt} $\bigl(\cktStartSwap,\; \req,\; \rid,\; \swapView[\rid],\; \Dshd,\;
  \notif(\vaultAddr, \rid, \swapMap),\; \comIO\bigr)$, where $\req$ exposes
  $\tokType, \tokTypeB, \feeTier, \amtOut, \amtMax$.
  \ADDEDINL{The view publishes $\tokType$, $\tokTypeB$, $\amtOut$ and $\amtMax$.}
\oitem{\wit} $(\usersk,\; \burnNonce,\; \coinsk)$.
\oitem{\Rel} as \cktref{ckt:startwithdraw}, over $\swapView$ and $\amtMax$.
\end{OutList}
\end{CktBox}


\begin{CktBox}{ckt:completeswap}{$\cktCompleteSwap$}
\fld{Arguments.} $\rid$, $\attSig$, $\evmOut$ of width $8$, $\freshrn{\mintNonce}$,
\ADDEDINL{$\freshrn{\changeNonce}$}.
\fld{Witness.} $\usersk$, $\mintNonce$, \ADDEDINL{$\changeNonce$}.
\fld{Reads.} $\initFlag$, $\mpcRespKey$, \ADDEDINL{$\swapView[\rid]$}, $\swapMap$.
\fld{Steps.}
\begin{enumerate}[nosep,leftmargin=1.7em,topsep=2pt]
\item $\assert \initFlag \geq 1$; \; \ADDEDINL{$\assert \changeNonce \neq \mintNonce$}.
\item $\assert \EcdsaVf\bigl(\lnk{fn:attDigest}{\attDigest(\rid, \evmOut)}\bc \attSig\bc
      \mpcRespKey\bigr)$.
\item \ADDEDINL{$\assert \rid \in \swapView$}; \; remove $\swapMap[\rid]$.
\item \ADDEDINL{$\assert \lnk{fn:refundComm}{\refundComm(\usersk, \rid)} = \swapView[\rid].\fComm$}
      (authenticating the swapper); \; remove $\swapView[\rid]$.
\item \ADDEDINL{$\amtOut \gets \swapView[\rid].\fAmtOut$,}
      \ADDEDINL{$\tokTypeB \gets \swapView[\rid].\fTokOut$,}
      \ADDEDINL{$\amtMax \gets \swapView[\rid].\fAmtMax$,}
      \ADDEDINL{$\tokType \gets \swapView[\rid].\fTokIn$}; \;
      $\amtIn \gets$ the attested value in $\evmOut$. \;
      $\amtRem \gets \amtMax - \amtIn$.
\end{enumerate}
\MidNote{\textbf{Midnight operations: two mints.} $\mintShieldedToken$ is called twice, both times
to the caller's own shielded account:
$\coinMint = \bigl(\mintNonce,\; \tokenTypeOf(\domSep(\tokTypeB), \vaultAddr),\; \amtOut\bigr)$ and
$\coinRem = \bigl(\changeNonce,\; \tokenTypeOf(\domSep(\tokType), \vaultAddr),\; \amtRem\bigr)$.
The second runs even when $\amtRem = 0$.}
\begin{OutList}
\oitem{\Dvault} $\swapMap[\rid]$ removed; \ADDEDINL{$\swapView[\rid]$ removed}.
\oitem{\Dshd} \bluedot mint record $\bigl(\domSep(\tokTypeB),\, \amtOut\bigr)$; \;
  \bluedot mint record $\bigl(\domSep(\tokType),\, \amtRem\bigr)$; \;
  \bluedot $\coinCommit(\coinMint, \text{caller})$; \;
  \bluedot $\coinCommit(\coinRem, \text{caller})$.
\oitem{\stmt} $\bigl(\cktCompleteSwap,\; \rid,\; \attSig,\; \amtIn,\; \Dvault,\; \Dshd\bigr)$.
\oitem{\wit} $(\usersk,\; \mintNonce,\; \ADDEDINL{\changeNonce})$.
\oitem{\Rel} $\refundComm(\usersk, \rid)$ $= \swapView[\rid].\fComm$ \;$\wedge$\;
  $\EcdsaVf(\attDigest(\rid,\evmOut)$, $\attSig$, $\mpcRespKey)$ \;$\wedge$\;
  $\changeNonce$ $\neq \mintNonce$ \;$\wedge$\;
  the first commitment opens to $(\mintNonce, \cdot, \amtOut, \text{caller})$ \;$\wedge$\;
  the second opens to $(\changeNonce, \cdot, \amtRem, \text{caller})$.
\end{OutList}
\end{CktBox}

\begin{CktBox}{ckt:refundswap}{$\cktRefundSwap$ }
\fld{Arguments.} $\rid$, $\attSig$, $\evmOut$ of with 5, $\mintNonce$.
\fld{Witness.} $\usersk$.
\fld{Reads.} $\initFlag$, $\mpcRespKey$, $\swapView[\rid]$, $\swapMap$.
\fld{Steps.}
\begin{enumerate}[nosep,leftmargin=1.7em,topsep=2pt]
\item Run \hyperref[ckt:failgate]{$\cktFailGate(\rid, \attSig, \evmOut)$}.
\item $\assert \rid \in \swapView$.
\item $\assert \lnk{fn:refundComm}{\refundComm(\usersk, \rid)} = \swapView[\rid].\fComm$ -
      \textbf{the swapper gate}.
\item Remove $\swapMap[\rid]$ and $\swapView[\rid]$.
\end{enumerate}
\MidNote{\textbf{Midnight operation.} Mint, to the caller's own shielded account,
$\coinMint = \bigl(\mintNonce,\; \mathsf{color},\;
\swapView[\rid].\fAmtMax\bigr).$}
where $\mathsf{color}= \lnk{fn:tokenTypeOf}{\tokenTypeOf(\domSep(\swapView[\rid].\fTokIn),
\vaultAddr)}$.
\begin{OutList}
\oitem{\Dvault} $\swapMap[\rid]$ removed; $\swapView[\rid]$ removed.
\oitem{\Dshd} \bluedot mint record
  $\bigl(\domSep(\swapView[\rid].\fTokIn),\, \swapView[\rid].\fAmtMax\bigr)$; \;
  \bluedot $\coinCommit(\coinMint, \text{caller})$.
\oitem{\stmt} $\bigl(\cktRefundSwap,\; \rid,\; \attSig,\; \Dvault,\; \Dshd\bigr)$.
\oitem{\wit} $(\usersk,\; \mintNonce)$.
\oitem{\Rel} $\refundComm(\usersk, \rid)$ $= \swapView[\rid].\fComm$ \;$\wedge$\; the relation of
  $\cktFailGate$ \;$\wedge$\; the published commitment opens to
  $(\mintNonce, \cdot, \swapView[\rid].\fAmtMax, \text{caller})$.
\end{OutList}
\end{CktBox}


\begin{PubBox}{What $\Swap$ publishes}
\textbf{On $\Lpub$:} the exchange in the clear: both tokens/colors and amounts, $\feeTier$, the recipient $\vaultEvmAddr$, $\amtMax$, $\amtOut$, and the resulting movements through the pool.

\smallskip
\textbf{On $\Lshd$:} as $\Withdraw$, over $\swapMap$ and $\swapMarkerMap$, with two mint
records at settlement. Mints do not leak any information about the payer and receiver.

\smallskip
\textbf{The executed price is public.} $\amtMax \in \req$ and $\amtIn \in \stmt_{\cktCompleteSwap}$
outright, and $\amtChg$ is published as a mint record as well. With both tokens, $\feeTier$ and the block position also public,every vault swap is a
fully transparent trade.  No address -  besides the vault address  - is involved in the trade.
\end{PubBox}


\subsection{\texorpdfstring{$\Supply$ and $\Redeem$}{Supply and Redeem}}
\label{ssec:lendinst}

The lending flow $\Supply$ and $\Redeem$ moves value between an ERC20 and its interest-bearing wrapper, both pinned at
$\Setup$ (\S\ref{ssec:state}), the token being lent is
at address $\underAddr$, the receipt token is at address $\wrapAddr$.
$\Supply$ and $\Redeem$ are both a $\Swap$.
$\Supply$ burns a coin of the underlying
color and mints one of the wrapper color; $\Redeem$ does the reverse.  They differ from
$\Swap$ in  that  \textbf{the minted amount is read from the attestation, not from
the request}. In $\Swap$, $\amtOut$ is a field of $\req$ that the caller fixed before
anything was burned; in $\Supply$ and $\Redeem$ the corresponding quantity is whatever the
lending market returned, and it reaches the circuit only through $\evmOut$.


Throughout the paper we use:
\begin{align*}
  \colrU &\;=\; \lnk{fn:tokenTypeOf}{\tokenTypeOf}\bigl(\domSep(\underAddr),\, \vaultAddr\bigr), \\
  \colrW &\;=\; \lnk{fn:tokenTypeOf}{\tokenTypeOf}\bigl(\domSep(\wrapAddr),\, \vaultAddr\bigr)
\end{align*}
for the two shielded colors involved: $\colrU$ the color of the lent ERC20, $\colrW$ the color of
the wrapper.

\subsubsection{Supply}

\begin{ProcBox}{$\Pi.\Supply^{\Lpub, \Lshd}\bigl(\shdAcct\bc \tokType\bc \amt\bc \auxx \,;\,
      \privIn\bigr)$}
\label{box:supply}
\sloppy
\emph{Private input.} $\privIn$ holds $\usersk$; the secret key $\coinsk$ associated to the
\zkswap\ account $\srcAcct$; a fresh \freshrn{$\burnNonce$}, used at Step~1; and a fresh
\freshrn{$\mintNonce$}, used at Step~3 on success or at Step~4 on failure.\\
\emph{Public arguments.} $\tokType = \underAddr$, read from pinned state rather than supplied;
$\auxx = (\evmNonce, \keyVer)$ only. The gas envelope is contract-fixed, because the
\emph{vault}'s own EVM account pays for this transaction. Requires $\cktApproveWrap$
(\cktref{ckt:approvestata}) to have been run once, so that the wrapper may pull the underlying.

\begin{Steps}
\item \textbf{Request call: $\cktStartSupply$} (\cktref{ckt:startsupply}). The caller hands the contract a
      coin descriptor $\coinIn = (\burnNonce,\, \colrU,\, \amt)$; the circuit burns it, files
      ``$\text{deposit}(\amt, \vaultEvmAddr)$ on $\wrapAddr$'' under $\rid$ in $\supplyMap$ with
      $\keyPath = \keyPathLit$, and pins
      $\supplyMarkerMap[\rid] \gets \lnk{fn:refundComm}{\refundComm(\usersk, \rid)}$. The caller address does not appear in the transaction.
\item \textbf{MPC.} The committee signs with the key derived at $\keyPath = \keyPathLit$, so from
      $\vaultEvmAddr$, the pool's own account; the relayer broadcasts; the committee observes the
      execution, reads the returned share count against $\schemaIn$, re-serialises it against
      $\schemaOut$, and posts $\attSig$ over
      $\lnk{fn:attDigest}{\attDigest(\rid, \evmOut)}$.
\item \textbf{Settlement call on success: $\cktCompleteSupply$} (\cktref{ckt:completesupply}).
      Supplier-gated; mints \emph{one} coin, of color $\colrW$ and value $\amtShares$ read from
      $\evmOut$.
\item \textbf{Settlement call on failure:  execute $\cktRefundSupply$}   (\cktref{ckt:refundsupply}):  Re-mints the \emph{full} $\amt$ of color $\colrU$ - the deposit never executed, so nothing was spent.
\end{Steps}

\medskip
Outcome: $\succeeded$ via Step~3, $\failed$ via Step~4. Private output: $\usersk$, and the new coins minted.
\end{ProcBox}


\begin{CktBox}{ckt:approvestata}{$\cktApproveStata$}
\fld{Arguments.} $\evmNonce$, $\keyVer$.
\fld{Witness.} none.
\fld{Reads.} $\initFlag$, $\underAddr$, $\wrapAddr$, $\vaultAddr$, $\sigNonce$, $\chainId$.
\fld{Steps.}
\begin{enumerate}[nosep,leftmargin=1.7em,topsep=2pt]
\item $\assert \initFlag \geq 1$.
\item $\evmTxOf \gets$ ``approve $\wrapAddr$ for an unlimited amount'', sent to $\underAddr$,
      contract-fixed gas envelope.
\item $\keyPath \gets \keyPathLit$; \;
      $\req \gets (\vaultAddr\bc \sigNonce\bc \keyVer\bc \keyPath\bc \evmTxOf\bc \chainId)$; \;
      $\rid \gets \reqIdOf(\req)$; \; $\assert \rid \notin \reqMap$.
\end{enumerate}
\begin{OutList}
\oitem{\Dvault} $\sigNonce \mathrel{+}= 1$; \; $\reqMap[\rid] \gets \req$.
\oitem{\Dshd} $\emptyset$.
\oitem{\stmt} $\bigl(\cktApproveStata,\; \req,\; \rid,\; \notif(\vaultAddr, \rid, \reqMap),\;
  \comIO\bigr)$.
\oitem{\wit} $\bot$.
\oitem{\Rel} the request is well formed.
\end{OutList}
\end{CktBox}


\begin{CktBox}{ckt:startsupply}{$\srcStartSupply: \cktStartSupply$}
\fld{Arguments.} $\evmNonce$, $\keyVer$, $\amt$, and
$\coinIn = (\burnNonce, \mathsf{color}, \mathsf{value})$.
\fld{Witness.} $\usersk$ (and $\coinsk$ at the wallet layer).
\fld{Reads.} $\initFlag$, $\underAddr$, $\wrapAddr$, $\vaultEvmAddr$, $\vaultAddr$, $\sigNonce$,
$\chainId$.
\fld{Steps.}
\begin{enumerate}[nosep,leftmargin=1.7em,topsep=2pt]
\item $\assert \initFlag \geq 1$, \; $0 < \amt \leq 2^{64}-1$.
\item $\assert \coinIn.\mathsf{color} = \colrU$ \; and \; $\coinIn.\mathsf{value} = \amt$.
\item $\evmTxOf \gets$ ``$\srcDepositSel$ for $\amt$, recipient $\vaultEvmAddr$'', sent to
      $\wrapAddr$, contract-fixed gas envelope.
\item $\keyPath \gets \keyPathLit$.
\item $\req \gets (\vaultAddr\bc \sigNonce\bc \keyVer\bc \keyPath\bc \evmTxOf\bc \chainId)$; \;
      $\rid \gets \lnk{fn:reqIdOf}{\reqIdOf(\req)}$; \; $\assert \rid \notin \supplyMap$.
\end{enumerate}
\MidNote{\textbf{Midnight operations: burn.}
$\mnhl{\code{receiveShielded}}(\coinIn)$ then
$\mnhl{\code{sendImmediateShielded}}(\coinIn, \burnAddr, \amt)$.}
\begin{OutList}
\oitem{\Dvault} $\sigNonce \mathrel{+}= 1$; \; $\supplyMap[\rid] \gets \req$; \;
  \ADDEDINL{$\supplyView[\rid] \gets \bigl(\lnk{fn:refundComm}{\refundComm(\usersk, \rid)},\;
  \amt\bigr)$}.
\oitem{\Dshd} \bluedot $\coinCommit(\coinIn, \vaultAddr)$, \;
  \bluedot $\coinNull(\coinIn, \vaultAddr)$, \;
  \bluedot $\coinCommit(\coinBurnOut, \burnAddr)$, \;
  \bluedot $\{\coinNull(\coinFund, \coinsk)\}$ and $\coinCommit(\coinChange, \srcAcct)$.
\oitem{\stmt} $\bigl(\cktStartSupply,\; \req,\; \rid,\; \supplyView[\rid],\; \Dshd,\;
  \notif(\vaultAddr, \rid, \supplyMap),\; \comIO\bigr)$.
\oitem{\wit} $(\usersk,\; \burnNonce,\; \coinsk)$.
\oitem{\Rel} $\coinIn.\mathsf{color} = \colrU$ \;$\wedge$\; $\coinIn.\mathsf{value} = \amt$
  \;$\wedge$\; $\supplyView[\rid].\fComm = \refundComm(\usersk, \rid)$.
\end{OutList}
\end{CktBox}


\begin{CktBox}{ckt:completesupply}{$\cktCompleteSupply$}
\fld{Arguments.} $\rid$, $\attSig$, $\evmOut$ of width $8$, $\freshrn{\mintNonce}$.
\fld{Witness.} $\usersk$, $\mintNonce$.
\fld{Reads.} $\initFlag$, $\mpcRespKey$, \ADDEDINL{$\supplyView[\rid]$}, $\supplyMap$, $\wrapAddr$.
\fld{Steps.}
\begin{enumerate}[nosep,leftmargin=1.7em,topsep=2pt]
\item $\assert \initFlag \geq 1$ and
      $\assert \EcdsaVf\bigl(\lnk{fn:attDigest}{\attDigest(\rid, \evmOut)}\bc \attSig\bc
      \mpcRespKey\bigr)$.
\item $\assert \rid \in \supplyMap$; \; remove $\supplyMap[\rid]$.
\item \ADDEDINL{$\assert \lnk{fn:refundComm}{\refundComm(\usersk, \rid)} =
      \supplyView[\rid].\fComm$} (authentication \textbf{supplier-only}); \; remove $\supplyView[\rid]$.
\item $\amtShares \gets$ the attested value in $\evmOut$, released publicly.
\end{enumerate}
\MidNote{\textbf{Midnight operation.} Mint
$\coinMintW = \bigl(\freshrn{\mintNonce},\; \colrW,\; \amtShares\bigr)$
to $\usrTag{\coinpk \text{ of } \setlAcct}$.}
\begin{OutList}
\oitem{\Dvault} $\supplyMap[\rid]$ removed; \ADDEDINL{$\supplyView[\rid]$ removed}.
\oitem{\Dshd} \bluedot mint record $\bigl(\domSep(\wrapAddr),\, \amtShares\bigr)$; \;
  \bluedot $\coinCommit(\coinMintW$, $\coinpk$ of $\shdAcct)$.
\oitem{\stmt} $\bigl(\cktCompleteSupply,\; \rid,\; \attSig,\; \amtShares,\; \Dvault,\;
  \Dshd\bigr)$.
\oitem{\wit} $(\usersk,\; \mintNonce)$.
\oitem{\Rel} $\refundComm(\usersk, \rid)$ $= \supplyView[\rid].\fComm$ \;$\wedge$\;
  $\EcdsaVf(\attDigest(\rid,\evmOut)$, $\attSig$, $\mpcRespKey)$ \;$\wedge$\;
  the commitment opens to
  $\bigl(\mintNonce,\, \colrW,\, \amtShares,\, \coinpk\bigr)$.
\end{OutList}
\end{CktBox}


\begin{CktBox}{ckt:refundsupply}{$\cktRefundSupply$}
\fld{Arguments.} $\rid$, $\attSig$, $\evmOut$ of with 5, $\mintNonce$.
\fld{Witness.} $\usersk$.
\fld{Reads.} $\initFlag$, $\mpcRespKey$, $\supplyView[\rid]$, $\supplyMap$.
\fld{Steps.}
\begin{enumerate}[nosep,leftmargin=1.7em,topsep=2pt]
\item Run \hyperref[ckt:failgate]{$\cktFailGate(\rid, \attSig, \evmOut)$}.
\item $\assert \rid \in \supplyView$.
\item $\assert \lnk{fn:refundComm}{\refundComm(\usersk, \rid)} = \supplyView[\rid].\fComm$ - (authentication:
      \textbf{supplier-only}).
\item Remove $\supplyMap[\rid]$ and $\supplyView[\rid]$.
\end{enumerate}
\MidNote{\textbf{Midnight operation.} Mint
$\coinMint = \bigl(\freshrn{\mintNonce},\; \colrU,\; \supplyView[\rid].\fAmt\bigr)$
to $\usrTag{\coinpk \text{ of } \setlAcct}$.}
\begin{OutList}
\oitem{\Dvault} $\supplyMap[\rid]$ removed; $\supplyView[\rid]$ removed.
\oitem{\Dshd} \bluedot mint record
  $\bigl(\domSep(\underAddr),\, \supplyView[\rid].\fAmt\bigr)$; \;
  \bluedot $\coinCommit\bigl(\coinMint,\, \usrTag{\coinpk \text{ of } \setlAcct}\bigr)$.
\oitem{\stmt} $\bigl(\cktRefundSupply,\; \rid,\; \attSig,\; \Dvault,\; \Dshd\bigr)$.
\oitem{\wit} $(\usersk,\; \mintNonce)$.
\oitem{\Rel} $\refundComm(\usersk, \rid)$ $= \supplyView[\rid].\fComm$ \;$\wedge$\; the relation of
  $\cktFailGate$ \;$\wedge$\; the commitment opens to
  $\bigl(\mintNonce,\, \colrU,\, \supplyView[\rid].\fAmt,\, \usrTag{\coinpk \text{ of }
  \shdAcct}\bigr)$.
\end{OutList}
\end{CktBox}

\begin{PubBox}{What $\Supply$ publishes}
\textbf{On $\Lpub$:} the supply in the clear: the wrapper contract called, the amount of the
underlying supplied, the receiver $\vaultEvmAddr$, and the share count returned. The vault's own account is the sender, so \emph{nothing tied to $\srcAcct$ appears}.

\smallskip
\textbf{On $\Lshd$:} $\stmt_{\cktSupply}$, the MPC signatures, and the settlement statement.
Of these, the components that involve the caller at all are: the marker
$\refundComm(\usersk, \rid)$; the coins burned to the contract: the nullifier and the new coin commitment destined to $\burnAddr$ (the private and random nonce $\burnNonce$ protects the identity of the caller of $\cktStartSupply$), the new private coins minted in the $\cktCompleteSupply/\cktRefundSupply$ destined  to $\shdAcct$ (the private and random nonce $\mintNonce$ protects the identity of the recipient).

\smallskip
\textbf{The exchange rate is public.} $\amt$ is a field of $\req$ and $\amtShares$ is
released by the settlement call and published again as a mint record, so
$\amt / \amtShares$. This is
information about $\Lpub$, not about the caller, and it is identical for every supplier in the same
block.

\smallskip
\textbf{Hidden:} which shielded account funded the call; which shielded account received the shares.\end{PubBox}


\subsubsection{Redeem}

\begin{ProcBox}{$\Pi.\Redeem^{\Lpub, \Lshd}\bigl(\shdAcct\bc \wrapOf{\tokType}\bc \amtShares\bc
      \auxx \,;\, \privIn\bigr)$}
\label{box:redeem}
\sloppy
\emph{Private input.} As $\Supply$: $\usersk$; $\coinsk$ of $\srcAcct$; a fresh
\freshrn{$\burnNonce$} at Step~1; and a fresh \freshrn{$\mintNonce$} at Step~3 on success or Step~4
on failure.\\
\emph{Public arguments.} $\wrapOf{\tokType} = \wrapAddr$, read from pinned state;
$\auxx = (\evmNonce, \keyVer)$ only. No approval leg is needed: the vault redeems its own shares,
so no allowance to a third party is involved.

\begin{Steps}
\item \textbf{Request call: $\cktRedeem$} (\cktref{ckt:startredeem}). The caller hands the contract a coin
      descriptor $\coinIn = (\burnNonce,\, \colrW,\, \amtShares)$; the circuit burns it, files
      ``$\text{redeem}(\amtShares, \vaultEvmAddr, \vaultEvmAddr)$ on $\wrapAddr$'' under $\rid$ in
      $\redeemMap$ with $\keyPath = \keyPathLit$, and pins
      $\redeemMarkerMap[\rid] \gets \lnk{fn:refundComm}{\refundComm(\usersk, \rid)}$.
\item \textbf{MPC.} As in $\Supply$, signing from $\vaultEvmAddr$. The committee reports
      $\amtAssets$, the amount of the underlying the market returned --- principal plus accrued
      interest.
\item \textbf{Settlement call on success: $\cktCompleteRedeem$} (\cktref{ckt:completeredeem}).
      Redeemer-gated; mints \emph{one} coin, of colour $\colrU$ and value $\amtAssets$.
\item \textbf{Settlement call on failure: $\cktRefundRedeem$}
      (\cktref{ckt:refundredeem}). Re-mints the \emph{full} $\amtShares$ of color $\colrW$.
\end{Steps}

\medskip
Outcome: $\succeeded$ via Step~3, $\failed$ via Step~4, $\pending$ otherwise. Private output:
$\usersk$.
\end{ProcBox}

% \begin{CktBox}{ckt:redeem}{$\cktRedeem$ --- the request call of $\Redeem$}
% \fld{Arguments.} $\evmNonce$, $\keyVer$, $\amtShares$, and the coin to surrender
% $\coinIn = (\burnNonce,\; \colr,\; \mathsf{value})$.
% \fld{Witness.} $\usersk$; and, at the wallet layer, $\coinsk$ of $\srcAcct$ and the openings of the
% coins the wallet spends.
% \fld{Reads.} $\initFlag$, $\wrapAddr$, $\vaultEvmAddr$, $\vaultAddr$, $\sigNonce$, $\chainId$.
% \fld{Steps.}
% \begin{enumerate}[nosep,leftmargin=1.7em,topsep=2pt]
% \item $\assert \initFlag \geq 1$, \; $0 < \amtShares \leq 2^{64}-1$. \textbf{This bounds the
%       \emph{input}, and only the input.} The refund path re-mints $\amtShares$, so the bound makes
%       the refund always representable; it says nothing about $\amtAssets$, which is what the success
%       path mints. That asymmetry is the subject of Warning~\ref{warn:redeembound}.
% \item $\assert \coinIn.\colr = \colrW$ \; and \; $\coinIn.\mathsf{value} = \amtShares$. As in
%       $\cktSupply$, both compared values are public and $\burnNonce$ is unconstrained.
% \item $\evmTxOf \gets$ ``$\srcRedeemSel$: redeem $\amtShares$ shares, receiver $\vaultEvmAddr$,
%       owner $\vaultEvmAddr$'', sent to $\wrapAddr$, on $\chainId$, no native value,
%       \textbf{contract-fixed} gas envelope. \emph{Both the receiver and the owner are read from
%       pinned state}: the shares being redeemed are the pool's, and the assets return to the pool.
% \item $\keyPath \gets \keyPathLit$.
% \item $\req \gets (\vaultAddr\bc \sigNonce\bc \keyVer\bc \keyPath\bc \evmTxOf\bc \chainId\bc
%       \schemaIn\bc \schemaOut)$, with the two contract-fixed redeem schemas; \;
%       $\rid \gets \lnk{fn:reqIdOf}{\reqIdOf(\req)}$; \; $\assert \rid \notin \redeemMap$.
% \end{enumerate}

% \MidNote{\textbf{Midnight operations: burn redeemed coin}, as in \cktref{ckt:supply}:
% $\mnhl{\opReceive}(\coinIn)$ then $\mnhl{\opSendImm}(\coinIn, \burnAddr, \amtShares)$.}

% \smallskip
% \noindent\textbf{The four coins this call touches.} Identical in structure to the table of
% \cktref{ckt:supply}, with $\colrW$ in place of $\colrU$ and $\amtShares$ in place of $\amt$:
% $\coinIn = (\burnNonce,\, \colrW,\, \amtShares)$;
% $\coinFund_i$ of colour $\colrW$ summing to $V \geq \amtShares$;
% $\coinChange = (\cdot,\, \colrW,\, V - \amtShares)$;
% $\coinBurnOut = \bigl(\evolveNonce(\burnNonce),\, \colrW,\, \amtShares\bigr)$.

% \begin{OutList}
% \oitem{\Dvault} $\sigNonce \mathrel{+}= 1$; \; $\redeemMap[\rid] \gets \req$; \;
%   $\redeemMarkerMap[\rid] \gets \lnk{fn:refundComm}{\refundComm(\usersk, \rid)}$.
% \oitem{\Dshd} exactly the five entries of \cktref{ckt:supply}, with $\colrW$ and $\amtShares$
%   substituted: \; \bluedot $\{\coinNull(\coinFund_i, \usrTag{\coinsk})\}_i$, \;
%   \bluedot $\coinNull(\coinIn, \ctrTag{\vaultAddr})$, \;
%   \bluedot $\coinCommit(\coinIn, \ctrTag{\vaultAddr})$, \;
%   \bluedot $\coinCommit(\coinBurnOut, \burnAddr)$, \;
%   \bluedot $\coinCommit(\coinChange, \usrTag{\coinpk})$.
% \oitem{\stmt} $\bigl(\cktRedeem,$ $\req,$ $\rid,$ $\refundComm(\usersk, \rid),$ $\Dshd,$
%   $\notif(\vaultAddr, \rid, \redeemMap),$ $\comIO\bigr)$, where $\req$ exposes
%   $\wrapAddr, \amtShares, \vaultEvmAddr, \evmNonce, \sigNonce$ and $\keyPath = \keyPathLit$. The
%   notification \textbf{names $\redeemMap$}.
% \oitem{\wit} $\bigl(\usersk,\; \burnNonce,\; \coinsk,\; \{\coinFund_i\}_i\bigr)$.
% \oitem{\Rel} $\redeemMarkerMap[\rid] = \refundComm(\usersk, \rid)$ \;$\wedge$\;
%   $\coinIn = (\burnNonce, \colrW, \amtShares)$ opens both
%   $\coinCommit(\coinIn, \ctrTag{\vaultAddr})$ and $\coinNull(\coinIn, \ctrTag{\vaultAddr})$
%   \;$\wedge$\; each $\coinFund_i$ opens its published nullifier under $\coinsk$ and is unspent
%   \;$\wedge$\; the colour balances.
% \end{OutList}
% \end{CktBox}

% \ADDEDMARK{Current source: \srcStartRedeem. Changes from Circuit~C15: the pending marker is now
% $\redeemView[\rid]$, a record holding the caller commitment and the share count, in place of
% $\redeemMarkerMap[\rid]$, which held the commitment alone.}
\begin{CktBox}{ckt:startredeem}{$\cktStartRedeem$ : (\srcStartRedeem)}
\fld{Arguments.} $\evmNonce$, $\keyVer$, $\amtShares$, and
$\coinIn = (\burnNonce, \mathsf{color}, \mathsf{value})$.
\fld{Witness.} $\usersk$ (and $\coinsk$ at the wallet layer).
\fld{Reads.} $\initFlag$, $\wrapAddr$, $\vaultEvmAddr$, $\vaultAddr$, $\sigNonce$, $\chainId$.
\fld{Steps.}
\begin{enumerate}[nosep,leftmargin=1.7em,topsep=2pt]
\item $\assert \initFlag \geq 1$, \; $0 < \amtShares \leq 2^{64}-1$.
\item $\assert \coinIn.\mathsf{color} = \colrW$ \; and \; $\coinIn.\mathsf{value} = \amtShares$.
\item $\evmTxOf \gets$ ``$\srcRedeemSel$ for $\amtShares$, receiver $\vaultEvmAddr$, owner
      $\vaultEvmAddr$'', sent to $\wrapAddr$, contract-fixed gas envelope.
\item $\keyPath \gets \keyPathLit$.
\item $\req \gets (\vaultAddr\bc \sigNonce\bc \keyVer\bc \keyPath\bc \evmTxOf\bc \chainId)$; \;
      $\rid \gets \lnk{fn:reqIdOf}{\reqIdOf(\req)}$; \; $\assert \rid \notin \redeemMap$.
\end{enumerate}
\MidNote{\textbf{Midnight operations: burn.}
$\mnhl{\code{receiveShielded}}(\coinIn)$ then
$\mnhl{\code{sendImmediateShielded}}(\coinIn, \burnAddr, \amtShares)$.}
\begin{OutList}
\oitem{\Dvault} $\sigNonce \mathrel{+}= 1$; \; $\redeemMap[\rid] \gets \req$; \;
  \ADDEDINL{$\redeemView[\rid] \gets \bigl(\lnk{fn:refundComm}{\refundComm(\usersk, \rid)},\;
  \amtShares\bigr)$}.
\oitem{\Dshd} \bluedot $\coinCommit(\coinIn, \vaultAddr)$, \;
  \bluedot $\coinNull(\coinIn, \vaultAddr)$, \;
  \bluedot $\coinCommit(\coinBurnOut, \burnAddr)$, \;
  \bluedot $\{\coinNull(\coinFund_i, \coinsk)\}_i$ and $\coinCommit(\coinChange, \srcAcct)$.
\oitem{\stmt} $\bigl(\cktStartRedeem,\; \req,\; \rid,\; \redeemView[\rid],\; \Dshd,\;
  \notif(\vaultAddr, \rid, \redeemMap),\; \comIO\bigr)$.
\oitem{\wit} $(\usersk,\; \burnNonce,\; \coinsk)$.
\oitem{\Rel} $\coinIn.\mathsf{color} = \colrW$ \;$\wedge$\; $\coinIn.\mathsf{value} = \amtShares$
  \;$\wedge$\; $\redeemView[\rid].\fComm = \refundComm(\usersk, \rid)$.
\end{OutList}
\end{CktBox}


\begin{CktBox}{ckt:completeredeem}{$\cktCompleteRedeem$}
\fld{Arguments.} $\rid$, $\attSig$, $\evmOut$ of width $8$, $\freshrn{\mintNonce}$.
\fld{Witness.} $\usersk$, $\mintNonce$.
\fld{Reads.} $\initFlag$, $\mpcRespKey$, \ADDEDINL{$\redeemView[\rid]$}, $\redeemMap$,
$\underAddr$.
\fld{Steps.}
\begin{enumerate}[nosep,leftmargin=1.7em,topsep=2pt]
\item $\assert \initFlag \geq 1$ and
      $\assert \EcdsaVf\bigl(\lnk{fn:attDigest}{\attDigest(\rid, \evmOut)}\bc \attSig\bc
      \mpcRespKey\bigr)$.
\item $\assert \rid \in \redeemMap$; \; remove $\redeemMap[\rid]$.
\item \ADDEDINL{$\assert \lnk{fn:refundComm}{\refundComm(\usersk, \rid)} =
      \redeemView[\rid].\fComm$} --- \textbf{redeemer-only}; \; remove $\redeemView[\rid]$.
\item $\amtAssets \gets$ the attested value in $\evmOut$, \textbf{released publicly}.
\end{enumerate}
\MidNote{\textbf{Midnight operation.} Mint
$\coinMintU = \bigl(\freshrn{\mintNonce},\; \colrU,\; \amtAssets\bigr)$
to $\coinpk$.}
\begin{OutList}
\oitem{\Dvault} $\redeemMap[\rid]$ removed; \ADDEDINL{$\redeemView[\rid]$ removed}.
\oitem{\Dshd} \bluedot mint record $\bigl(\domSep(\underAddr),\, \amtAssets\bigr)$; \;
  \bluedot $\coinCommit\bigl(\coinMintU,\, \coinpk\bigr)$.
\oitem{\stmt} $\bigl(\cktCompleteRedeem,\; \rid,\; \attSig,\; \amtAssets,\; \Dvault,\;
  \Dshd\bigr)$.
\oitem{\wit} $(\usersk,\; \mintNonce)$.
\oitem{\Rel} $\refundComm(\usersk, \rid)$ $= \redeemView[\rid].\fComm$ \;$\wedge$\;
  $\EcdsaVf(\attDigest(\rid,\evmOut)$, $\attSig$, $\mpcRespKey)$ \;$\wedge$\;
  the commitment opens to
  $\bigl(\mintNonce,\, \colrU,\, \amtAssets,\, \coinpk\bigr)$.
\end{OutList}
\end{CktBox}


\begin{CktBox}{ckt:refundredeem}{$\cktRefundRedeem$}
\fld{Arguments.} $\rid$, $\attSig$, $\evmOut$ of width $5$, $\mintNonce$.
\fld{Witness.} $\usersk$.
\fld{Reads.} $\initFlag$, $\mpcRespKey$, $\redeemView[\rid]$, $\redeemMap$.
\fld{Steps.}
\begin{enumerate}[nosep,leftmargin=1.7em,topsep=2pt]
\item Run \hyperref[ckt:failgate]{$\cktFailGate(\rid, \attSig, \evmOut)$}.

\item $\assert \rid \in \redeemView$.
\item $\assert \lnk{fn:refundComm}{\refundComm(\usersk, \rid)} = \redeemView[\rid].\fComm$ (authenticating
      \textbf{redeemer-only}).
\item Remove $\redeemMap[\rid]$ and $\redeemView[\rid]$.
\end{enumerate}
\MidNote{\textbf{Midnight operation.} Mint
$\coinMint = \bigl(\freshrn{\mintNonce},\; \colrW,\; \redeemView[\rid].\fShares\bigr)$
to $\usrTag{\coinpk \text{ of } \setlAcct}$.}
\begin{OutList}
\oitem{\Dvault} $\redeemMap[\rid]$ removed; $\redeemView[\rid]$ removed.
\oitem{\Dshd} \bluedot mint record
  $\bigl(\domSep(\wrapAddr),\, \redeemView[\rid].\fShares\bigr)$; \;
  \bluedot $\coinCommit\bigl(\coinMint,\, \usrTag{\coinpk \text{ of } \setlAcct}\bigr)$.
\oitem{\stmt} $\bigl(\cktRefundRedeem,\; \rid,\; \attSig,\; \Dvault,\; \Dshd\bigr)$.
\oitem{\wit} $(\usersk,\; \mintNonce)$.
\oitem{\Rel} $\refundComm(\usersk, \rid)$ $= \redeemView[\rid].\fComm$ \;$\wedge$\; the relation of
  $\cktFailGate$ \;$\wedge$\; the commitment opens to
  $\bigl(\mintNonce,\, \colrW,\, \redeemView[\rid].\fShares,\, \usrTag{\coinpk \text{ of }
  \setlAcct}\bigr)$.
\end{OutList}
\end{CktBox}
\begin{PubBox}{What $\Redeem$ publishes}
\textbf{On $\Lpub$:} the redemption in the clear: the wrapper called, the share count burned,  and the amount of the underlying returned. The vault's account $\vaultEvmAddr$ is involved in all transactions on $\Lpub$ so nothing tied to $\srcAcct$ appears.

\smallskip
\textbf{On $\Lshd$:} as $\Supply$, over $\redeemMap$ and $\redeemMarkerMap$, with the two colors exchanged.

\smallskip
\textbf{The accrued yield is public.} $\amtShares$ is a field of $\req$ and $\amtAssets$ is released
by the settlement call, so the realised rate is recoverable, and with it the interest the position
earned.

\smallskip
\textbf{Hidden:} which shielded account funded the call; which received the assets.
\end{PubBox}

%-----------------------------------------------------------------------------------------------------
%Section: Cryptographic assuptions
%-----------------------------------------------------------------------------------------------------
\section{Cryptographic Assumptions}
\label{sec:assumptions}

Before presenting a formal security proof for the unlikability (\S~\ref{sec:proof}) and soundness (\S~\ref{sec:soundness})
properties of the protocol described in Section~\ref{sec:instantiation},
we define the cryptographic assumptions that we use in the proofs.
Most of the cryptographic assumptions we list below are inherited from the Midnight shielded-circuit and shielded-coin layer..


\vspace{5pt}
\noindent
\textbf{Notation}
$\secpar$ is the security parameter; PPT means probabilistic polynomial time in $\secpar$;
$\negl(\cdot)$ denotes a negligible function. $x \gets \mathcal{D}$ is sampling, $y \gets F(x)$ the output of a randomized algorithm, $\Adv^{\mathcal{O}}$ an adversary with access to the oracle set
$\mathcal{O}$.
Notation  $p_a\approx_{\negl}  p_b$ means that probability distribution $p_a$ is equal to $p_a$ up to a negligible factor.


\subsection{(Multi-use)  Pseudorandom functions}
\label{ssec:prf}

Let $\PRFfam : \prfKeySp \times \prfDom \to \prfRng$ be a family of functions.

\begin{definition}[PRF security]
\label{def:prf}
For a PPT adversary $\Adv$ define
\[
\AdvPRF_{\PRFfam}(\Adv,\secpar) \;=\;
\Bigl|\ \Pr\bigl[\,\prfkey \gets \prfKeySp \ :\ \Adv^{\PRFfam(\prfkey,\cdot)}(1^{\secpar}) = 1\,\bigr]
\;-\;
\Pr\bigl[\,\RandFun \gets \mathsf{Funs}(\prfDom,\prfRng) \ :\ \Adv^{\RandFun(\cdot)}(1^{\secpar}) = 1\,\bigr]\ \Bigr|,
\]
where $\mathsf{Funs}(\prfDom,\prfRng)$ is the set of all functions from $\prfDom$ to $\prfRng$.
$\PRFfam$ is a \emph{pseudorandom function} if $\AdvPRF_{\PRFfam}(\Adv,\secpar) \le \negl(\secpar)$
for every PPT $\Adv$.
\end{definition}


\begin{definition}[Multi-user PRF Security~\cite{BBM00,BBT16}]
\label{def:muprf}
For $m = m(\secpar)$ users and a PPT adversary $\Adv$ with access to $m$ oracles, define
\[
\AdvMUPRF_{\PRFfam,m}(\Adv,\secpar) \;=\;
\Bigl|\ \Pr\bigl[\,\prfkey_1,\dots,\prfkey_m \gets \prfKeySp \ :\
\Adv^{\PRFfam(\prfkey_1,\cdot),\dots,\PRFfam(\prfkey_m,\cdot)}(1^{\secpar}) = 1\,\bigr]
\]
\[
-\;
\Pr\bigl[\,\RandFun_1,\dots,\RandFun_m \gets \mathsf{Funs}(\prfDom,\prfRng) \ :\
\Adv^{\RandFun_1(\cdot),\dots,\RandFun_m(\cdot)}(1^{\secpar}) = 1\,\bigr]\ \Bigr|.
\]
$\PRFfam$ is \emph{multi-user secure} if $\AdvMUPRF_{\PRFfam,m}(\Adv,\secpar) \le \negl(\secpar)$
for every PPT $\Adv$ and every polynomial $m$.
\end{definition}

\begin{remark}[On the need of multi-user PRF security]
\label{rem:whymuprf}
This property is not strictly needed, as any an hybrid over the $m$ users reduces any  multi-input secure PRF (Definition~\ref{def:muprf})  into a secure PRF (Definition~\ref{def:prf}) with a
loss of a factor $m$. However,   the multi-user version  is preferred to avoid this loss.
\end{remark}


\subsection{Hash Functions}
\begin{definition}[Collision-Resistant Hash Function~\cite{KatzLindell}]
\label{def:crhf}
A hash function family $\mathcal{H} = \{H_{pp} : \{0,1\}^* \to \{0,1\}^{\ell(\secpar)}\}_{pp \in \{0,1\}^{\secpar}}$ parameterized by public parameters $pp$ and security parameter $\secpar$ is \emph{collision-resistant} if for every PPT adversary $\mathcal{A}$, there exists a negligible function $\negl(\secpar)$ such that:

    \[
\Pr\left[
  pp \leftarrow \mathrm{Gen}(1^{\secpar}), \, (x, x') \leftarrow \mathcal{A}(1^{\secpar}, pp) : x \neq x' \land H_{pp}(x) = H_{pp}(x')
\right] \le \negl(\secpar)
\]
\end{definition}



\subsection{Non-interactive zero-knowledge arguments}
\label{ssec:nizkdef}
The Midnight layer produces, for every circuit call, a non-interactive argument for the relation
$\Rel_C$. We use the standard notion of a preprocessing non-interactive
argument system in the common-reference-string model~\cite{BCCT12,Gro16}.

\begin{definition}[Non-interactive argument system]
\label{def:nizk}
Let $\Rel$ be an NP relation with statements $\stmt$ and witnesses $\wit$. A triple
$\NIZK = (\zkSetup, \zkProve, \zkVerify)$ of PPT algorithms, where
$\crs \gets \zkSetup(1^{\secpar}, \Rel)$, $\zkproof \gets \zkProve(\crs, \stmt, \wit)$ and
$\zkVerify(\crs, \stmt, \zkproof) \in \{0,1\}$, is a \emph{non-interactive argument system} for
$\Rel$ if it satisfies completeness, computational soundness and zero knowledge as defined below.
\end{definition}

\begin{definition}[Completeness]
\label{def:zkcomplete}
For every $(\stmt,\wit) \in \Rel$,
$\Pr[\crs \gets \zkSetup(1^{\secpar},\Rel);\ \zkproof \gets \zkProve(\crs,\stmt,\wit) :
\zkVerify(\crs,\stmt,\zkproof) = 1] = 1$.
\end{definition}

\begin{definition}[Computational soundness]
\label{def:zksound}
For every PPT $\Adv$,
\[
\AdvSnd_{\NIZK}(\Adv,\secpar) \;=\;
\Pr\bigl[\,\crs \gets \zkSetup(1^{\secpar},\Rel);\ (\stmt,\zkproof) \gets \Adv(\crs)\ :\
\zkVerify(\crs,\stmt,\zkproof) = 1 \ \wedge\ \stmt \notin \mathcal{L}_{\Rel}\,\bigr]
\;\le\; \negl(\secpar),
\]
where $\mathcal{L}_{\Rel} = \{\stmt : \exists \wit \text{ s.t. } \Rel(\stmt,\wit) = 1\}$.
\end{definition}

\begin{definition}[Knowledge soundness]
\label{def:zkksound}
$\NIZK$ is \emph{knowledge sound} if for every PPT $\Adv$ there is a PPT extractor $\mathcal{E}$
such that
\[
\Pr\bigl[\,\crs \gets \zkSetup(1^{\secpar},\Rel);\ (\stmt,\zkproof) \gets \Adv(\crs);\
\wit \gets \mathcal{E}^{\Adv}(\crs,\stmt)\ :\
\zkVerify(\crs,\stmt,\zkproof) = 1 \ \wedge\ \Rel(\stmt,\wit) = 0\,\bigr] \le \negl(\secpar).
\]
Knowledge soundness implies Definition~\ref{def:zksound}.
\end{definition}

\begin{definition}[Zero knowledge]
\label{def:zk}
$\NIZK$ is \emph{(computationally) zero knowledge} if there is a PPT simulator
$\zkSim = (\zkSim_1, \zkSim_2)$ such that for every PPT $\Adv$,
\begin{align*}
  \AdvZK_{\NIZK}(\Adv,\secpar) &\;=\; \biggl| \Pr\bigl[\crs \gets \zkSetup(1^{\secpar},\Rel) : \Adv^{\mathcal{P}(\crs,\cdot,\cdot)}(\crs) = 1\bigr] \\
  &\qquad -\; \Pr\bigl[(\crs,\td) \gets \zkSim_1(1^{\secpar},\Rel) : \Adv^{\mathcal{S}(\crs,\td,\cdot,\cdot)}(\crs) = 1\bigr] \biggr|
\end{align*}
is negligible, where on a query $(\stmt,\wit)$ with $\Rel(\stmt,\wit)=1$ the oracle
$\mathcal{P}$ returns $\zkProve(\crs,\stmt,\wit)$ and $\mathcal{S}$ returns
$\zkSim_2(\crs,\td,\stmt)$,  that is, the simulated proof is produced \textbf{from the statement
alone}, without the witness. Both oracles return $\bot$ when $\Rel(\stmt,\wit) = 0$.
\end{definition}

\subsection{Commitment schemes}
\label{ssec:comdef}

\begin{definition}[Commitment scheme]
\label{def:commit}
A \emph{commitment scheme} $\CS = (\csGen, \csCom)$ consists of  parameter-generation function $\ck \gets \csGen(1^{\secpar})$, taking the security parameter $\secpar$ and outputting a public key $\ck$ and
a commitment function $\csCom(\ck, m; r)$ taking the parameters $\ck$ a message $m$ and randomness $r \gets \bits{\secpar}$, and outputting a commitment.
\end{definition}

\begin{definition}[Computational hiding]
\label{def:hiding}
$\CS$ is \emph{computationally hiding} if for every PPT $\Adv$,
\[
\AdvHide_{\CS}(\Adv,\secpar) =
\Bigl|\ \Pr\bigl[\ck \gets \csGen(1^{\secpar});\ (m_0,m_1,\mathsf{st}) \gets \Adv(\ck);\
r \gets \bits{\secpar}\ :\ \Adv(\mathsf{st}, \csCom(\ck,m_0;r)) = 1\bigr]
\]
\[
-\ \Pr\bigl[\ck \gets \csGen(1^{\secpar});\ (m_0,m_1,\mathsf{st}) \gets \Adv(\ck);\
r \gets \bits{\secpar}\ :\ \Adv(\mathsf{st}, \csCom(\ck,m_1;r)) = 1\bigr]\ \Bigr|
\;\le\; \negl(\secpar).
\]
\end{definition}

\begin{definition}[Computational binding]
\label{def:binding}
$\CS$ is \emph{computationally binding} if for every PPT $\Adv$,
\begin{align*}
  \AdvBind_{\CS}(\Adv,\secpar) &\;=\; \Pr\Bigl[\ck \gets \csGen(1^{\secpar});\ (m,r,m',r') \gets \Adv(\ck) : \\
  &\qquad\qquad m \neq m' \ \wedge\ \csCom(\ck,m;r) = \csCom(\ck,m';r')\Bigr] \;\le\; \negl(\secpar).
\end{align*}
\end{definition}

\noindent The coin commitment $\coinCommit$ of \S\ref{ssec:hashcomputation} is used as a commitment
in this sense, with the coin nonce playing the role of $r$ and the tuple $(\colr, \mathsf{value}, \rcpt)$ the role of $m$.
% In Midnight's instantiation $\ck$ are the parameters
% used to instantiate hash function $\PersistantHash$.
% Standard treatments of both properties are
in~\cite{KatzLindell}.

\subsection{Message authentication codes}
\label{ssec:macdef}

\begin{definition}[MAC and existential unforgeability]
\label{def:mac}
A \emph{message authentication code} is a pair $\MAC = (\macTag, \macVer)$ with
$t \gets \macTag(\prfkey, m)$ and $\macVer(\prfkey, m, t) \in \{0,1\}$, correct in the usual sense.
$\MAC$ is \emph{existentially unforgeable under chosen-message attack} if for every PPT $\Adv$,
\[
\AdvMAC_{\MAC}(\Adv,\secpar) =
\Pr\bigl[\prfkey \gets \prfKeySp;\ (m^*,t^*) \gets \Adv^{\macTag(\prfkey,\cdot)}(1^{\secpar})\ :\
\macVer(\prfkey,m^*,t^*) = 1 \ \wedge\ m^* \notin Q\bigr] \le \negl(\secpar),
\]
where $Q$ is the set of messages $\Adv$ queried. A pseudorandom function
(Definition~\ref{def:prf}) with $\macTag(\prfkey,m) = \PRFfam(\prfkey,m)$ and verification by
recomputation is such a MAC~\cite{KatzLindell}.
\end{definition}



\subsection{Key-Private Public Key Encryption}
\begin{definition}\label{def:keyprivate}[Key-Private Encryption / IK-CPA~\cite{ASIACRYPT:BelBolDes01}]
Let $\Pi = (\mathsf{KeyGen}, \mathsf{Enc}, \mathsf{Dec})$ be a public-key encryption scheme. We say $\Pi$ is key-private under chosen-plaintext attack (\textsf{IK-CPA}) if for all probabilistic polynomial-time (\textsf{PPT}) adversaries $\mathcal{A} = (\mathcal{A}_1, \mathcal{A}_2)$, the advantage
\[
\mathbf{Adv}_{\Pi, \mathcal{A}}^{\mathrm{ik\text{-}cpa}}(\lambda) = \left| \Pr\left[ \mathbf{Exp}_{\Pi, \mathcal{A}}^{\mathrm{ik\text{-}cpa}\text{-}1}(\lambda) = 1 \right] - \Pr\left[ \mathbf{Exp}_{\Pi, \mathcal{A}}^{\mathrm{ik\text{-}cpa}\text{-}0}(\lambda) = 1 \right] \right|
\]
is negligible in the security parameter $\lambda$, where the experiment is defined as follow:

\vspace{7pt}


\underline{$\mathbf{Exp}_{\Pi, \mathcal{A}}^{\mathrm{ik\text{-}cpa}\text{-}b}(\lambda)$ for $b \in \{0,1\}$}
\begin{enumerate}
    \item $(pk_0, sk_0) \gets \mathsf{KeyGen}(1^\lambda)$ and $(pk_1, sk_1) \gets \mathsf{KeyGen}(1^\lambda)$
    \item $(m, \mathrm{st}) \gets \mathcal{A}_1(pk_0, pk_1)$
    \item $c \gets \mathsf{Enc}(pk_b, m)$
    \item $b' \gets \mathcal{A}_2(c, \mathrm{st})$
    \item Return $b'$
\end{enumerate}
\end{definition}

\subsection{Honest Verifier ZK Proofs}
\begin{definition}\label{def:hvzk}[Honest-Verifier Zero-Knowledge~\cite{JC:GolOre94}]
Let $(P, V)$ be an interactive proof system for a language $L$ associated with an $\mathbf{NP}$-relation $R(x, w)$. Let $\mathrm{View}_V\langle P(w), V \rangle(x)$ denote the random variable representing $V$'s view during an execution of $\langle P(w), V \rangle$ on input $x$ with witness $w$. We say $(P, V)$ is Computational Honest-Verifier Zero-Knowledge (\textsf{c-HVZK}) if there exists a \textsf{PPT} simulator $\mathcal{S}$ such that:
\[
\{ \mathcal{S}(x) \}_{x \in L} \stackrel{c}{\approx} \{ \mathrm{View}_V\langle P(w), V \rangle(x) \}_{(x,w) \in R}
\]
where $\stackrel{c}{\approx}$ denotes computational indistinguishability over sequences indexed by $x \in L$.
\end{definition}





\subsection{Signatures with key derivation and pre-signatures}
\label{ssec:sigdef}

\begin{definition}[EUF-CMA under additive key derivation and pre-signatures]
\label{def:sig}
Let $\SIG$ = $(\sigKG$, $\sigSign$, $\sigVer)$ be a signature scheme with a public key-derivation function
$\Derive$ mapping a master public key and a derivation path to a derived public key. $\SIG$ is
\emph{unforgeable under additive key derivation with pre-signatures} if for every PPT $\Adv$ that
receives the master public key, may request signatures under any derived key on messages of its
choice, and may additionally request pre-signatures (first-round signing material generated before
the message is known),
\begin{multline*}
\AdvSIG(\Adv,\secpar) \;=\;
\Pr\bigl[\ \Adv \text{ outputs } (\keyPath^*, m^*, \sigma^*) \text{ such that }\\
\sigVer\bigl(\Derive(\mpcRootPk, \vaultAddr, \keyPath^*),\, m^*,\, \sigma^*\bigr) = 1
\ \text{ and } m^* \text{ was never queried at path } \keyPath^*\ \bigr]
\end{multline*}
is negligible. The security of ECDSA in  this model -- additive key derivation together with
pre-signatures,  is analyzed
in~\cite{GrothShoup22}.
\end{definition}








% Reference Citation
% Adapted from: Katz, J., & Lindell, Y. (2021). Introduction to Modern Cryptography (3rd ed.). CRC Press. (Chapter 5, Definition 5.2).
%-----------------------------------------------------------------------------------------------------
%-----------------------------------------------------------------------------------------------------

\section{Proof of Withdrawer Unlinkability}
\label{sec:proof}

 \newversion{ The same  proof  \emph{strategy} described in this section applies  to the latest release.
This is because the circuits it adds ($\cktFlush$, $\cktSend$, $\cktQueueAtt$) are simply \textbf{bookkeeping} circuits that label the requests differently.
 This labeling  \textbf{does not add any user-dependent information} to the
 requests, as such they do not introduce any link to the caller.
 Furthermore the circuits are \emph{permissionless} -- they can be called by anyone, without needing proof of any secret and  \emph{anonymous} -- the
identity of the caller is not revealed at application layer.
Two changes to the circuits this section does analyse must be carried through the proof: the
marker $\refundComm(\usersk, \rid)$ becomes $\ownComm(\inIndex, \usersk)$, a keyed hash of the
same secret over the caller-chosen random index $\inIndex$ in place of $\rid$ (the source
requires $\inIndex$ to be fresh per request, which the argument must assume), and
$\cktCompleteWithdraw$ proves that marker on every verdict, so its success branch is no
longer witness-free.
}


\subsection{Proof Overview}
As we discussed in Section~\ref{def:unlink}, withdrawer unlinkability informally says the following:
a cross-chain protocol is unlinkable if, for any users $U_0,U_1$, with shielded accounts, $\candZero,\candOne$  who are  \emph{equally likely to  have called a  withdraw} at time $t$ for amount $x$,   and given the \emph{transcript of all messages exchanged up to that point\footnote{In the actual definition the adversary see the transcript of all messages exchanged before and after that particular withdraw request, but the actions after the withdraw are normalized to the event that both $\candZero$ and $\candOne$ would have performed that withdraw.}}, no one should be able to infer if the withdrawer was $\candZero$ or $\candOne$ (in polynomial time).
The same definition can be extended to every action (i.e., Swap, Supply, Redeem), except Deposit.

In this section, we provide a formal proof that  protocol $\Pi$ (Section~\ref{sec:instantiation}) achieves  withdrawer unlinkability \S~\ref{def:unlink}, assuming cryptographic assumptions stated in \S~\ref{sec:assumptions} and assuming the privacy and soundness of the Midnight shielded-coin and shielded-computation layer.

\vspace{8pt}
\noindent
\emph{Proof Intuition and Approach.}
The intuition of the proof is simple: to prove that our protocol achieves unlinkability, we need to prove that  the protocol messages generated upon execution of any action call (Withdraw, Swap, Supply, Redeem) -- except Deposit --   do not leak any information about the identity of the caller.

The approach to provide such proof is to analyze the \emph{transcript} of each call: look at every message generated by our protocol,  identify what  messages \emph{encode the identity of the caller}, and then use the security assumption of the cryptographic tools employed, to conclude that such encoding hides the identity from any polynomial time adversary.


\subsection{Intuition for the Proof of  Withdrawer-unlinkability.}
A withdraw request consists of two parts: (1) the transaction that will be executed on the public chain (2) the transactions executed on the private chain. Let us call $\shdAcct_i$, the user who is executing the $\Withdraw$  flow (i.e., $\cktStartWithdraw,\cktCompleteWithdraw$).

\vspace{5pt}
\emph{Transcript on public chain $\Lpub$ (Ethereum).}
The transaction executed  on the public chain $\req$ contains the sender $\vaultEvmAddr$, the recipient of the withdraw $\destAddr^*$ and the amount. The sender is a fixed constant $\vaultEvmAddr$, independent of the caller $\shdAcct_i$.
Hence, \emph{the request on the public chain does not contain any information about the identity of the withdrawer.}
\textbf{Note. If a user re-uses  existing  evm address for $\destAddr$ instead of a fresh one, our protocol provides no unlinkability guarantees.}


\vspace{5pt}
\emph{Transcript on private chain $\Lshd$ (Midnight).}
When a user $\shdAcct_i$ of  $\Lshd$ (Midnight) starts a withdraw of $\amt$ token,  it will first create a private payment to a public burn address $\burnAddr$.
We use Midnight libraries \MidNoteT{$\mnhl{\code{receiveShielded}}(\coinIn)$} then
\MidNoteT{$\mnhl{\code{sendImmediateShielded}}(\coinIn, \burnAddr, \amt)$}.
Under the assumption that the randomness included in $\coinIn$ is secret and unpredictable to the adversary, the  Midnight shielded-coin layer guarantee that the transcript of the above calls does not leak any information about the identity of the payer ($\shdAcct_i $).
In addition to the coins generated from the Midnight-coin layer, a withdraw call, also generates an \textbf{authentication tag} $\refundComm(\usersk_i,\rid)$, where $\usersk_i$ is the application key associate to the user who is calling withdraw.  This authentication tag is computed using a pseudorandom function $\Hash. $ By the definition (Def~\ref{def:prf}, Def~\ref{def:muprf}), the outputs of $\Hash$ do not reveal any information about the secret key $\usersk_i$ used for the evaluation.
Yet, $\Hash$ is a deterministic function, so, if we call $\refundComm(\usersk_i,\rid)$ on the same $\usersk_i$ and reuse an  $\rid$ in more than one withdraw request, then the two requests are readily linked.
This is why,  our implementation relies on the fact that there is a one-to-one mapping between a requestid $\rid$ and a request.
Hence, due to the  \textbf{pseudorandomness} of $\Hash$ and the  \textbf{uniqueness} of $\rid$, $\refundComm(\usersk_i,\rid)$ does not leak any information about the caller.
Finally, if a withdraw request fails, the caller must obtain a refund from the Midnight chain.
This is done by  \MidNoteT{$\mnhl\mintShieldedToken(\mathsf{newCoin})$}. Assuming that $\mathsf{newCoin}$ contains a fresh random nonce  ($\mintNonce$) that remains private, the privacy property of Midnight shielded-coin layer guarantees that
the recipient of the refund is protected.
In addition, every execution above is attested by a zero-knowledge proof, via the Midnight private-computation layer (Kachina).
Using the privacy guarantees of the latter, we can assume that the the transcript generated by the zk proofs, does not leak any information about the secret used in the execution of every circuit above.


\vspace{5pt}
\emph{Transcript of other actions (i.e., Swap, Supply, Redeem).}
The transcript of other actions, all follow a similar structure as Withdraw, and use the same cryptographic assumptions as withdraw. Hence we informally state that these additional transacripts do not leak any information about the caller of a withdrawal.
Formally, in the proof this is implicitly obtained in our hybrid-argument proofs.
In our hybrid experiments, in every experiment we remove a secret from the transcript, and the change affects every call (not just withdrawal).


\subsection{ Assumption on Consolidated Coins in Midnight Wallets (External to Our protocol)}
\label{sub:nofingerprint}
Our proof of security  (see hybrid $\hybrid_3$ in the formal proof), uses an assumption about user's wallet spending convention.
In the proof we assume that  \textbf{a user will always keep the coins consolidated in one coin}.

Let us explain.
At very high level, a coin on Midnight is a tuple (nonce, type, value); creating it publishes a commitment, spending it publishes a nullifier.
Spending a coin on Midnight  involves  (1) publishing the nullifiers (derived from the sender's secret key);  (2) creating a new commitment for the \emph{recipient}'s coin, and  (3) attaching a zero-knowledge proof of validity of (1) and (2).
Nullifiers do not reveal anything about the amount and the spender (to a PPT Adversary), however  the \emph{number} of nullifiers used to make a shielded payment reveal  \emph{how many coins} were spent. This information could help  identifying who is the user who is spending those coins.
So, for instance, if an adversary somehow knows that Alice has 3 coins and Bob 2 (even if he does not know the value of them), by looking at the number of nullifiers released in a  payment transaction, the adversary could determine if Alice is paying.

The above power is captured by our security game, where we give the adversary  the power to choose which actions and  transactions  honest parties perform  as part of her side channel information (See Section~\ref{sec:definition}). So, an adversary could instruct Alice to receive payments of amount 2, and then 3; and Bob to receive one payment of 5.
If Alice does not always keep her coins  consolidated in one, when Alice will spend 5, she will use two nullifers (while Bob will use only 1). The number of coins leak the fact that Alice is paying.
Therefore, in this security proof \textbf{we assume that all honest users keep their wallet consolidated}.
This means, at all time, each user keeps \textbf{one private coin }(shielded UTXO). And when it spends any amount of that coin,  she always mints a change coin.


\subsection{Formal Proof}
\begin{theorem}
\emph{If $\Hash$ satisfies Multi-user PRF Security (Definition~\ref{def:muprf}), Midnight's layer commitments satisfy Hiding (Definition~\ref{def:hiding}), Midnight's circuit outputs satisfy  zero-knowledge (Definition~\ref{def:nizk}), and Midnight's nullifiers satisfy Pseudorandomness (Definition~\ref{def:prf}),  and assuming that wallet satisfies the consolidated-coin assumption (Section~\ref{sub:nofingerprint}),
then system $\Pi$ described in Section~\ref{sec:instantiation}
achieves  withdrawer unlinkability (Definition ~\ref{def:unlink}).}
\end{theorem}


\begin{proof}

Our objective is to prove that any PPT adversary $\Adv$

$$\Pr[\unlinkgame_{\Pi,\Adv}(1^{\secpar}) = 1] \le \frac{1}{2} + \negl(\secpar)$$

with $\unlinkgame$ described in Fig.~\ref{fig:unlink}.
\newcommand{\last}{6}
We prove this statement with a standard  hybrid arguments.
We start with $\hybrid_0$ corresponding to the  real world, where
experiment  $\unlinkgame_{\Pi,\Adv}$ is executed with the real protocol $\Pi$ (described in \S~\ref{sec:instantiation}). We start by assuming that $\Adv$ wins the real-world game $\unlinkgame_{\hybrid_0,\Adv}$ with probability $p_{0}$.
Then, our goal is to reach a \emph{simulated} world where the challenger is executing a protocol that is completely simulated -- no identities/secrets are used anywhere -- and claim that $\Adv$ wins in
this simulated world with the ``same'' probability  $p_{0}$ (negligibly close). Which allows us to conclude that $p_{0}$ must be close to $1/2$ since in the simulated world the protocol messages contain no secret information, and no adversary can guess which account was picked in the game with better than $1/2$ probability.

The core of the proof is to reach the \emph{simulated world} through a sequence of hybrid worlds (or hybrid experiments),
 $\hybrid_1, \hybrid_2, \ldots$  where we run the game
 $\unlinkgame_{\hybrid_i,\Adv}$ with a modified protocol $\Pi_{\hybrid_i}$.
 In  each $\Pi_{\hybrid_i}$, the protocol is modified by removing the secrets used to compute a  cryptographic object. In each hybrid world $\hybrid_i$,
we use the cryptographic properties to claim that the probability  of $\Adv$ winning the game in this world  must be  close to  the previous world $\hybrid_{i-1}$,  up to a negligible factor.
In the last hybrid game $\hybrid_{\last}$, we reach the \emph{simulated} world, a world where  no secret is used in any computation. Since secrets are not used, in such a world, the protocol messages do not leak any information about the caller -- only the public information about the transactions do.
\emph{Since the public information do not contain the identity of the caller of a withdraw, a swap, a supply or a  redeem,}, the transcript overall does not leak any information.
In this simulated world, the probability that the adversary $\Adv$ wins the game no better than just guessing.
Hence $p_{\last}= 1/2$, and the sequence of hybrid games allows use to conclude that $p_0\approx_{\negl} p_1 \approx\ldots \approx_{\negl} 1/2$ and therefore $p_0 \approx_{\negl} 1/2$, where $\approx_{\negl}$ means equal up to a negligible factor.
We provide the formal hybrid-argument proof below.


\subsubsection{Hybrid 0}
This hybrid corresponds o the game $\unlinkgame_{\Pi,\Adv}(1^{\secpar})$ described in Figure \ref{ssec:unlink}. This is the \emph{real world}, where every oracle is implemented faithfully following protocol $\Pi$.
When participating in the real-world  experiment, $\Adv$  obtains the following transcript.


\begin{itemize} [label=$\blacksquare$]
    \item\textbf{Transcript from the oracle calls  pre-challenge}:
    The adversary has oracle access to both ledger and to cross-chain oracles  (as explained in Section~\ref{sec:definition}).

          \[(\candZero\bc \candOne\bc \tokType\bc \amt\bc \auxx^*\bc \outc^*\bc
      \mathsf{st}) \gets \Adv^{\mathcal{O}}(\pp),\]


    As such $\Adv$ can request creations of honest accounts on $\Lpub$ and $\Lshd$, and
    instruct actions on behalf of those users with the cross-chain oracles: deposit, withdraw, swap, supply, redeem.
     First,  recall that a user participating in the cross-chain system holds the following keys (these are privet to the user (challenger) and not known to the adversary).:
\begin{itemize}
    \item { \bf Application secret key:} $\usersk_i$: a private key for the cross-chain system: $\usersk_i\leftarrow\{0,1\}^{\lambda}$. This is randomly chosen upon the first Deposit request.
    \item Shielded Account  ($\callerZKey_i$, $\callerZKeySk_i$) on $\Lshd$; (\zkswap\ wallet on Midnight).
    \item Public account $\pubAddr_i$  on $\Lpub$; (used to fund the derived account).
\end{itemize}
The secret keys of the simulated users are held by the challenger who simulates the experiment.


\medskip

For convenience of the reader below we report an abridged version of the the transcripts that $\Adv$ would obtain from the interactions with the cross-chain oracles $\ODep$, $\OWd$, $\OSw$, $\OSup$, $\ORed$.

For each oracle call, we highlight:
\begin{itemize}
    \item $\Dvault$: messages posted on the contract state  on the Midnight chain.
    \item $\Dshd$: message posted on the shielded-coin layer on the Midnight chain.
    \item statements, zk proofs and relations proved by the Midnight private-circuit layer.
\end{itemize}
\begin{itemize}
    \item $\ODep(\pubAddr_i, \shdAcct_i, \tokType, \amt, \auxx, \outc)$:


\begin{TxtBox}{trn:deposit}{Transcript for the $\Deposit$ flow}
Circuit \upd{$\cktStartDeposit$} (see Circuit~\ref{ckt:deposit}) \begin{OutList}
\oitem{\Dvault} $\sigNonce \mathrel{+}= 1$; \; \upd{$\depositMap[\rid] \gets \req$; \;
  $\depositView[\rid] \gets \bigl(\userComm(\secret{\usersk_i}),\, \tokType,\, \amt\bigr)$}.
\oitem{\Dshd} $\emptyset$.
\oitem{\stmt} $(\upd{\cktStartDeposit},\; \req,\; \rid,\; \upd{\depositView[\rid]},\;
  \notif(\vaultAddr, \rid, \upd{\depositMap}),\; \comIO)$, where
  $$\req = \tokType, \amt, \vaultEvmAddr, \evmNonce, \gasEnv, \sigNonce, \keyPath$$
  and
  $\keyPath_i= \userComm(\secret{\usersk_i})$.
   \oitem{\zkproofs} \bluedot
\oitem{\Rel}  \MidNoteT{$\req.\keyPath = \userComm(\secret{\usersk_i})$
  \upd{$\;\wedge\; \depositView[\rid].\fComm = \userComm(\secret{\usersk_i})$}.}
%\oitem{\wit} $\usersk_i$.
\end{OutList}
Circuit: $\cktCompleteDeposit$ (see Circuit~\ref{ckt:completedeposit})
\begin{OutList}
\oitem{\Dvault} \upd{$\depositMap[\rid]$ removed; $\depositView[\rid]$ removed}.
  \oitem{\Dshd} \bluedot
\oitem{\stmt} $\bigl(\cktCompleteDeposit,\; \rid,\; \upd{\depositView[\rid]},\; \attSig,\;
  \tokType, \amt,\; \coinClaim\bigr)$  where
  $\coinClaim=\lnk{fn:coinCommit}\coinCommit \bigl(\freshrn{\mintNonce}, \tokType,\amt,\shdAcct_i)$.
  \oitem{\zkproofs} \bluedot
 \MidNote{\oitem{\Rel} $\userComm(\usersk_i) = \upd{\depositView[\rid].\fComm}$ \;$\wedge$\;
  $\EcdsaVf( \attSig, \mpcRespKey)$ \;$\wedge$\; and $\coinClaim=\lnk{fn:coinCommit}\coinCommit \bigl(\mintNonce, \tokType,\amt,\shdAcct)$.}
  \oitem{\wit} $(\usersk_i,\; \mintNonce,\; \shdAcct)$
\end{OutList}
\end{TxtBox}

\begin{ProcBox}
    { Note:} We expect that users will re-use  the same  application secret $\usersk_i$, and hence the same $\keyPath_i$  and $\derivedAddr_i$ for every Deposit call.
\end{ProcBox}


\item $\OWd(\shdAcct_i, \pubAddr_j, \tokType, \amt, \auxx, \outc)$
%---------------------------------------------------------------------
%  $\OWd$
%---------------------------------------------------------------------
\begin{TxtBox}{trn:withdraw}{Transcript for the $\Withdraw$ flow}
 \upd{$\cktStartWithdraw$} (see Figure~\ref{ckt:startwithdraw})
\begin{OutList}
\oitem{\Dvault} $\sigNonce \mathrel{+}= 1$; \; $\reqMap[\rid] \gets \req$; \;\\
  \upd{$\wdView[\rid] \gets \bigl(\lnk{fn:refundComm}{\refundComm(\secret{\usersk_i}, \rid)},\,
  \tokType,\, \amt\bigr)$}.
\oitem{\Dshd} \emph{Nullifiers published} (one coin spent):
  \begin{itemize}[nosep,leftmargin=1.3em,topsep=1pt]
  \item \bluedot $\coinNull\bigl(\secret{\coinFund},\; \mathsf{user}(\secret{\coinsk_i})\bigr)$ (we assume consolidated-coin):
        keyed by the caller's \textbf{secret};
  \item \bluedot $\coinNull\bigl(\freshrn{\coinIn},\; \mathsf{contract}(\vaultAddr)\bigr)$: the contract
        spending what it has just received, keyed by a \textbf{public} address.
  \end{itemize}
  \emph{Commitments published} (one coin created):
  \begin{itemize}[nosep,leftmargin=1.3em,topsep=1pt]
  \item \bluedot $\coinCommit\bigl(\freshrn{\coinIn},\; \mathsf{contract}(\vaultAddr)\bigr)$ : the wallet's caller paying the contract;
  \item \bluedot $\coinCommit\bigl(\coinBurnOut,\; \burnAddr\bigr)$, over
        $\bigl(\evolveNonce(\burnNonce),\, \mathsf{color}^{*},\, \amt\bigr)$: the contract
        paying the burn address;
  \item \bluedot $\coinCommit\bigl(\coinChange,\; \mathsf{user}(\coinpk)\bigr)$, of value
        $(V - \amt)$, recipient is the caller's public coin key.
  \end{itemize}
 \MidNote{
\oitem{\stmt} $\bigl(\upd{\cktStartWithdraw},$ $\req,$ $\rid,$ $\upd{\wdView[\rid]},$
$\Dshd,$
  $\notif(\vaultAddr, \rid, \reqMap),$ $\comIO\bigr)$, where $\req$ exposes
  $\tokType, \amt, \destAddr, \evmNonce, \sigNonce$ and $\keyPath = \keyPathLit$.}
%\WitNote{\oitem{\wit} $\bigl(\usersk,\; \burnNonce,\; \coinsk,\; \{\coinFund_i\}_i\bigr)$.}
 \MidNote{\oitem{\Rel} \upd{$\wdView[\rid].\fComm = \refundComm(\usersk, \rid)$} \;$\wedge$\;
  $\coinIn = (\burnNonce, \mathsf{color}^{*}, \amt)$ opens both
  $\coinCommit(\coinIn, \mathsf{contract}(\vaultAddr))$ and
  $\coinNull(\coinIn, \mathsf{contract}(\vaultAddr))$ \;$\wedge$\;
  each $\coinFund$ opens its published nullifier under $\coinsk$ and is unspent \;$\wedge$\;
  the color balances: $V + \amt$ in, $\amt + (V-\amt) + \amt$ out.
  }
\end{OutList}

\vspace{5pt}
$\cktCompleteWithdraw$ (See circuit~\ref{ckt:completewithdraw})
\begin{OutList}
\oitem{\Dvault} $\reqMap[\rid]$ removed; \upd{$\wdView[\rid]$ removed}.
\oitem{\Dshd} $\okbit = 1$: $\emptyset$.   $\okbit = 0$: \upd{\bluedot mint record
  $\bigl(\lnk{fn:domSep}{\domSep(\wdView[\rid].\fErc)},\, \wdView[\rid].\fAmt\bigr)$; \;
  \bluedot $\lnk{fn:coinCommit}{\coinCommit(\freshrn{\mintNonce}, \wdView[\rid].\fErc,
  \wdView[\rid].\fAmt, \shdAcct_i)}$}.
\oitem{\stmt} $\bigl(\cktCompleteWithdraw,\; \rid,\; \attSig,\; \okbit,\;
  \upd{\wdView[\rid]},\; \Dshd\bigr)$.
%\oitem{\wit} $\okbit = 1$: $\bot$. \quad $\okbit = 0$: $(\usersk, \mintNonce)$.
\oitem{\Rel} $\okbit = 1$: $\EcdsaVf(\attDigest(\rid,\evmOut), \attSig, \mpcRespKey)$. \quad
  $\okbit = 0$: the same, $\wedge$ \upd{$\refundComm(\secret{\usersk_i}, \rid) =
  \wdView[\rid].\fComm$ $\wedge$ commitment opening
  $(\freshrn{\mintNonce}, \cdot, \wdView[\rid].\fAmt, \shdAcct_i)$}.
\end{OutList}

\vspace{10pt}

\upd{$\cktRefundWithdraw$} (See circuit~\ref{ckt:refundwithdraw})
\begin{OutList}
\oitem{\Dvault}  \upd{$\reqMap[\rid]$ removed; $\wdView[\rid]$ removed}.
\oitem{\Dshd} \bluedot mint record $\bigl(\domSep(\upd{\wdView[\rid].\fErc}),\,
  \upd{\wdView[\rid].\fAmt}\bigr)$; \;
  \bluedot $\lnk{fn:coinCommit}{\coinCommit(\freshrn{\mintNonce},\tokType,\amt, \shdAcct_i)}$.
\oitem{\stmt} $\bigl(\upd{\cktRefundWithdraw},\; \rid,\; \attSig,\; \upd{\wdView[\rid]},\;
  \Dvault,\; \Dshd\bigr)$.
%\oitem{\wit} $(\usersk,\; \mintNonce)$.
\oitem{\Rel} $\refundComm(\usersk, \rid) = \upd{\wdView[\rid].\fComm}$ \;$\wedge$\;
  \upd{the relation of $\cktFailGate$} \;$\wedge$\; the commitment opens to
  $(\freshrn{\mintNonce}, \cdot, \upd{\wdView[\rid].\fAmt}, \secret{\shdAcct_i})$.
\end{OutList}
\end{TxtBox}

\item $\OSw(\shdAcct\bc \tokType\bc \amt\bc \tokTypeB\bc \amtB\bc \auxx\bc \outc)$
%---------------------------------------------------------------------
%  $\OSw$
%---------------------------------------------------------------------
\begin{TxtBox}{trn:swap}{Transcript for the $\Swap$ flow.}
 Circuit \upd{$\cktStartSwap$} (see Table~\ref{ckt:startswap})
\begin{OutList}
\oitem{\Dvault} $\sigNonce \mathrel{+}= 1$; \; $\swapMap[\rid] \gets \req$; \;
  \upd{$\swapView[\rid] \gets \bigl(\lnk{fn:refundComm}{\refundComm(\usersk, \rid)},\,
  \tokType,\, \tokTypeB,\, \amtOut,\, \amtMax\bigr)$}.
\oitem{\Dshd} as \cktref{ckt:startwithdraw}, with $\amtMax$ in place of $\amt$: \;
  \bluedot $\coinCommit(\coinIn, \vaultAddr)$, \; \bluedot $\coinNull(\coinIn, \vaultAddr)$, \;
  \bluedot $\coinCommit(\coinBurnOut, \burnAddr)$, \;
  \bluedot $\{\coinNull(\coinFund_i, \coinsk_i)\}$ and $\coinCommit(\coinChange, \secret{\srcAcct_i})$.
\oitem{\stmt} $\bigl(\upd{\cktStartSwap},$ $\req,$ $\rid,$ $\upd{\swapView[\rid]},$ $\Dshd,$
  $\notif(\vaultAddr, \rid, \swapMap),$ $\comIO\bigr)$, where $\req$ exposes
  $\tokType, \tokTypeB, \feeTier, \amtOut, \amtMax$.
\oitem{\wit} $(\usersk,\; \burnNonce,\; \coinsk)$.
\oitem{\Rel} as \cktref{ckt:startwithdraw}, over \upd{$\swapView$} and $\amtMax$.
\end{OutList}

\vspace{5pt}
Circuit $\cktCompleteSwap$ (see Table~\ref{ckt:completeswap})
\begin{OutList}
\oitem{\Dvault} $\swapMap[\rid]$ removed; \upd{$\swapView[\rid]$ removed}.
\oitem{\Dshd} \bluedot mint records $\bigl(\domSep(\tokTypeB), \amtOut\bigr)$ and
  $\bigl(\domSep(\tokType), \amtChg\bigr)$; \;
  \bluedot $\coinCommit(\freshrn{\mintNonce}, \shdAcct_i)$ and
  \upd{$\coinCommit(\freshrn{\changeNonce}, \shdAcct_i)$}.
\oitem{\stmt} $\bigl(\cktCompleteSwap,\; \rid,\; \attSig,\; \amtIn,\;
  \upd{\swapView[\rid]},\; \Dshd\bigr)$.
%\oitem{\wit} $(\usersk,\; \freshrn{\mintNonce},\; \freshrn{\changeNonce})$
\oitem{\Rel} $\refundComm(\usersk, \rid) = \upd{\swapView[\rid].\fComm}$
  \;$\wedge$\; the ECDSA statement \;$\wedge$\; \upd{$\changeNonce \neq \mintNonce$}
  \;$\wedge$\; \upd{$\amtIn \leq \amtMax$} \;$\wedge$\; the two commitments open under
  $\mintNonce$ and \upd{$\changeNonce$}.
\end{OutList}

\vspace{5pt}
\upd{Circuit $\cktRefundSwap$ (see Table~\ref{ckt:refundswap})}
 \begin{OutList}
\oitem{\Dvault} \upd{$\swapMap[\rid]$ removed; $\swapView[\rid]$ removed}.
\oitem{\Dshd} \upd{\bluedot mint record $\bigl(\domSep(\swapView[\rid].\fTokIn),\,
  \swapView[\rid].\fAmtMax\bigr)$; \;
  \bluedot $\lnk{fn:coinCommit}{\coinCommit(\freshrn{\mintNonce}, \tokType, \amtMax, \shdAcct_i)}$}.
\oitem{\stmt} \upd{$\bigl(\cktRefundSwap,\; \rid,\; \attSig,\; \swapView[\rid],\; \Dvault,\;
  \Dshd\bigr)$}.
%\oitem{\wit} $(\usersk,\; \mintNonce)$.
\oitem{\Rel} \upd{$\refundComm(\usersk, \rid) = \swapView[\rid].\fComm$ \;$\wedge$\;
  the relation of $\cktFailGate$ \;$\wedge$\; the commitment opens to
  $(\freshrn{\mintNonce}, \cdot, \swapView[\rid].\fAmtMax, \secret{\shdAcct_i})$}.
\end{OutList}
\end{TxtBox}


%%%%%%%------------


\item $\OSup(\shdAcct\bc \tokType\bc \amt\bc \auxx\bc \outc)$

\begin{TxtBox}{trn:supply}{Transcript for the $\Supply$ flow.}

 Circuit \upd{$\cktStartSupply$} (see Table~\ref{ckt:startsupply})
\begin{OutList}
\oitem{\Dvault} $\sigNonce \mathrel{+}= 1$; \; $\supplyMap[\rid] \gets \req$; \;
  \upd{$\supplyView[\rid] \gets \bigl(\lnk{fn:refundComm}{\refundComm(\secret{\usersk_i}, \rid)},\,
  \amt\bigr)$}.
\oitem{\Dshd} as \cktref{ckt:startwithdraw}, with $\colrU$ in place of $\mathsf{color}^{*}$: \;
  \bluedot $\coinCommit(\coinIn, \vaultAddr)$, \; \bluedot $\coinNull(\coinIn, \vaultAddr)$, \;
  \bluedot $\coinCommit(\coinBurnOut, \burnAddr)$, \;
  \bluedot $\{\coinNull(\coinFund_i, \coinsk_i)\}$ and $\coinCommit(\coinChange, \secret{\srcAcct_i})$.
\oitem{\stmt} $\bigl(\upd{\cktStartSupply},$ $\req,$ $\rid,$ $\upd{\supplyView[\rid]},$ $\Dshd,$
  $\notif(\vaultAddr, \rid, \supplyMap),$ $\comIO\bigr)$, where $\req$ exposes
  $\amt$, and  $\underAddr$, $\wrapAddr$, $\vaultEvmAddr$, $\evmNonce$, $\sigNonce$ and
  $\keyPath = \keyPathLit$. The notification \textbf{names $\supplyMap$}, so it alone identifies
  the request as a supply.
\oitem{\wit} $(\usersk,\; \burnNonce,\; \coinsk)$.
\oitem{\Rel} as \cktref{ckt:startwithdraw}, over \upd{$\supplyView$} and $\colrU$.
\end{OutList}

\vspace{5pt}
Circuit $\cktCompleteSupply$ (see Table~\ref{ckt:completesupply})
\begin{OutList}
\oitem{\Dvault} $\supplyMap[\rid]$ removed; \upd{$\supplyView[\rid]$ removed}.
\oitem{\Dshd} \bluedot mint record $\bigl(\lnk{fn:domSep}{\domSep(\wrapAddr)},\, \amtShares\bigr)$; \;
  \bluedot $\lnk{fn:coinCommit}{\coinCommit(\freshrn{\mintNonce}, \wrapAddr, \amtShares, \shdAcct_i)}$.
\oitem{\stmt} $\bigl(\cktCompleteSupply,\; \rid,\; \attSig,\; \amtShares,\;
  \upd{\supplyView[\rid]},\; \Dvault,\;
  \Dshd\bigr)$, where $\amtShares$ is read from $\evmOut$.
%\oitem{\wit} $(\usersk,\; \freshrn{\mintNonce})$
\oitem{\Rel} $\refundComm(\usersk, \rid) = \upd{\supplyView[\rid].\fComm}$
  \;$\wedge$\; the ECDSA statement \;$\wedge$\; the commitment opens to
  $(\freshrn{\mintNonce},\, \colrW,\, \amtShares,\, \secret{\shdAcct_i})$.
\end{OutList}

\vspace{5pt}
\upd{Circuit $\cktRefundSupply$ (see Table~\ref{ckt:refundsupply})}
\begin{OutList}
\oitem{\Dvault} \upd{$\supplyMap[\rid]$ removed; $\supplyView[\rid]$ removed}.
\oitem{\Dshd} \upd{\bluedot mint record $\bigl(\lnk{fn:domSep}{\domSep(\underAddr)},\,
  \supplyView[\rid].\fAmt\bigr)$; \;
  \bluedot $\lnk{fn:coinCommit}{\coinCommit(\freshrn{\mintNonce}, \underAddr, \amt, \shdAcct_i)}$}.
\oitem{\stmt} \upd{$\bigl(\cktRefundSupply,\; \rid,\; \attSig,\; \supplyView[\rid],\; \Dvault,\;
  \Dshd\bigr)$}.
%\oitem{\wit} $(\usersk,\; \mintNonce)$.
\oitem{\Rel} \upd{$\refundComm(\usersk, \rid) = \supplyView[\rid].\fComm$ \;$\wedge$\;
  the relation of $\cktFailGate$ \;$\wedge$\; the commitment opens to
  $(\freshrn{\mintNonce},\, \colrU,\, \supplyView[\rid].\fAmt,\, \secret{\shdAcct_i})$}.
\end{OutList}

\end{TxtBox}


\item $\ORed(\shdAcct\bc \wrapOf{\tokType}\bc \amtShares\bc \auxx\bc \outc)$
\begin{TxtBox}{trn:redeem}{Transcript for the $\Redeem$ flow.}
 Circuit \upd{$\cktStartRedeem$} (see Table~\ref{ckt:startredeem})
\begin{OutList}
\oitem{\Dvault} $\sigNonce \mathrel{+}= 1$; \; $\redeemMap[\rid] \gets \req$; \;
  \upd{$\redeemView[\rid] \gets \bigl(\lnk{fn:refundComm}{\refundComm(\secret{\usersk_i}, \rid)},\,
  \amtShares\bigr)$}.
\oitem{\Dshd} as \cktref{ckt:startsupply}, with $\colrW$ in place of $\colrU$ and $\amtShares$ in place
  of $\amt$: \;
  \bluedot $\coinCommit(\coinIn, \vaultAddr)$, \; \bluedot $\coinNull(\coinIn, \vaultAddr)$, \;
  \bluedot $\coinCommit(\coinBurnOut, \burnAddr)$, \;
  \bluedot $\{\coinNull(\coinFund_i, \coinsk_i)\}$ and $\coinCommit(\coinChange, \secret{\srcAcct_i})$.
\oitem{\stmt} $\bigl(\upd{\cktStartRedeem},$ $\req,$ $\rid,$ $\upd{\redeemView[\rid]},$ $\Dshd,$
  $\notif(\vaultAddr, \rid, \redeemMap),$ $\comIO\bigr)$, where $\req$ exposes
   $\amtShares$ and $\wrapAddr$, $\vaultEvmAddr$, $\evmNonce, \sigNonce$ and $\keyPath = \keyPathLit$. The
  notification \textbf{names $\redeemMap$}.
\oitem{\wit} $(\usersk,\; \burnNonce,\; \coinsk)$.
\oitem{\Rel} as \cktref{ckt:startwithdraw}, over \upd{$\redeemView$} and $\colrW$.
\end{OutList}

\vspace{5pt}
Circuit $\cktCompleteRedeem$ (see Table~\ref{ckt:completeredeem})
\begin{OutList}
\oitem{\Dvault} $\redeemMap[\rid]$ removed; \upd{$\redeemView[\rid]$ removed}.
\oitem{\Dshd} \bluedot mint record $\bigl(\lnk{fn:domSep}{\domSep(\underAddr)},\, \amtAssets\bigr)$; \;
  \bluedot $\lnk{fn:coinCommit}{\coinCommit(\freshrn{\mintNonce}, \underAddr, \amtAssets, \shdAcct_i)}$.
\oitem{\stmt} $\bigl(\cktCompleteRedeem,\; \rid,\; \attSig,\; \amtAssets,\;
  \upd{\redeemView[\rid]},\; \Dvault,\;
  \Dshd\bigr)$, where $\amtAssets$ is read from $\evmOut$.
%\oitem{\wit} $(\usersk,\; \freshrn{\mintNonce})$
\oitem{\Rel} $\refundComm(\usersk, \rid) = \upd{\redeemView[\rid].\fComm}$
  \;$\wedge$\; the ECDSA statement \;$\wedge$\; the commitment opens to
  $(\freshrn{\mintNonce},\, \colrU,\, \amtAssets,\, \secret{\shdAcct_i})$.
\end{OutList}

\vspace{5pt}
\upd{Circuit $\cktRefundRedeem$ (see Table~\ref{ckt:refundredeem})}
\begin{OutList}
\oitem{\Dvault} \upd{$\redeemMap[\rid]$ removed; $\redeemView[\rid]$ removed}.
\oitem{\Dshd} \upd{\bluedot mint record $\bigl(\lnk{fn:domSep}{\domSep(\wrapAddr)},\,
  \redeemView[\rid].\fShares\bigr)$; \;
  \bluedot $\lnk{fn:coinCommit}{\coinCommit(\freshrn{\mintNonce}, \wrapAddr, \amtShares, \shdAcct_i)}$}.
\oitem{\stmt} \upd{$\bigl(\cktRefundRedeem,\; \rid,\; \attSig,\; \redeemView[\rid],\; \Dvault,\;
  \Dshd\bigr)$}.
%\oitem{\wit} $(\usersk,\; \mintNonce)$.
\oitem{\Rel} \upd{$\refundComm(\usersk, \rid) = \redeemView[\rid].\fComm$ \;$\wedge$\;
  the relation of $\cktFailGate$ \;$\wedge$\; the commitment opens to
  $(\freshrn{\mintNonce},\, \colrW,\, \redeemView[\rid].\fShares,\, \secret{\shdAcct_i})$}.
\end{OutList}

\end{TxtBox}

\end{itemize}

    \item \textbf{Transcript from the challenge phase:}

    In the challenge phase, the adversary $\Adv$ selects
    two honest shielded accounts  $\candZero, \candOne$ and an amount $\amt$ to be withdrawn.

    The challenge chooses one of the two users at random ($\candB$), and  executes the withdraw call honestly using that account. Namely:
          $$\Withdraw(\candB, \freshrn{\destAddr^*}, \tokType, \amt, \auxx \,;\, \usersk_b) $$

          where \freshrn{$\destAddr^*$} is a fresh address created on the public chain.


As result of this call, the adversary will observe the following transcript


\begin{CktBox}{trn:withdrawchal}{Transcript for the $\Withdraw$ flow for user $\candB$}
   \upd{$\cktStartWithdraw$}
\begin{OutList}
\oitem{\Dvault} $\sigNonce \mathrel{+}= 1$; \; $\reqMap[\rid] \gets \req$; \;
  \upd{$\wdView[\rid] \gets \bigl(\lnk{fn:refundComm}{\refundComm(\secret{\usersk_b}, \rid)},\,
  \tokType,\, \amt\bigr)$}.
\oitem{\Dshd} \emph{Nullifiers published} (exactly one coin spent.):
  \begin{itemize}[nosep,leftmargin=1.3em,topsep=1pt]
  \item \bluedot $\coinNull\bigl(\secret{\coinFund},\; \mathsf{user}(\secret{\coinsk_b})\bigr)$,
        keyed by the caller's \textbf{secret};
  \item \bluedot $\coinNull\bigl(\freshrn{\coinIn},\; \mathsf{contract}(\vaultAddr)\bigr)$.
  \end{itemize}
  \emph{Commitments published} (one  coin created):
  \begin{itemize}[nosep,leftmargin=1.3em,topsep=1pt]
  \item \bluedot $\coinCommit\bigl(\freshrn{\coinIn},\; \mathsf{contract}(\vaultAddr)\bigr)$ - the wallet
        paying the contract; \textbf{public address again};
  \item \bluedot $\coinCommit\bigl(\coinBurnOut,\; \burnAddr\bigr)$, over
        $\bigl(\evolveNonce(\burnNonce),\, \mathsf{color}^{*},\, \amt\bigr)$ --- the contract
        paying the burn address;
  \item \bluedot $\coinCommit\bigl(\coinChange,\; \mathsf{user}(\secret{\coinpk_b})\bigr)$, of value
        $V - \amt$ - keyed by the caller's public coin key.
  \end{itemize}
 \MidNote{
  \oitem{\stmt} $\bigl(\upd{\cktStartWithdraw},$ $\req,$ $\rid,$ $\upd{\wdView[\rid]},$
 $\Dshd,$
   $\notif(\vaultAddr, \rid, \reqMap),$ $\comIO\bigr)$, where $\req$ exposes
   $\tokType, \amt, \destAddr^*, \evmNonce, \sigNonce$ and $\keyPath = \keyPathLit$.

 \oitem{\Rel} \upd{$\wdView[\rid].\fComm = \refundComm(\secret{\candB}, \rid)$} \;$\wedge$\;
   $\coinIn = (\freshrn{\burnNonce}, \mathsf{color}^{*}, \amt)$ opens both
   $\coinCommit(\freshrn{\coinIn}, \mathsf{contract}(\vaultAddr))$ and
   $\coinNull(\freshrn{\coinIn}, \mathsf{contract}(\vaultAddr))$ \;$\wedge$\;
    $\coinFund_i$ opens its published nullifier under $\coinsk$ and is unspent \;$\wedge$\;
  the color balances: $V + \amt$ in, $\amt + (V-\amt) + \amt$ out.
  }
\end{OutList}
 % \oitem{\wit} $\bigl(\secret{\usersk_b},\; \burnNonce,\; \secret{\coinsk_b},\; \{\coinFund_i\}\bigr)$.
\vspace{5pt}
$\cktCompleteWithdraw$ (See circuit~\ref{ckt:completewithdraw})
\begin{OutList}
\oitem{\Dvault} $\reqMap[\rid]$ removed; \upd{$\wdView[\rid]$ removed}.
\oitem{\Dshd} $\okbit = 1$: $\emptyset$.   $\okbit = 0$: \upd{\bluedot mint record
  $\bigl(\lnk{fn:domSep}{\domSep(\wdView[\rid].\fErc)},\, \wdView[\rid].\fAmt\bigr)$; \;
  \bluedot $\lnk{fn:coinCommit}{\coinCommit(\freshrn{\mintNonce}, \tokType, \amt, \secret{\candB})}$}.
\oitem{\stmt} $\bigl(\cktCompleteWithdraw,\; \rid,\; \attSig,\; \okbit,\;
  \upd{\wdView[\rid]},\; \Dshd\bigr)$.
$(\secret{\candB}, \mintNonce)$.
\oitem{\Rel} $\okbit = 1$: $\EcdsaVf(\attDigest(\rid,\evmOut), \attSig, \mpcRespKey)$. \quad
  $\okbit = 0$: the same, $\wedge$ \upd{$\refundComm(\secret{\candB}, \rid) =
  \wdView[\rid].\fComm$ $\wedge$ commitment opening
  $(\freshrn{\mintNonce}, \cdot, \wdView[\rid].\fAmt, \secret{\candB})$}.

%  \oitem{\wit} [private] $\okbit = 1$: $\bot$. \quad $\okbit = 0$:
\end{OutList}

\vspace{10pt}

\upd{$\cktRefundWithdraw$} (See circuit~\ref{ckt:refundwithdraw})
\begin{OutList}
\oitem{\Dvault}  \upd{$\reqMap[\rid]$ removed; $\wdView[\rid]$ removed}.
\oitem{\Dshd} \bluedot mint record $\bigl(\domSep(\upd{\wdView[\rid].\fErc}),\,
  \upd{\wdView[\rid].\fAmt}\bigr)$; \;
  \bluedot $\lnk{fn:coinCommit}{\coinCommit(\coinMint, \secret{\candB})}$.
\oitem{\stmt} $\bigl(\upd{\cktRefundWithdraw},\; \rid,\; \attSig,\; \upd{\wdView[\rid]},\;
  \Dvault,\; \Dshd\bigr)$.

\oitem{\Rel} $\refundComm(\secret{\candB}, \rid) = \upd{\wdView[\rid].\fComm}$ \;$\wedge$\;
  \upd{the relation of $\cktFailGate$} \;$\wedge$\; the commitment opens to
  $(\mintNonce, \cdot, \upd{\wdView[\rid].\fAmt}, \secret{\candB})$.

%  \oitem{\wit} [private] $(\secret{\candB},\; \mintNonce)$.
\end{OutList}
\end{CktBox}


 \item\textbf{transcript from the oracle calls post-challenge:}
    The transcript here is the same as in the pre-challenge phase under the condition specified in the definition (See Definition~\ref{sec:definition}).
\end{itemize}


Let $p_0$ be the probability that any PPT  $\Adv$ wins the game $\unlinkgame_{\hybrid_{0},\Adv}$, i.e.


$$Pr[\unlinkgame_{\hybrid_{0},\Adv} \rightarrow 1] =p_0$$

\subsubsection{\texorpdfstring{\MidNoteT{Hybrid 1  - Midnight  Shielded Computation  and Shielded Coin Layer and }}{Hybrid 1 - Midnight Shielded Computation and Shielded Coin Layer and}}

In this hybrid experiment, we use the properties offered by the Midnight Layer.
We consider a sequence of sub-hybrid, and in each one of them we replace the secrets used in a Midnight component, and use
the underlying security of the Mindight layer to claim that
our replacement is sound.
While we state precisely which assumptions we are using, we
do not provide formal reductions here, since it requires to prove the security of the Midnight components, which we assume to be secure by hypothesis assumption.


In Hybrid 1  we will use the following cryptographic propoerties that are implicit in the Midnight shielded comptuation layer.

\begin{enumerate}
    \item \textbf{ZK Proofs protect the secrets (Zero-Knowledge Property~\ref{def:nizk})}. The zk proofs associated to the statements  $\stmt$ and relation $\Rel$ satisfy the \emph{zero-knowledge property}. Technically this means that we can replace every proof  generated using the user's secrets: $\usersk_i$, $\coinsk_i$, $\mintNonce$, $\shdAcct_i$ etc, with a simulated proof  that is calculated without any secret.

    \item \textbf{Commitments protect the receivers 's of the coins (Hiding~\ref{def:hiding}).} In circuits
    $\cktCompleteDeposit$, $\cktCompleteWithdraw$, $\cktCompleteSwap$ coins are calculated as the output of
    \MidNoteT{$\opMint$},   using the recipient's identity ($\shdAcct_i$ 's public key coin $\coinpk_i$):
    \begin{align*}
    {\coinCommit(c, \mathsf{rcpt})} &= \PersistantHash\bigl(\text{cc-tag} \concat \mathsf{nonce} \concat
  \mathsf{color} \concat \mathsf{value} \concat \coinpk_i\bigr)
\end{align*}

  Assuming that the above  Midnight Commitment functions satisfy hiding when
the  $\mathsf{nonce}$ is picked at random, we can replace the recipient identity $\coinpk_i$ with the string $0$
and no adversary can detect the change.


\item\textbf{Nullifers do not leak information about the recipients (PRF~\ref{def:prf}).} Nullifiers are computed
as part of {\MidNoteT{$\opReceive$}} code,  when the user sends a coin to the contract (in $\Withdraw$ and $\Swap$ flow).
They are computed as follows:
\begin{align*}
{\coinNull(c, \mathsf{sndr})} &= \PersistantHash\bigl(\text{cn-tag} \concat \mathsf{nonce} \concat
  \mathsf{color} \concat \mathsf{value} \concat \mathsf{sndr}\bigr)
\end{align*}

Assuming that the hash function $\PersistantHash$ behave as a PRF and assuming that $\mathsf{nonce}$ remains private throughout, the nullifer protects the user $\mathsf{sndr} (\coinsk_i)$.

\item  \textbf{Additional ZSwap layer artifacts do not leak any information about the sender and receivers of a shielded transaction. }
A ZSwap transaction contains several cryptographic layers that provides the correctness and functionality guarantees for a shielded payment to execute and land correctly.
For instance, as part of the soundness of the payments,
the value involved in the coins are also committed
in Pedersen's commitments whose correctness is proved in zk via specifically designed Schnorr's zk proofs (Definition~\ref{def:hvzk}) (see~\cite{Zswap22}).
Finally, the shielded transaction contains a ciphertext, that allows the recipient to detect that the coin generated in the transaction is for them. This ciphertext is computed with a \emph{key-private}  public key encryption scheme (See Definition~\ref{def:keyprivate}. Key private means that the ciphertext does not leak any information about the public key that was used to compute it. This property is necessary for unlinkability.

\end{enumerate}


Using the above  properties,  in Hybrid 1, we make the following changes in the way the experiment  $\unlinkgame_{\Pi,\Adv}$.

\begin{enumerate}
    \item  \textbf{Hybrid $\hybrid_{1.1}$} \textbf{Replace hones tproof with  Simulated ZK proofs (use no witness)}:  In $\unlinkgame_{\hybrid_{1.1},\Adv}$ we replace all the zk proofs  created in the system (zk proofs of the circuits, and zk proofs of the SZSwap wallets),  with \emph{simulated} zk  proofs  $\zkSim_2(\crs,\td,\stmt)$  (with  $(\crs,\td) \gets \zkSim_1(1^{\secpar},\Rel)$ ) that are calculated  \emph{without any secret.}
    as per Definition~\ref{def:nizk}).
     This means that the transcript of every zk proof, \emph{does not carry any information about the identity of the caller} of the circuits executed on Midnight.

     If $\Adv$ wins the game $\unlinkgame_{\hybrid_{1.1},\Adv}$ with a probability $p_{1.1}$ that is non-negligible different than $\unlinkgame_{\hybrid_0,\Adv}$,  then $\Adv$ can be used to break the zk property of each proof.   Therefore, due the zero-knowledge Property (\S~\ref{def:nizk}), it follows:

    $$Pr[\unlinkgame_{\hybrid_{1.1},\Adv} \rightarrow 1] = p_{1.1}\approx_{\negl(\lambda)} p_0$$

    \item \textbf{Hybrid $\hybrid_{1.2}$} \textbf{Remove Recipient Identity from Coin Commitments}: in this hybrid game, the challenger replaces all the commitments representing the coins that follow the form:
\begin{align*}
{\coinCommit(c, \mathsf{rcpt})} &= \Hash\bigl(\text{cc-tag} \concat \mathsf{nonce} \concat
  \mathsf{color} \concat \mathsf{value} \concat \mathsf{rcpt}\bigr)
\end{align*}

which takes as input the identity of the recipient  $\mathsf{rcpt}$, with  the following:

\begin{align*}
{\coinCommit(\mathsf{nonce} \concat
  \mathsf{color} \concat \mathsf{value} \concat  0)}
\end{align*}
Note, in this experiment, the challenger will still maintains the private list of the raw coin-info (nonce, color, value) for each unspent coin. Such information, however,  is not inputted  in the commitment function.

     If $\Adv$ wins the game $\unlinkgame_{\hybrid_{1.2},\Adv}$ with a probability $p_{1.2}$ that is non-negligible different than $\unlinkgame_{\hybrid_{1.1},\Adv}$,  then $\Adv$ can be used to  break the hiding of the commitment scheme used by Midnight shielded-coin layer.  The reduction to hiding can go through since the proofs are simulated and since our protocol invokes the Midnight library  with secret and random nonces for each new coin minted (evolved nonces are generated by Midnight).  Therefore,  assuming that $\coinCommit$ satisfies computationally hiding property (Definition~\ref{def:hiding}, with $\csGen(1^{\secpar})$ outputing the parameters for the hash function) it follows that: $$Pr[\unlinkgame_{\hybrid_{1.2},\Adv} \rightarrow 1]=  p_{1.2}\approx p_{1.1}$$


 \item \textbf{Hybrid $\hybrid_{1.3}$} Nullifier are computed without any sender's information ($\coinsk_i$).

 Namely, in   $\hybrid_{1.1}$ the challenge replaces this computation:

\begin{align*}
{\coinNull(c, \mathsf{sndr})} &= \PersistantHash\bigl(\text{cn-tag} \concat \mathsf{nonce} \concat
  \mathsf{color} \concat \mathsf{value} \concat \mathsf{sndr}\bigr)
\end{align*}

with  this:
\begin{align*}
{\coinNull(c, \mathsf{sndr})} &= \PersistantHash\bigl(\text{cn-tag} \concat \mathsf{nonce}  \concat \mathsf{value}\concat  0 \bigr)
\end{align*}

This computation is exactly as the coin commitment with the only difference in the tag   \textsf{cn-tag} (vs \textsf{cc-tag}). Here $\mathsf{nonce}$  is the randomness used to compute the coin commitment is some earlier invocation of $\opMint$.

 If $\Adv$ wins the game $\unlinkgame_{\hybrid_{1.3},\Adv}$ with a probability $p_{1.3}$ that is non-negligible different than $\unlinkgame_{\hybrid_{1.2},\Adv}$,  then $\Adv$ can be used to break the security of Midnight's nullifers.
 This is part of the Midnight design and we defer the reader to their proofs of security.

Therefore, $$Pr[\unlinkgame_{\hybrid_{1.3},\Adv}  \rightarrow 1]=  p_{1.3}\approx_{\negl(\lambda)}  p_{1.2}$$
\end{enumerate}


 \item \textbf{Hybrid $\hybrid_{1.4}$ \MidNote{ZSwap Wallet Layer (Pedersen Commitments and Key-Private Public key encryption.)} }
A shielded payment involves the use of other cryptographic objects  such as Pedersen's Commitments (Def.~\ref{def:commit}), Key-Private Public key Ciphertexts (Def.~\ref{def:keyprivate}), and Schnorr's proofs (Def.~\ref{def:hvzk}), as described in~\cite{Zswap22}.
Pedersen's Commitments  and Schnorr's proofs are used as part of the Midnight shielded-coin layer for the soundness of the proof, whereas Key-Private Public key Ciphertexts are used to
guarantee discoverability of the coin  to the receiver.
In this hybrid experiments, we  remove the secrets used in any of such objects. We can do so, since all zk proofs are simulated.
Our code does not interfere with this layer, neither provides the randomness used in the wallet itself, so we inherit the security properties guaranteed by the Midnight shielded coin layer.
The  resulting transcript should  be indistinguishable from the transcript generated in  $\hybrid_{1.3}$,  due to the zk property of Schnorr's zk proofs and  the Key-Privacy of the  Public key Ciphertexts.


Therefore, $$Pr[\unlinkgame_{\hybrid_{1.4},\Adv}  \rightarrow 1] = p_{1.3}\approx_{\negl(\lambda)} p_{1.3}$$


\subsubsection{Hybrid 2. Replace \texorpdfstring{$\Hash(\usersk_i, \cdot)$}{HashK(userSk\_i, .)} with a Random Function \texorpdfstring{$\RandFun_i$}{RF\_i} (uses no secret)}
In hybrid $\hybrid_2$ we change the calculation of $\userComm$ and $\refundComm$ that are directly  dependent on the user secret $\usersk_i$.
We replace the function $\Hash_{\usersk_i}$ keyed with the secret key of the user $\usersk_i$, with a random function $\RandFun_i$. Precisely,  the following computations (See section~\ref{ssec:userskeys}):
\begin{align*}
{\userComm(\usersk)_i} &= \Hash\bigl(\text{``vault:user:''}\bc \usersk_i\bigr)
\\
{\refundComm(\usersk_i, \rid)} &= \Hash\bigl(\text{``vault:refund:''}\bc \usersk_i\bc \rid\bigr)
  &&\text{\small the gate on refunds and swap settlement}
\end{align*}

are replaced with:


\begin{align*}
 {\userComm^{\RandFun_i}} &= \RandFun_i\bigl(\text{``vault:user:''}\bigr)\\
{\refundComm^{\RandFun_i}(\rid)} &= \RandFun_i\bigl(\text{``vault:refund:''}\bc  \rid\bigr)
 \end{align*}


The only difference between the hybrid  $\unlinkgame_{\hybrid_{1.4},\Adv}$  and $\unlinkgame_{\hybrid_{2},\Adv}$
is in the use of truly random functions
$\RandFun_1(\cdot)$, $\RandFun_2(\cdot), \dots, \RandFun_{\numusers}(\cdot)$ instead of
$\Hash_{\usersk_1}(\cdot)$, $\Hash_{\usersk_2}(\cdot)$, $\ldots,\Hash_{\usersk_{\numusers}}(\cdot)$.

Assume that $\Adv$ wins the game $\unlinkgame_{\hybrid_{2},\Adv}$  (when $\Hash$  was used ) with a probability $p_{2}$,  that is non-negligible different from $\unlinkgame_{\hybrid_{1.4},\Adv}$ (where independent random functions $\RandFun_i$ are used), which is $p_{1.4}$.
Then we can use $\Adv$ to build a  distinguisher $\distprf$ that violates the multi-user PRF security of $\Hash$  (Definition~\ref{def:muprf}) with  the same probability.

\begin{tcolorbox}[
  breakable,
  colback=boxbg,
  colframe=boxrule,
  arc=2pt,                 % Corner radius
  boxrule=0.8pt,           % Border thickness
  top=8pt, bottom=8pt,     % Inner padding
  left=10pt, right=10pt
]
\underline{\textbf{Reduction $\distprf^{\oracle(\cdot,\cdot)}$}}\\
Activate adversary $\Adv$ and simulated game $\unlinkgame_{\hybrid_{1.4},\Adv}$ to it as follows:
\begin{itemize}
    \item Simulate the execution of every oracle using the same approach as in Experiment $\hybrid_{1.4}$ with the following exception:
        \begin{enumerate}
            \item For every new user $U_i$ created for the cross-chain system, instead of picking the application secret key $\usersk_i$, the challenger calls the oracle on input $i$: $\oracle(i,\cdot)$
            \item For every deposit flow for user $U_i$ compute function $\userComm$ as:
            \begin{align*}
                & \userComm^{i} = \oracle(i, \text{``vault:user:''})\\
            \end{align*}
            \item For every withdraw,  Swap,  Supply, Redeem flow for user $U_i$ compute function as
            \begin{align*}
                &\refundComm^{i}(\rid) = \oracle(i,\text{``vault:refund:''}|| \rid)\\
            \end{align*}
        \end{enumerate}
    \item Challenge \[(\candZero\bc \candOne\bc \tokType\bc \amt\bc \auxx^*\bc \outc^*\bc \mathsf{st}) \gets \Adv^{\mathcal{O}}(\pp),\]
      where $\candZero, \candOne$ are honest shielded accounts.

    \textbf{Challenge}
    \begin{enumerate}
        \item If $\candZero$ or $\candOne$ is not honest, output $0$;
        \item if $\bal[\candZero][\tokType] < \amt$ or $\bal[\candOne][\tokType] < \amt$, output $0$;
    %    \item if $\Leak(\candZero, \tokType, \amt) \neq \Leak(\candOne, \tokType, \amt)$, output $0$.
        \item otherwise, pick a random bit $b$ and execute:
        \[ \Withdraw(\candB, \destAddr^*, \tokType, \amt, \auxx \,;\, \privIn_b) \]
        for a fresh public-chain address $\destAddr^*$, replacing the function $\Hash$ for $\candB$ with the evaluation of its correspondent oracle.
\end{enumerate}

\item Post-Challenge: simulate all the cross-chain oracles as above.
\item Output: upon receiving a bit $b'$ from $\Adv$.
\item If $b'=b$ output $1$, else output $0$.
\end{itemize}
\end{tcolorbox}


\vspace{5pt}
\noindent
\textbf{Analysis of $\distprf$'s probability of success:}
\begin{itemize}
    \item Case $\oracle(i,\cdot)= \RandFun_i(\cdot)$.

    If the Oracle is implemented as a collection of truly random functions, then  $\distprf$ is simulating the experiment $\unlinkgame_{\hybrid_{2},\Adv}$.
    Hence,


$$Pr[\distprf^{\RandFun_i(\cdot)}\rightarrow 1 ] = Pr[\unlinkgame_{\hybrid_{2},\Adv}\rightarrow 1 ]= p_{2}$$


     \item Case $\oracle(i,\cdot)= \Hash(sk_i,\cdot)$
      If the Oracle is implemented as the function $\Hash$ evaluated on  independently sampled secret keys $\usersk_i$, then the transcript generated by $\distprf$ is identical to $\unlinkgame_{\hybrid_{1.4},\Adv}$. Hence,


     $$Pr[\distprf^{\Hash(sk_i, \cdot)}\rightarrow 1] = Pr[\unlinkgame_{\hybrid_{1.4},\Adv}\rightarrow 1 ]= p_{1.4}$$

\end{itemize}


Hence:

$$ \left| Pr[\distprf^{\RandFun_i(\cdot)}\rightarrow 1] - Pr[\distprf^{\Hash(sk_i, \cdot)}\rightarrow 1] \right|  = p_{2}- p_{1.4}  $$

Since we assume that $\Hash$ satisfies the property defined in Definition~\ref{def:muprf}, it follows that $p_{2}- p_{1.4}$ must  be negligible.


\vspace{10pt}
We have removed the secret key of the caller from the authentication tags $\userComm$ and $\refundComm$.
In this hybrid world, the application secret key  $\usersk_i$ is not used in \emph{any computation}.


Note, however, that even  if we are using a random function $\RandFun_i$, the output  is still deterministic -- i.e., when queried with the same input, $\RandFun_i$ will return the same output.
So, for instance, since in Deposit,  $\userComm_i$ is computed on the same string ``vaul:user' every time, the adversary can easily detect that the same user is performing various deposits. This is by design: we are not protecting the linkability of the deposit.

We are interested in guaranteed that  withdrawals/swaps/supply/redeem  are unlinkable to each other and unlinkable to the deposits. Recall that for these action, the user authenticate with $\refundComm_i$.
And recall that:

$${\refundComm^{\RandFun_i}(\rid)} = \RandFun_i\bigl(\text{``vault:refund:''}\bc  \rid\bigr)$$


 where $\rid$ is unique per request (and $\rid = \Kec(\req)$).

This means that every $\refundComm_j$ is  different from any other $\refundComm_j$ for any $i,j$, and carries no information about the recipient.
Finally the request itself: $\req$
\begin{enumerate}
    \item $\req$ does not contain information about the identity of the withdrawer/swapper/supplier/redeemer.
    \item each $\req$ is assumed to be  unique due to the correctness of the underlying  EVM nounce.
\end{enumerate}


Putting everything together, in $\hybrid$ 2 we have:

\begin{enumerate}[nosep,leftmargin=1.7em,topsep=3pt]
\item[(i)]  no identity is leaked from the cryptographic components. In the previous hybrid we removed the uses of the secret inputs from the coin and circuit computation, and in this hybrid we removed the identity $\usersk_i$.
So, in hybrid $\hybrid_2$ we are left with  no secret and no account identity is used in producing any component of the view: proofs are simulated, coin commitments carry a dummy recipient, nullifiers carry a  dummy sender, and $\userComm$ and $\refundComm$ are evaluations of a random function;

\item[(ii)] no information identity is  leaked from the withdraw/swap/supply/redeem request $\req$ (this is by design);
\item[(ii)] no information is leaked in the spending  due to our consolidate-coin assumption (\S~\ref{sub:nofingerprint}).
\end{enumerate}


This means that in  the world $\hybrid$ 2, the protocol messages    are computed without any user-secret,  the public transcript contain no information that identify the user who is chosen in the challenge phase by the challenger to execute $\Withdraw$
Consequently, in this world, no adversary can guess which user was chosen  with a probability better than  $\frac{1}{2}$.
Hence:  $$\Pr[\unlinkgame_{\hybrid_2,\Adv} \rightarrow 1] = \frac{1}{2}$$


and due to the sequence of indistinguishable hybrid arguments,  we conclude:

$$\Pr[\unlinkgame_{\hybrid_0,\Adv} \rightarrow 1] = \frac{1}{2}+\negl(\secpar)$$


\qed
\end{proof}


\subsection{Swap/Supply/ Redeem}
\label{sub:swapsupplyredeemproof}


The proof of Withdraw Unlinkability develops a sequence of arguments
that removes  \emph{all user-dependent secrets $\usersk_i$ and $\shdAcct_i$ from every computation.}
The final hybrid experiment of the proof is a \emph{simulated} world where no secrets are used (in any of the procedures).
We notice that that Swap, Supply and Redeem follow the same structure of the Withdraw, and the same sequence of arguments apply.
In the following we examine each call and explain why the same proof argument applies.


\medskip
\noindent\textbf{$\Swap$.}  We notice that the circuits  $\cktStartSwap$/ $\cktCompleteSwap$
 follows the same structure  of  $\cktStartWithdraw$/ $\cktCompleteWithdraw$.
Specifically, $\cktStartSwap$ is the same as $\cktWithdraw$ with $\swapMap$ and
$\swapMarkerMap$ in place of $\reqMap$ and $\wdMarkerMap$. The logic is also
almost identical: a coin is surrendered and burned in the same two moves, the same authentication tag $\refundComm(\usersk_b, \rid)$  is computed.
The only difference is that in swap, the transaction on the public chain involves only the vault account and concerns swapping coin \emph{types}, whereas in the withdraw request, the transaction must have an explicit evm destination address $\destAddr$.
The settlement call   $\cktCompleteSwap$
differs in two ways from $\cktCompleteWithdraw$, however, neither of them depends on the user's identity.
$\cktCompleteWithdraw$  mints on the \emph{success} path rather than the failure path, and it mints
\emph{two} coins rather than one.
Minting more coins does not change our arguments in  $\hybrid_{1.1}$, since in that hybrid we replace  every proof in the transcript and $\hybrid_{1.2}$ over \emph{every coin commitment}, so a second commitment is replaced by the same hybrid and the same reduction.
The amounts released and swapped, are in the public state, just like the amount withdrawn.
 Every step of $\hybrid_0$ through $\hybrid_2$ therefore applies verbatim, with
$\lendunlinkgame$ in place of $\unlinkgame$, and Definition~\ref{def:lendunlink}.

\medskip
\noindent\textbf{$\Supply$.}  $\Supply$ is a mirror of $\Swap$ with the following differences:
(1) in $\Supply$ the token types and color are hardwired in the call (whereas in $\Swap$ they are an input chosen by the caller);  since types and color are public in all the calls, this change does not affect the analysis.
(2) $\Supply$ returns $\amtShares$ shares instead of an swapped amount.
$\amtShares$ is not chosen by the challenger, it is determined by $\amt$ and by the market's exchange rate at the
executing block, and the rate is a property of $\Lpub$, global to all suppliers, not of the shielded
account that paid.
Both differences do not affect the creation of the cryptographic objects and do not affect the analysis.

\medskip
\noindent\textbf{$\Redeem$.}  $\Redeem$ is  identical to $\Supply$ but with the color swapped, with the  released $\amtAssets$ being a function of $\amtShares$ and the same global rate.


\remove{
We formalize this in the following Lemma~\ref{lem:transfer}
\begin{lemma}[The chain of hybrid is  ``Procedure Independent'']
\label{lem:transfer}
Let $\mathsf{P}$ in ($\Swap, \Supply, \Redeem$)  be a procedure of $\Pi$ , and let
$\mathsf{Exp}_{\mathsf{P}}$ be the experiment of Fig.~\ref{fig:unlinkGeneral} with $\Withdraw$ replaced by $\mathsf{P}$. If:
\begin{enumerate}[leftmargin=2.2em,itemsep=2pt,topsep=3pt]
\item[\textbf{(C1)}] the request call writes $\keyPath = \keyPathLit$, so that every field of the stored record $\req$ is a function of the transaction arguments and of public state, and \emph{no field of it depends on the caller};
\item[\textbf{(C2)}] the only value the call publishes that depends on the caller's long-term secret is $\refundComm(\usersk, \rid)$, stored in a marker map whose identity is determined by the product and not by the caller;
\item[\textbf{(C3)}] every value released by the settlement call is a function of $\req$ and of the  public state of $\Lpub$, and the coin-layer components of the view consist of coin commitments  and nullifiers in numbers fixed  (by assumption on the spending fingerprint (Section ~\ref{sub:nofingerprint}))
\end{enumerate}
Then $\mathsf{P}$ satisfies unlinkability under the same assumptions as Theorem~1, with the same chain of hybrids and the same reductions.
\end{lemma}

\begin{proof}[Proof sketch]
In the proof \S\ref{sec:proof} each hybrid  is a substitution applied to a class of published objects and does not depend on which procedure produce them.
For instance,  $\hybrid_{1.1}$ replaces every proof by a simulated
one -- across all procedures that invoke zk proofs--;  $\hybrid_{1.2}$ replaces the recipient slot of \emph{every coin commitment} with $0$, $\hybrid_{1.3}$ replaces the sender
slot of every nullifier with $0$, $\hybrid_{1.4}$ replaces the identity from every ZSwap wallet objects (Pedersen's commitment, Schorr's proof  and Ciphertexts),
and finally  $\hybrid_2$  replaces every evaluation of $\Hash_{\usersk_i}$ with a random function $\RandFun_i$.
The only object that are specific to each procedure, are not the public parameters of the request and the public consequence of each request and do not depend on $b$. Conditions
(C1)--(C3) are exactly this claim. (C1) removes the request record, (C2) confines the caller's secret
to a single value of the form the hybrids already retire, and (C3) removes the settlement release and
fixes the coin-layer cardinalities. Clause (ii) of $\hybrid_3$ is discharged by  the no-spending fingerprint assumption \S~\ref{sub:nofingerprint}.
\end{proof}
}

\section{Soundness}
\label{sec:soundness}

In this section we analyze the soundness of the system.
When defining soundness, we identify the soundness requirement for each action, and then defined the soundness of the system as 
the conjunction of the soundness of each properties.
An alternative approach would have been to define an ideal functionality for our cross-chain system payment, and proving UC security of it, and is intended for
future analysis.

This section is organized as follows: (1) Definition (Section~\ref{sec:soundnessdefinition}): we define soundness of the system; (2) Assumptions (Section~\ref{sec:soundnessassumptions}) we state the trust, model and cryptographic assumption we rely on to argue soundness (3) Proof (\S~\ref{subsec:soudnesscommon}-\S~\ref{susec:soundnessrest}) we prove that each action achieves soundness under the stated assumptions.


\subsection{Definition of Soundness} 
\label{sec:soundnessdefinition}
For simplicity,  in this first formal treatment, we define soundness  of each procedure : Deposit, Withdraw, Swap, Supply, Redeem individually. 
The soundness of the overall system is then defined as the soundness of all of its procedures.
In the future, we will consider using the Universal Composability framework to have a more comprehensive approach.


\begin{definition}
\label{def:depositss}
\emph{Deposit Soundness:} $\Pi.\Deposit^{\Lpub, \Lshd}\bigl(\pubAddr\bc \shdAcct\bc \tokType\bc \amt\bc
      \auxx \,;\, \privIn\bigr)$ is \emph{sound} if no PPT algorithm $\Adv$ can  deposit  $\amt$, $\tokType$ on its own account $\shdAcct*$ on $\Lshd$  and have transferred $\leq v$ of the same token to the vault on $\Lpub$ from its own account $\pubAddr$ with non-negligible probability.
 \end{definition}


\begin{definition}
\label{def:withdrawss}
    
\emph{Withdraw Soundness:} $\Pi.\Withdraw^{\Lpub, \Lshd}\bigl(\cdot\bc \destAddr*\bc \tokType\bc \amt\bc
      \auxx \,;\, \privIn\bigr)$ is \emph{sound} if 
      two conditions are satisfied:
      (1) eligibility: no PPT algorithm $\Adv$ can  withdraw  $\amt$, $\tokType$ without controlling 
      an account $\shdAcct^*$  on Midnight  that  holds at least value ($\amt$, $\tokType$) with non-negligible probability.
      (2) no double-spending: no PPT algorithm $\Adv$, owning a shielded account $\shdAcct^*$ with amount $v^*$  can successfully withdraw  $\amt$ of $\tokType$ from $\shdAcct^*$, and at its completion own 
      an amount $\geq v^*- \amt$.
      
    
% \end{definition}
\end{definition}

\begin{definition}
\label{def:swap} \emph{Swap Soundness:} 
$\Pi.\Swap^{\Lpub, \Lshd}\bigl(\cdot\bc \tokType\bc \amt\bc
    \tokTypeB\bc \amtB\bc \auxx \,;\, \privIn\bigr)$ is \emph{sound} 
    if   two conditions are satisfied:
      (1) eligibility: no PPT algorithm $\Adv$ can  swap $\amt$ of $\tokType$ unless $\Adv$ controls an account $\shdAcct^*$ that owns at least $\amt$ of $\tokType$, with more than non-negligible probability. 
      (2) no double-spending: no PPT algorithm $\Adv$,   owning a shielded account $\shdAcct^*$ with amount $v^*_A$  of 
      $\tokType_A$ and  $v^*_B$ of $\tokTypeB_B$ can successfully swap  $\amt$ of $\tokType_A$  for $\tokType_B$ and then, upon swap completion own $\geq v^*_A - \amt$  of $\tokType_A$, with non-negligible probability. 
\end{definition}

\begin{definition}
\label{def:supplysound}
\emph{Supply soundness}
$\Pi.\Supply^{\Lpub, \Lshd}\bigl(\shdAcct^*\bc \tokType\bc \amt\bc \auxx;\, \privIn\bigr)$
is \emph{sound}  if   two conditions are satisfied: (1) eligibility: no PPT algorithm $\Adv$ can  supply $\amt$ of the underlying
color $\colrU$ (a function of $\tokType$) without controlling an account $\shdAcct^*$ on Midnight that owns at least $\amt$ of
$\colrU$, with more than non-negligible probability.  (2) no double-spending: no PPT algorithm $\Adv$,   owing a shielded account $\shdAcct^*$ with amount $v^*_U$  of   $\colrU$ can successfully supply  $\amt$ of $\colrU$ and, upon supply completion own $\geq v^*_U-\amt$ of $\colrU$ tokens,  with non-negligible probability. 
\end{definition}

\begin{definition}
\label{def:redeemsound}
\emph{Redeem soundness.}
$\Pi.\Redeem^{\Lpub, \Lshd}\bigl(\shdAcct\bc \tokType'\bc \amtShares\bc \auxx
      \,;\, \privIn\bigr)$ is \emph{sound} if no PPT $\Adv$ can 
      redeem $\amtShares$ of coins of $\colrW$ (a function of $\tokTypeB$ without having supplied an $\amt$ of USD  through $\Supply$, such that 
      $\amtShares$ = $\amt+ \mathsf{interest}$ with non-negligible probability.
      
 \end{definition}

\begin{definition}[\textbf{Cross-Chain \Payment System  Soundness}]
\label{def:totalsoundness}
    A \emph{public-to-private cross-chain \Payment system} $\Pi$, relative to $\Lpub$ and $\Lshd$ is  \emph{sound} it achieves  Deposit, Withdraw, Swap, Supply and Redeem soundness.
\end{definition}
\subsection{Assumptions}
\label{sec:soundnessassumptions}
In this proof we use the following assumptions

\vspace{5pt}
\noindent
\textbf{System Assumptions:}
\begin{itemize}
    \item Ethereum  and Midnight are secure and sound  instantiations of $\Lpub$ and $\Lshd$, respectively.

    \item MPC $\SigSys$ is trusted. It satisfies:
    \begin{enumerate}
        \item  \textbf{MPC committee integrity:}  the committee signs only what the protocol authorises. 
        \item \textbf{Attestation faithfulness:}  MPC only provide attestations for events that have been executed on $\Lpub$ (Ethereum).
        $\SigSys$ attests $\failed$ \emph{only} for a request that did \emph{not} execute on $\Lpub$.
     
    \end{enumerate}
\end{itemize}

\vspace{5pt}
\noindent
\textbf{Network Assumptions:}  Communication Channels are authenticated.


\vspace{5pt}
\noindent
\textbf{Cryptographic Assumptions:}
\begin{itemize}
    \item \textbf{ECDSA is EUF-CMA under additive key derivation and pre-signatures}
    (Definition~\ref{def:sig}). In our proof we use this property as a guarantee that no adversary forges the MPC
    attestation signature $\attSig$ that the settlement circuits verify against the pinned
    $\mpcRespKey$. This assumption is also implictly used under the system assumption that  Ethereum  ($\Lpub$)  is sound,  so, no adversary forges a transaction signature under the vault's
    derived EVM account, which would move pooled funds with no Midnight call at all. Additive
    key derivation must be inside the game because the vault's account is a deterministic
    function of $(\vaultAddr, \keyPath)$; pre-signatures must be inside it because the MPC
     produces signatures using a pre-processing step.

    \item \textbf{$\Kec$ is collision resistant} (Definition~\ref{def:crhf}). We use this property for the computation of 
    $\rid = \Kec(\req)$,  to claim that an attestation on $\rid$ for $\req$ cannot be re-used to settle a different request $\req'$ that collides with  $\rid$.
    We also use in the computation, $\attDigest = \Kec(\rid \concat \evmOut)$, to avoid that a collision lets a signature issued for one attested output be replayed as a signature for another.


    \item \textbf{$\Hash$ is a secure PRF} (Definition~\ref{def:prf}) \textbf{and therefore a MAC
    on fixed-length inputs} (Definition~\ref{def:mac}). Soundness uses the MAC direction: the
    caller gates accept only a witness $\usersk$ satisfying
    $\refundComm(\usersk, \rid) = \swapMarkerMap[\rid]$, so an adversary not holding $\usersk$
    cannot claim another user's refund, and the same assumption protects the depositor gate
    $\userComm(\usersk)$.

    \item \textbf{Midnight's shielded-circuit layer is sound} (Definition~\ref{def:zksound} and Definition~\ref{def:zkksound}).
    No adversary produces an accepting $\zkproof$ for a transcript $\stmt$ that no honest
    execution of the circuit could have produced. This is what makes every assertion in the
    circuit boxes binding on the published state update, and what stops any authentication step to be skipped.

    \item \textbf{Midnight's shielded-coin  and circuit-layer  achieves soundness and  zero-knowledge}
    (Definitions~\ref{def:zksound} and~\ref{def:binding}). Soundness of the Zswap proofs gives
    well-formedness and balance of the offer, so no transaction creates value except through a
    declared mint. Binding of $\coinCommit$ gives that a published commitment opens to at most
    one $(\mathsf{nonce}, \mathsf{color}, \mathsf{value})$, so a surrendered coin cannot be
    re-opened at a larger value and a minted coin cannot be claimed at an amount other than the
    one the circuit fixed. Zero-knowledge  allows us to claim that honest users' secret keys 
    used as witness in computing the proof are protected.

% \item \textbf{$\Hash$ is collision-resistant.} We use  two properties implied by 
% collision-resistance.
% {\bf One-Wayness}: we use the fact that $\gamma = \userComm(\usersk)$ or
%     $\gamma = \refundComm(\usersk, \rid)$ published by an honest user, no adversary finds the pre-image
%     $\usersk^*$.


    \item \textbf{$\Hash$ and $\Kec$  are collision resistant}
    (Definition~\ref{def:crhf}). 
    This property is vital in all authentication gates.
    For instance, an adversary who can find collisions in the digest of the requestID signed by $\SigSys$, can easily re-use valid signatures, on requests that were never executed on $\Lpub$.
    
Note that collision resistance implies one-wayness. We rely on the one-way property of $\Hash$ to prove that an authentication tag ($\gamma=\userComm(\usersk)$ or $\gamma=\refundComm(\usersk,\rid)$ published by an honest user cannot be used by an adversary to pass the authentication gate of the \code{complete} circuits. 
    
    % Used for
    % $\domSep(\tokType)$ = $\Kec(\text{``erc20:vault:''} $ $\concat$ $\tokType)$: two ERC20 addresses that
    % collided would receive the \emph{same} shielded color, so a coin of the cheap token would pass
    % the color check of the expensive one. 
\end{itemize}

\subsection{Proof}
\label{sound:proof}

\begin{theorem}
\emph{If ECDSA is   EUF-CMA under additive key derivation and pre-signatures signature scheme (Definition~\ref{def:sig}), 
$\Kec$ is a collision-resistant hash function (Definition~\ref{def:crhf}), $\Hash$ is a secure PRF (Definition~\ref{def:prf}) and collision-resistant hash function, Midnight shielded-circuit layer achieves soundness (Definition~\ref{def:zksound}) , Midnight shielded-coin layer (proofs and coin-commitments) satisfies soundness (Definition~\ref{def:zksound}) and binding  (Definition~\ref{def:binding}),  then protocol $\Pi$ described in Section~\ref{sec:instantiation} satisfy Soundness (Definition~\ref{def:totalsoundness}.}
\end{theorem}


The proof of this theorem requires proving  soundness of each action, from $\Deposit$\S~\ref{def:depositss} to  Supply\S~\ref{def:redeemsound}.

The main observation is that every action $\Deposit$, $\Withdraw$, $\Swap$, $\Supply$, $\Redeem$  share the same logic and implementation, with minor differences that reflect only in the color  of shielded transaction executed on $\Lshd$.

Every procedure follows the same template:

\begin{enumerate}
    \item \textbf{Start Action:} Every action $X$ is started with the  $\mathtt{startX}$ circuit: $\cktStartDeposit,\cktWithdraw,$ etc.  Starting an action entails in the following steps 
    \begin{itemize}[label=\small$\blacksquare$]
        \item \textbf{Create EVM request}  $\req$ and $\rid$: circuit  crafts the transaction/call that must be executed on $\Lpub$. $\rid=\reqIdOf(\req)$ is created.

    \item \newversion{In the latest release, $\rid$ is not created in the Start Action and
    there is no request counter $\sigNonce$. The Start Action queues the request, the
    permissionless circuit $\cktFlush$ assigns its EVM nonce (the vault account's next nonce
    $\vaultNonce$ for every action but Deposit, whose caller names the nonce of their own
    derived account), and the permissionless circuit $\cktSend$ computes
    $\rid=\reqIdOf(\req)$ over the completed request and notifies the MPC. Distinctness of
    $\rid$ now follows from the EVM nonce, which the vault account never reuses, and a second
    identical request is refused while the first is open.}
        
        \item \textbf{Create authentication tag}: circuit creates $\userComm(\usersk_i)$ (for $\cktStartDeposit$) or $\refundComm(\usersk_i,\rid)$  for all other actions, is created.
         \begin{itemize} [label=\small$\diamond$]
        \item \textbf{* Private Coins on $\Lshd$ are burned}  (* Not present in $\cktStartDeposit$). In a          Widthdraw/Swap/Supply/Redeem request,  circuit privately sends their private coins to a burn       address through the calls $\mnhl{\code{receiveShielded}}$  and $\mnhl{\code{sendImmediateShielded}}$ 
            \end{itemize}
        \item \textbf{$\rid,\refundComm$ added to the map of pending requests for that action} At this point a pending action is identified by the pair $(\rid,\refundComm)$  ( or  $(\rid,\userComm)$ only for $\Deposit$.
    \end{itemize}

    \item \textbf{Complete Action:} Every action $X$ is completed with the  $\mathtt{completeX}$ circuit.
    To complete an action identified by $\rid,\refundComm$, the user must possess a valid attestation
    on $\rid$ and outcome $\evmOut$ that verifies under  $\SigSys$'s public key $\mpcRespKey$.
    The outcome $\evmOut$ can be $\succeeded$ or $\failed$.
    All complete circuits follows these steps:
    \begin{itemize}[label=\small$\blacksquare$]
        \item \textbf{Verify Attestation (i.e., ECDSA signature).} The circuit runs the ECDSA verification algorithm for attestation $\attSig$ for input $(\rid,\evmOut)$ under public key $\mpcRespKey$.

    For this check we use the unforgeability of ECDSA (in the pre-signature and additive-key derivation  setting) and the collision-resistance properties of $\Kec$.
        \item \textbf{Check that $\rid$ is pending.} The circuit check that $\rid$ is pending in the action-dependent map (e.g.,  $\depositView$, $\wdView$, etc).
        % \item \newversion{In the new version we use $\rid$ and  the timing information $\lastseen$ to determine if a request is pending or resolved.}

        \item \textbf{Perform Action: if transaction was successful $\evmOut$= $\succeeded$}
        \begin{itemize}[label=\small$\square$]
             \item \textbf{Verify authentication tag ($\userComm/\refundComm$).} The circuit recompute $\userComm(\usersk_i)$ or $\refundComm(\usersk_i,\rid)$ from the secret $\usersk_i$ and check if it is consistent with the one stored in the action-dependent map, keyed by $\rid$. 

             For this check we rely on the one-wayness property of $\Hash$, and the hiding properties of the  Midnight layer.
             \item \textbf{Perform action.} For Deposit/Swap/Supply/Reedem the action is to  \textbf{mint new coins}. For Withdraw no action is required.
       \end{itemize}
 \item \textbf{Refund is performed: if transaction failed $\evmOut$= $\failed$.}
      \begin{itemize}[label=\small$\square$]
       \item \textbf{Verify authentication tag ($\userComm/\refundComm$).} The circuit recompute $\userComm$ $(\usersk_i)$ or $\refundComm$( $\usersk_i,\rid)$ from the secret $\usersk_i$ and check if it is consistent with the one stored in the action-dependent map, keyed by $\rid$. 
        For this check we rely on the one-wayness property of $\Hash$, and the hiding properties of the  Midnight layer.
        
        \item \textbf{Refund.} For Withdraw/Swap/Supply/Reedem the action is to  \textbf{mint new coins}. For Deposit no action is required.
        \end{itemize}
    \end{itemize}
\end{enumerate}


Between Step 1 (start request) and Step 2  (complete request)  the MPC system $\SigSys$ attempts to execute the 
request $\rid$ on $\Lpub$. When the request is executed, $\SigSys$ signs the outcome (attestation).
%% Chainging point: here we change and add information about last seen block


\vspace{5pt}
Now, recall that at high-level, and $\Adv$  \emph{violates soundness} if it is able to get out of the system more value than what they contribute in.
In complete Deposit, this means to obtain shielded coins  in $\Lshd$ for an amount $v^*$ higher than what $\Adv$ actually deposited on $\Lpub$; in complete Withdraw this means to cash-out an amount $v^*$ that is higher than what was burned in the request for withdrawing, etc.
This also includes the case where $\Adv$ steals in-flight attestations that are intended for an honest user, and uses them to mint shielded transaction on its own account on $\Lshd$. 

$\Adv$  can successfully execute  \code{completeX}  if and only if it presents  (1) an  attestation $\attSig$, that verifiers under the $\SigSys$'s public key, on a pair $(\rid,\evmOut)$  (2)  the request id $\rid$  is in the map of pending requests,   and (3) a proof of knowledge of the pre-image of the authentication tag ($\userComm$ (for Deposit) or $\refundComm$ for any other call associated to the request id $\rid$).


\subsection{Soundness Arguments Common to all Circuits}
\label{subsec:soudnesscommon}

Let $(\rid^*, \evmOut^*, \attSig^*)$ be what $\Adv$ presents to a \texttt{completeX} circuit, and suppose the circuit accepts. 
(Recall that in every  \code{complete}  action, the circuit will use $\rid^*$ to look up the request $\req^*$
associated to it;   where $\rid=\Kec(\req)$, and  $\attSig$ is a signature over  $(\rid, \evmOut)$: i.e.,  $\attDigest(\rid, \evmOut) = \Kec(\rid \concat \evmOut)$.


% Two events we must rule out.
% First,  that $\SigSys$ never signed $\attDigest(\rid^*, \evmOut^*)$, so the outcome $\Adv$ claims is not the outcome that occurred on
% $\Lpub$ (Case~1). 
% Second, that $\SigSys$ did sign it, but for a request belonging to an honest user, so  $\Adv$ settles value it never surrendered (Case~2).

We now enumerate the ways that adversary $\Adv$ could violate soundness of the authentication patter of  \code{complete}, and we rule them out, using  the cryptographic assumptions that we rely on.


\vspace{8pt}
\noindent
\textbf{Case 1. $\Adv$ claims a false outcome  $\evmOut^*$ for her request $\req^*$.}

In this case, we consider an adversary that performed the \texttt{start} step, and created a request $\req^*$ with id $\rid^*$ to be performed on $\Lpub$, and later it attempts to run a \texttt{completeX} circuit for a false outcome $\evmOut^*$. False means, it does not reflect what truly happened on the chain, e.g., $\Adv$ could claim that the transactions failed when in reality it succeeded, or the opposite.

To pass the verification on a false  $\evmOut^*$ that was  not signed by the MPC (here we are using the \emph{attestation faithfulness assumption} --  $\SigSys$ only signs truth), $\Adv$ might have done one of the following:
\begin{enumerate}

     \item \textbf{Forge $\SigSys$'s  attestation on $\evmOut^*$ $\rightarrow$  forge ECDSA signatures.}  $\Adv$ could  forge an attestation $\sigma^*$  for the pair $\rid^*,\evmOut^*$ of her choice, such that 
    $$\assert \EcdsaVf\bigl(\lnk{fn:attDigest}{\attDigest(\rid^*, \evmOut^*)}\bc \sigma*\bc \mpcRespKey\bigr) \rightarrow 1$$
 regardless of whether the request transaction  $\rid^*$ was actually executed on Ethereum.

  Assuming that ECDSA unforgeability with pre-signatures and key-derivation, this event happens with   negligible probability.


  
    \item \textbf{Re-use a valid signature on a previous input, by breaking the collision-resistance of $\attDigest(\cdot)$.}

    Assume the MPC had signed an attestation for a pair $(\rid,\evmOut)$ at any point in time, and published  $\attSig$.
    If the adversary can craft a pair  $\rid^*$, $\evmOut^*$ such that 
     
      $$\attDigest(\rid, \evmOut) =\attDigest(\rid*, \evmOut*) $$ 

      
      the adversary, could   \emph{re-use}  the  valid signature $\attSig$ for  $(\rid*, \evmOut*)$.

      Due to the collision-resistance property of $\Kec$, this event happens with negligible probability. 

     \item \textbf{Re-use a valid signature on a previous input by break collision-resistance of $\reqIdOf$.}
       Similarly,  an adversary could attempt to  \emph{re-use} $\attSig$ emitted on   $\attDigest$, by crafting a request $\req^*$, such that $\reqIdOf(\req^*) = \reqIdOf(\req)$ for an honest pending request $\req$ published  on the chain.
       The adversary would intercept the signature $\attSig$ and present it to the \code{completeX} function before the honest user.
       
        Due to the collision-resistance property of $\Kec$, this event happens with negligible probability. 


\end{enumerate}


\vspace{8pt}
\noindent
\textbf{Case 3.  $\Adv$ uses a valid attestation for action A, to complete  Action B.}
   $\Adv$ uses an attestation for an action (e.g., Swap) to settle  a different action (e.g., Redeem). 
   Specifically, $\Adv$ use the valid attestation of $(\rid,\evmOut)$  for a request $\rid^*=\reqIdOf(\req^*)$ where $\req^*$ is a swap request, to pass the signature verification of $\cktCompleteSupply$.

   First note that in the implementation each requestID $\rid$ lives in the data structure associated  to the action requested. So, 
   the $\rid$ of swap request lives in the map $\swapView$.
    When $\Adv$ attempts to run the Redeem circuit on
    $(\rid,\evmOut)$, the circuit check if $\rid$ is pending in the map $\redeemView$, which is not existing, and the circuit execution will fail -- unless the adversary managed to create a swap request $\req_{\mathsf{swap}}$ and a redeem request $\req_{\mathsf{redeem}}$ such that the hash $\reqIdOf(\req_{\mathsf{swap}})=\req_{\mathsf{redeem}}$.

    By the \textbf{collision-resistance} property of $\reqIdOf$, the probability of this attack is negligible.
  

 \vspace{8pt}
\noindent
 \textbf{Case 4. $\Adv$ mints more than it paid.}
      $\amt$ and $\tokType$ are not chosen at settlement time: they are read from $\req$, which was
      fixed at request time and whose digest is $\rid$. The same $\amt$ is the amount of the EVM
      transfer the committee signed. Changing it means finding a second record with the same $\rid$, which violated  the  \textbf{collision-resistance}  property of $\reqIdOf()$.

\vspace{8pt}
\noindent
\textbf{Case 5.  $\Adv$ attempt to complete the same request twice (Replay Attack).} 
   $\Adv$ replays a valid $(\rid, \evmOut, \attSig)$. Every complete circuit removes both the record
   entry and the marker entry at $\rid$. A second call reads the
   same entries, and at apply time the node re-executes the transcript against real state, where
   they are gone. This replay attack is addressed at ledger level, and it does not rest on the protocol messages and cryptographic assumptions.
   \vspace{5pt}
  \newversion{The latest release keeps this ledger-level protection (completion removes the
  request's queue entry, its attestation and its event) and adds an explicit mechanism that
  relies on timing information. Time is defined as the relative order of processed
  transactions on the public blockchain, and is tracked via the \emph{block-height} of the
  \emph{last-seen} transaction attested by the MPC: $\cktFlush$ stamps every request with the
  current $\lastseen$ and raises $\lastseen$ to the height of every attestation it folds, and
  an attestation settles a request only if its block height is strictly greater than the
  request's stamp.}

\vspace{8pt}
\noindent
\textbf{Case 6.  $\Adv$ claims an honest user's request.}
In this case, we consider the scenario where $\Adv$, observing the 
requests posted by honest users, and the attestations that 
the MPC posts back, attempts to use a valid attestation on behalf of the honest user.
In other words, $\Adv$ tries  to cash-in, swap, redeem, refund to her account,   money that was deposited by an honest user.

 An adversary could observe when a certain request $\rid$ associated to a certain tag
     is executed on $\Lpub$;  obtain the attestation on it from $\SigSys$ (recall that $\SigSys$ answers to anyone who asks), and use it to  \texttt{completeX}  the request on their \emph{own} account on Midnight.

To successful execute a \texttt{completeX} (or \texttt{refundX}) circuit  for requests posted by an honest user $U_i$, $\Adv$ should do one of the following:

\begin{enumerate}
    \item \textbf{Pass the authentication check for markers  $\userComm(\usersk_i, \cdot)$  (for $\cktCompleteDeposit)$  or  $\refundComm(\usersk_i,\cdot)$ (for all other complete circuits)}

     To pass the authentication check present in any $\mathtt{complete}$ circuit, $\Adv$ must prove  \emph{knowledge of the pre-image } ($\usersk_i)$ of  the authentication tag:
     $\userComm(\usersk_i, \cdot)$ or  $\refundComm(\usersk_i, \cdot)$.
     We note that  secret  $\usersk_i$ is used  in several messages of the protocol: not only to compute the authentication tags  $\userComm(\usersk_i, \cdot)$  and 
     $\refundComm(\usersk_i,\rid)$, but also as witness in the zk proofs that honest user $U_i$ has  published over time. 
   

    Due to the \textbf{zero-knowledge property} (Definition~\ref{def:zk}) of the Midnight shielded-circuit layer,  $\Adv$ does not learn any information about $\usersk_i$ from the proofs generated by the user's honest circuit computation (with all but negligible probability).

      \vspace{3pt}
    Due to the \textbf{pre-image resistance} (Definition \ref{def:crhf}) of $\Hash$ used to computed  $\userComm$ and $\refundComm$,  the adversary cannot prove knowledge of the pre-image of those functions.

    \vspace{3pt}
    Due to the \textbf{knowledge soundness property}  (Definition~\ref{def:zkksound}), the $\Adv$ must know the witness ($\usersk_i$) in order to generate a proof that verifies.    

\end{enumerate}
%Soundness Table -- Omitted
% \begin{table}[t]
% \centering\scriptsize
% \setlength{\tabcolsep}{5pt}
% \renewcommand{\arraystretch}{1.25}
% \begin{tabular}{@{}
%   >{\raggedright\arraybackslash}p{4cm}
%   >{\raggedright\arraybackslash}p{1.9cm}
%   >{\raggedright\arraybackslash}p{4cm}
%   >{\raggedright\arraybackslash}p{2cm}
%   >{\raggedright\arraybackslash}p{4.5cm}@{}}
% \toprule
% & \textbf{gate 1:} $\evmOut$ & \textbf{gate 2:} identity & \textbf{gate 3:} consumed & \textbf{gate 4:} the mint opens to \\
% \midrule
% $\cktCompleteDeposit$ & $\succeeded$ & $\userComm(\usersk) = \req.\keyPath$ & $\reqMap$ & $(\mintNonce, \colr(\tokType), \amt)$, both read from $\req$ \\
% $\cktCompleteWithdraw$ & success bit, either value & \emph{none} if the bit is 1; $\refundComm(\usersk,\rid) = \wdMarkerMap[\rid]$ if 0 & $\reqMap$, $\wdMarkerMap$ & nothing if the bit is 1; else as $\cktRefundValue$ \\
% $\cktCompleteSwap$ & gives $\amtSpent$ & $\refundComm(\usersk,\rid) = \swapMarkerMap[\rid]$ & $\swapMap$, $\swapMarkerMap$ & $(\mintNonce, \colr(\tokTypeB), \amtB)$ and $(\changeNonce, \colr(\tokType), \amtMax - \amtSpent)$ \\
% $\cktCompleteSupply$ & gives $\amtShares$ & $\refundComm(\usersk,\rid) = \supplyMarkerMap[\rid]$ & $\supplyMap$, $\supplyMarkerMap$ & $(\mintNonce, \colrW, \amtShares)$ \\
% $\cktCompleteRedeem$ & gives $\amtAssets$ & $\refundComm(\usersk,\rid) = \redeemMarkerMap[\rid]$ & $\redeemMap$, $\redeemMarkerMap$ & $(\mintNonce, \colrU, \amtAssets)$ \\
% $\cktRefund$ & $= \failed$ & $\refundComm(\usersk,\rid) = $ the marker of the branch taken & the two maps of the branch taken & $(\mintNonce, \colr, v)$, $\colr$ and $v$ read from $\req$ \\
% \bottomrule
% \end{tabular}
% \caption{ The conditions each complete circuit checks. Gate~1 is
% $\EcdsaVf(\attDigest(\rid,\evmOut), \attSig, \mpcRespKey)$ in every row; the column says what else
% is required of $\evmOut$. Gate~3 is enforced by ledger re-execution, the other three by $\zkproof$.
% Here $\colr(a) = \tokenTypeOf(\domSep(a), \vaultAddr)$. Values read from $\req$ were fixed before
% anything was burned; the others come from $\evmOut$.}
% \label{tab:gates}
% \end{table}

% \begin{table}[t]
% \centering\scriptsize
% \setlength{\tabcolsep}{3.5pt}
% \renewcommand{\arraystretch}{1.25}
% \begin{tabular}{@{}
%   >{\raggedright\arraybackslash}p{5.2cm}
%   >{\raggedright\arraybackslash}p{4.3cm}
%   >{\raggedright\arraybackslash}p{2.8cm}
%   >{\raggedright\arraybackslash}p{2.0cm}@{}}
% \toprule
% \textbf{What $\Adv$ tries} & \textbf{What must break} & \textbf{Assumption} & \textbf{Argued in} \\
% \midrule
% Complete with an outcome $\SigSys$ never attested & forge ECDSA under $\mpcRespKey$ & Def.~\ref{def:sig} & Case~1.3 \\
% Re-use one attestation for another $(\rid, \evmOut)$ & collision in $\Kec$ inside $\attDigest$ & Def.~\ref{def:crhf} & Case~1.1 \\
% Point a valid $\rid$ at a different record & collision in $\Kec$ & Def.~\ref{def:crhf} & Case~1.2 \\
% Pass an honest user's identity gate & invert $\Hash$ on a published tag & one-wayness of $\Hash$ & Case~2.1 \\
% Learn $\usersk_i$ from honest proofs & zero-knowledge of the circuit layer & Def.~\ref{def:zk} & Case~2.1 \\
% Skip a gate and still produce a proof & soundness of the argument system & Def.~\ref{def:zksound} & Case~2.2 \\
% Settle one request twice & ledger re-execution & \emph{none (ledger)} & Case~2.3 \\
% Settle a vault-path request as a deposit & invert $\Hash$ at the fixed target $\keyPathLit$ & one-wayness of $\Hash$ & Case~2.4 \\
% Surrender a coin of the wrong color or value & binding of $\coinCommit$, circuit soundness & Def.~\ref{def:binding}, \ref{def:zksound} & \S\ref{susec:soundnessrest} \\
% Make two ERC20s share one color & collision in $\Hash$ inside $\domSep$ & Def.~\ref{def:crhf} & \S\ref{susec:soundnessrest} \\
% Get a settlement \emph{and} a refund & failure faithfulness of $\SigSys$ & system assumption & \S\ref{susec:soundnessrefund} \\
% Inflate the amount minted by $\cktCompleteSupply$/$\cktCompleteRedeem$ & attestation faithfulness & system assumption & \S\ref{susec:soundnessrest} \\
% \bottomrule
% \end{tabular}
% \caption{Every attack considered in this proof and the assumption that excludes it. The last two
% rows are excluded by trust in $\SigSys$, not by cryptography.}
% \label{tab:attackmap}
% \end{table}


Next we analyze the attack vector that is specific to each call.


% \item {\bf Case 4. $\Adv$ settles twice, or settles a non-deposit request.}
%       $\reqMap[\rid]$ is removed at settlement, so a replay fails at apply time. Withdraw, swap,
%       supply, redeem and approve requests also live in $\reqMap$ but carry
%       $\keyPath = \keyPathLit$; settling one needs $\userComm(\usersk^*) = \keyPathLit$. This is the
%       only routing argument that rests on a cryptographic assumption rather than on disjoint maps.
% \end{enumerate}


\subsection{Soundness Argument Specific for Withdraw}
\label{susec:soundnesswithdraw}

In our protocol $\Withdraw$ (Procedure \ref{box:withdraw})  consists of two steps: withdraw's request (Circuit $\cktWithdraw$~\ref{ckt:startwithdraw}) and withdraw's complete (Circuit $\cktCompleteWithdraw$~\ref{ckt:completewithdraw}).
In  withdraw's request, a Midnight user holding account $\shdAcct$ requesting to withdraw  $(\tokType\bc \amt)$ to address  $\destAddr*$  is required to  make a shielded payment to a \emph{burn} address of the value $(\tokType\bc \amt)$.
This is the pre-requisite to claim a withdrawal.
This payment is performed by calling Midnight circuits $\mnhl{\code{receiveShielded}}(\coinIn)$ and $\mnhl{\code{sendImmediateShielded}}(\coinIn, \burnAddr, \amtMax)$, where $(\coinIn)$ contains the secret nounce associated to the coins that the user want to spend, and  $\code{sendImmediateShielded}$ requires  knowledge of the secret  key $\coinsk$  associated to the $\shdAcct$ that is attempting to withdraw.
If both $\mnhl{\code{receiveShielded}}(\coinIn)$  and $\mnhl{\code{sendImmediateShielded}}$  are  successful, then 
an authentication component $\refundComm$ is created for the request ID $\rid$ associated to the withdraw request.
$\refundComm$ will be used only if the withdrawal request is not completed successfully 
and this is indicated by  an attestation $\attSig$ on a $\failed$ string. In such a case, the user executes and use $\attSig$  and the knowledge of the pre-image of $\refundComm$  as a proof 
of entitlement for a coin of value $(\tokType\bc \amt)$.


A  malicious user $\Adv$ violates withdraw soundness if they are able to either withdraw money that they do not  own, or obtain a refund on withdrawal that was completed successfully.
We analyze both cases, and show that they can only happen with negligible probability.
\begin{enumerate}
    \item \textbf{Case 1.  $\Adv$ attempts to withdraw money it does not own.}
     Recall, the withdraw request ($\req \gets (\vaultAddr\bc \sigNonce\bc \keyVer\bc \keyPath\bc \evmTxOf\bc \chainId)$ is considered  by $\SigSys$ only if the user has successfully executed the circuit $\cktWithdraw$.
     This implies that circuits      
     $\mnhl{\code{receiveShielded}}$ $(\coinIn)$ and $\mnhl{\code{sendImmediateShielded}}$ $(\coinIn, \burnAddr, \amtMax)$ have been successfully computed. 
     Consider the following two cases:
            \begin{enumerate}
                \item  \textbf{$\Adv$ owns a coin with amount $\leq\amtMax$ } but still manages to successfully run 
                  $\mnhl{\code{receiveShielded}}$ $(\coinIn)$ for the amount $\amt$ asserted in the circuit $\cktWithdraw$.

                  
                 Due to the \textbf{soundness} of the execution of the  Midnight circuit $\mnhl{\code{receiveShielded}}$, this case happens with  negligible probability.
                 \item  \textbf{$\Adv$ burns  coins owned by an honest user.}
                   An adversary $\Adv$ could observe  honest user's executions of  $\mnhl{\code{receiveShielded}}$  in multiple deposits/swaps/supply/redeem executions and attempt to use/manipulate the collected information into a successful execution of  $\mnhl{\code{receiveShielded}}$ $(\coinIn^H)$ with an honest coin $\coinIn^H$.  
                   
                   Due to \textbf{the hiding, and zero-knowledge properties} of the primitives used   by Midnight's library, this event could only happen with negligible probability.  

           \end{enumerate}
    \item \textbf{Case 2. $\Adv$'s withdraw request $\rid$ is successfully executed on Ethereum \emph{and} a refund on Midnight for $\rid$.}  This case corresponds to \textbf{Case 1} analyzed in \S~\ref{subsec:soudnesscommon}.
\end{enumerate}


\subsection{Soundness Arguments Specific to Swap/Supply/Redeem Circuits}
\label{susec:soundnessrest}

Swap/Supply/Redeem have the same logic.
A user who owns $\amt$ coins of \emph{type A} requests to swap them 
for $\amtB$ coins of \emph{type B}.
They differ in three respects: the \emph{admissible types}, how the returned amount is calculated,
and how many coins the settlement mints; $\Swap$ mints two (the bought coin and the change),
$\Supply$ and $\Redeem$ mint one. 
A fourth, smaller difference: $\cktCompleteSwap$ reads the stored
request, while $\cktCompleteSupply$ and $\cktCompleteRedeem$ obtain the information from 
the attested outcome.

In Swap,  user chooses the type \emph{type B}, that  she wants to swap to,  and there is additional input $\amtMax$,  chosen by the user, which  defines  the most she is willing to spend
She surrenders $\amtMax$ up front, the router spends some
$\amtSpent \leq \amtMax$, and $\cktCompleteSwap$ returns $\amtMax - \amtSpent$ as change. If the
router would need more than $\amtMax$, the EVM call reverts and the flow goes to $\cktRefundSwap$.
In Supply, and Redeem, the types are fixed and hardwired in the contract definitions, and the amounts  are chosen by the market. 
The types are defined in two addresses $\underAddr$ and $\wrapAddr$ written in the public state of the vault by $\cktInit$ and are  never chosen by the caller.
In Supply,  \emph{Type A} is a fixed $\colrU$  and type B is the fixed $\colrW$, and result of the call, the user gets  $\amtShares$ of $\colrW$.
In Redeem,  \emph{Type A} is $\colrW$  and the amount redeemed  is decided by the market and communicated to the circuit directly by the blockchain record.  
The Swap/Supply/Redeem   follow  the same authentication pattern of Widthdraw flow, for the refund path. As such, the  arguments of soundness for  Swap  follow very closely the arguments on Withdrawal. The additional attack vector a  malicious user can leverage in Swap/Supply/Redeem, is to initiate a request $\req$ for token types and amounts that do not correctly match the types and amounts of the coins it possesses. 

We identify the following cases:
%

\begin{enumerate}[leftmargin=1.8em,itemsep=5pt]
    \item \textbf{$\Adv$ surrenders a coin that does not match the request it posted.} $\Adv$
could attempt to burn a coin of a cheap color, or of a value below $\amt$, while filing a request
for $\amt$ of an expensive $\tokType$, and so withdraw on $\Lpub$ more than it gave up on
$\Lshd$.

This attack is prevented by the checks performed in the circuit.
$\assert \coinIn.\mathsf{color}$ = ${\tokenTypeOf}(\domSep(\tokType), \vaultAddr)$,   and $\coinIn.\mathsf{value} = \amtMax$. 
% Due to the soundness of the zk proofs generated by the execution of Midnight circuit layer,  an adversary cannot attempt to swap/supply/redeem  coins that  are not consistent with the request/ or with the colors allowed by the system.

The surrendered coin's color and value are asserted \emph{equal}
to the values that go into $\req$, and by binding of the coin commitment
(Definition~\ref{def:binding}) the descriptor $\coinIn$ the circuit checks is the same coin whose
commitment and nullifier are published. By the \textbf{soundness of the  zk argument system}
(Definition~\ref{def:zksound}) a proof that passes verification cannot correspond to an execution in
which those assertions failed. Hence the event happens with probability negligible in $\secpar$.

\noindent The same  case applies to $\Supply$, with 
with $(\colrU, \amt)$, and to $\Redeem$ with $(\colrW, \amtShares)$. In all four the assertion is Step~2 of the request circuit and the argument
is the same two-line appeal to binding and to circuit soundness.


\item \textbf{$\Adv$ surrenders a coin of the wrong color or value.} Prevented by  Step~2 of
      \cktref{ckt:startsupply} and \cktref{ckt:startredeem}. Note that the
      two colors here are \emph{pinned}, not caller-supplied: $\colrU$ and $\colrW$ are functions
      of $\underAddr$ and $\wrapAddr$ read from state, so unlike $\cktWithdraw$, where the caller
      names $\tokType$, there is no color argument to get wrong.


\item \textbf{$\Adv$ inflates the minted amount.} This is specific to Supply and Redeem only.
      $\cktCompleteSupply$ and $\cktCompleteRedeem$ take the minted value from $\evmOut$ and read no
      field of $\req$. The only thing binding that value is that
      $\attSig$ verifies over $\attDigest(\rid, \evmOut)$, so an inflated mint requires a forged or
      re-used attestation, excluded as above, \emph{or} a committee that attests a share count
      the market did not return, which is excluded by assumption on the trustworthiness of the routing system and the chain,  and not by the protocol. 
\end{enumerate}



\section{Take Away from this document}
\label{sec:conclusions}



\subsection{Take away for a Midnight User}

\begin{itemize}
    \item \textbf{This  cross-chain system  involves a  public chain so all amounts are public!}.
Even if you can call actions  such as Withdraw/Swap/Supply/Redeem anonymously on Midnight (i.e., without revealing your identity),  the amounts of those actions   must be executed on the public chain, so those amounts  are still public, even if your identity is not.
\item Because amounts are public, \textbf{you have to be careful when you choose the amounts and timing when performing  Withdraw/Swap/Supply.}
For instance,  if you Deposit  5 USD, and then immediately Withdraw 5 USD on a fresh, and therefore unlinkable to the address you used to deposit,  ethereum wallet, the amounts and timings still gives a way a link between the deposit and the withdraw.
Similar arguments apply for Swap/Supply/Redeem.



\item \textbf{Withdraw is unlikable to you,   only if you provide a fresh Ethereum account.}
If you withdraw funds into an account that you previously used, then the withdraw is connected to your previous activity.
In the worst case, if you withdraw to an Ethereum account that you used to fund your  derived address used in Deposit, then you have linked deposit and withdraw.

\item \textbf{Keep your private coins consolidate in one coin. }
The number of coins you need to collect in order to pay for a Withdraw/Swap/Supply
leaks information about who you are.
This is true regardless of whether you use our tool, or you simply operate on Midnight.
You should always have your coins consolidated in one coin, and when you pay you should always have a change coin.



\item  \textbf{Deposit requests are all linkable to you. }
Our system assigns a fixed identity $\userComm(\usersk_i)$ to each user that is used for each deposit.
This makes all deposits requested by the owner of  $\usersk_i$ linkable.
(all other actions are not linkable since the authentication tag that we use, takes a one-time string $\rid$ that makes each tag unique.)


\item For privacy,  \textbf{secret $\usersk_i$  must be kept secret at all times.}
If someone learns $\usersk_i$, they can then link all your previous activities.

\item On the bright side, if you leak $\usersk_i$,
the adversary can at most steal one pending request.



\item Your anonymity set  in Swap/Supply/Redeem is restricted to  the number of people of are transacting  on those particular coins.

\end{itemize}





\subsection{For the Developers}

\begin{itemize}
    \item \textbf{A fixed $\keyPath_i$ per user is only a design choice.} In our implementation,  the user's software picks a random application secret $\usersk_i$ as a permanent secret used across all deposits (and all calls). This design choice prioritizes user's convenience over stronger privacy. Convenience come from the fact that the user will only need to fund a single derived account (dictated by  $\userComm_i$).
    This gives weaker privacy  because (1) all deposits to the user's account are strictly linkable  -- even if they came from different ethereum accounts.  (2) leakage of the application secret $\usersk_i$ deanonymize all user's past activities.
\end{itemize}

\subsection{For the Cryptographers}

This treatment focuses on the erc20 vault contract,  it is not an attempt to analyze the entire system rigorously.
As such, we made a few simplifications, and some arguments are  informal.


\begin{itemize}
    \item Unlinkability: in our systems the identities  of the caller are used in two layers: in our protocol messages and in Midnight's protocol messages.
    In this proof, we identify which Midnight's items contributed to our security and we rely on them, but we do not prove formally that those objects achieve the property that are claimed.


    \item Soundness: our definition of soundness could be  made more comprehensive using the UC framework. We defer it to the next iteration of the document.


\end{itemize}


% Every entry is listed, cited or not, as further reading.
\nocite{*}
\bibliographystyle{plain}
\bibliography{references}

\end{document}
